% Lutte contre les comportements préjudiciables à la santé reproductive — SVT 1ère TI — Cours signé OnBuch+
% Programme officiel MINESEC. Compiler avec : tectonic NCT05-lutte-comportements-sante-reproductive.tex
\newcommand{\DOCMATIERE}{SVT}
\newcommand{\DOCNIVEAU}{1ère TI}
\newcommand{\DOCMODULE}{Module 2 : L'Éducation à la Santé}
\newcommand{\DOCLECON}{NCT05}
\newcommand{\DOCTITRE}{Lutte contre les comportements préjudiciables à la santé reproductive}
% OnBuch+ — préambule « cours signés » ENRICHI (compilé avec tectonic / XeLaTeX)
% Système visuel « Pop » d'OnBuch+ : fond crème (jamais blanc), bordures encre
% épaisses, ombres « dures » décalées, titres Archivo Black en capitales.
% Version MATHÉMATIQUES : graphes (pgfplots/tikz), boîtes pédagogiques
% (définition, propriété, méthode, attention, le savais-tu, exercice résolu,
% exercice, TP, bilan), page de garde et signature OnBuch+ ancrée en bas.
%
% Macros attendues dans main.tex AVANT \input{preamble} :
%   \DOCMATIERE  \DOCNIVEAU  \DOCMODULE  \DOCLECON (numéro)  \DOCTITRE
\documentclass[11pt]{article}

\usepackage{fontspec}
\usepackage[french]{babel}
\usepackage[a4paper,margin=2cm,top=3.4cm,bottom=2.8cm,headheight=1.4cm,footskip=1.3cm]{geometry}
\usepackage{amsmath,amssymb,amsfonts}
\usepackage{mathtools} % \xmapsto, \coloneqq... souvent utilisés par les modèles
\usepackage{cancel} % simplifications barrées (bilans de forces)
% \abs{z} (module, valeur absolue) et \norm{u} (norme), utilisés par les modèles
\providecommand{\abs}{}\let\abs\relax\DeclarePairedDelimiter{\abs}{\lvert}{\rvert}
\providecommand{\norm}{}\let\norm\relax\DeclarePairedDelimiter{\norm}{\lVert}{\rVert}
\usepackage[table]{xcolor} % [table] : fournit aussi \rowcolors (alternance de lignes)
\usepackage{pagecolor}
\usepackage{tikz}
\usetikzlibrary{arrows.meta,positioning,calc,shapes.geometric,shapes.misc,shapes.arrows,decorations.pathmorphing,decorations.pathreplacing,decorations.markings,patterns,shadows,mindmap,trees,fit,backgrounds,matrix,intersections,angles,quotes}
\usepackage{pgfplots}
\usepackage[version=4]{mhchem} % équations nucléaires \ce{^{235}_{92}U}
\usepackage[european,siunitx]{circuitikz} % schémas électriques
\pgfplotsset{compat=1.18}
\usepgfplotslibrary{fillbetween,groupplots} % aires entre courbes (calcul intégral)
\usepackage{enumitem}
\usepackage{fancyhdr}
% [most] charge toutes les bibliothèques tcolorbox (listings, minted, raster,
% documentation...) et gonfle inutilement la mémoire TeX sur un document long
% (~300 boîtes) : on ne charge que ce qui est réellement utilisé.
\usepackage[skins,breakable]{tcolorbox}
\usepackage{graphicx}
\usepackage{colortbl}
\usepackage{booktabs}
\usepackage{tabularx}
\usepackage{diagbox} % case d'en-tête diagonale des tableaux à double entrée
% Colonnes extensibles centrée / gauche / droite (C, L, R), souvent utilisées par les modèles
\newcolumntype{C}{>{\centering\arraybackslash}X}
\newcolumntype{L}{>{\raggedright\arraybackslash}X}
\newcolumntype{R}{>{\raggedleft\arraybackslash}X}
\usepackage{array}
\usepackage{multicol}
\usepackage{siunitx}
% parse-numbers=false : accepte aussi \SI{2,5 \times 10^{-3}}{...}
\sisetup{locale=FR,output-decimal-marker={,},group-minimum-digits=4,parse-numbers=false}
\DeclareSIUnit\ohm{\ensuremath{\symup{\Omega}}} % U+2126 absent de la police
\DeclareSIUnit\division{div} % réglages d'oscilloscope (ms/div, V/div)
\DeclareSIUnit\tr{tr} % tours (vitesse de rotation, ex. \SI{1500}{\tr\per\minute})
\DeclareSIUnit\ch{ch} % cheval-vapeur (puissance moteur, unité usuelle hors SI)
\DeclareSIUnit\franccfa{FCFA} % franc CFA (coûts, contexte camerounais)
\DeclareSIUnit\dioptre{\ensuremath{\delta}} % vergence d'une lentille (dioptrie)
\usepackage{qrcode}
% \degree : commande fréquemment utilisée par les modèles pour un angle ou une
% température en dehors de \SI{}{} (non fournie par les paquets chargés).
\providecommand{\degree}{^{\circ}}
% \qed (fin de démonstration), utilisé par les modèles sans amsthm
\providecommand{\qed}{\hfill$\square$}
% \barre : double barre des valeurs interdites dans un tableau de variations (tkz-tab)
\providecommand{\barre}{\ensuremath{\|}}
% environnement proof (amsthm non chargé) utilisé par les modèles
\ifdefined\proof\else\newenvironment{proof}[1][Démonstration]{\par\noindent\textit{#1.}\ }{\qed\par}\fi
\usepackage{lastpage}
\usepackage{hyperref}
\hypersetup{hidelinks}

% ============================================================
% Palette « Pop » OnBuch+ — mêmes hex que la classe P (plus_ui.dart)
% ============================================================
\definecolor{poporange}{HTML}{F59321}
\definecolor{popdark}{HTML}{E07A0C}
\definecolor{popcream}{HTML}{FFF6E8}
\definecolor{popink}{HTML}{1C1714}
\definecolor{poppurple}{HTML}{7A5AE0}
\definecolor{popgreen}{HTML}{2E9E6A}
\definecolor{poppink}{HTML}{E0457B}
\definecolor{popblue}{HTML}{2F7DE1}
\definecolor{popgold}{HTML}{C98A12}
\definecolor{popmuted}{HTML}{8A7F76}
\definecolor{popline}{HTML}{EADFD2}
\definecolor{popgray}{HTML}{8A7F76}
\colorlet{popgrey}{popgray}
% Alias tolérants (le modèle nomme parfois les couleurs différemment)
\colorlet{popwhite}{white}
\colorlet{popblack}{popink}
\colorlet{popred}{poppink}
\colorlet{popyellow}{popgold}
% Teintes claires pour fonds de tableaux / zones
\colorlet{popblueL}{popblue!10!white}
\colorlet{poporangeL}{poporange!14!white}
\colorlet{popgreenL}{popgreen!12!white}
\colorlet{poppurpleL}{poppurple!12!white}
\colorlet{poppinkL}{poppink!12!white}
\colorlet{popgoldL}{popgold!14!white}

\pagecolor{popcream}

% ============================================================
% Typographie
% ============================================================
\defaultfontfeatures{Ligatures=TeX}
\setmainfont{PlusJakartaSans-Medium.ttf}[Path=../../../fonts/,BoldFont=PlusJakartaSans-Bold.ttf]
% Maths en sans-serif (Fira Math) avec chiffres et lettres droites de Plus
% Jakarta, pour que les formules se fondent dans le texte.
\usepackage[math-style=ISO,bold-style=ISO]{unicode-math}
\setmathfont{FiraMath-Regular.otf}[Path=../../../fonts/]
\setmathfont{PlusJakartaSans-Medium.ttf}[Path=../../../fonts/,range={up/num,up/{Latin,latin},\mathup}]
\usepackage{icomma} % virgule décimale sans espace en mode math
% Caractères mathématiques Unicode absents des polices de texte (Plus Jakarta,
% Archivo Black) : rendus par leur équivalent mathématique, en texte comme en
% formule — sinon ils s'affichent en carré vide (ex. « Arithmétique dans ℤ »).
\usepackage{newunicodechar}
\newunicodechar{ℝ}{\ensuremath{\mathbb{R}}}
\newunicodechar{ℤ}{\ensuremath{\mathbb{Z}}}
\newunicodechar{ℂ}{\ensuremath{\mathbb{C}}}
\newunicodechar{ℕ}{\ensuremath{\mathbb{N}}}
\newunicodechar{ℚ}{\ensuremath{\mathbb{Q}}}
\newunicodechar{𝔻}{\ensuremath{\mathbb{D}}}
\newunicodechar{α}{\ensuremath{\alpha}}
\newunicodechar{β}{\ensuremath{\beta}}
\newunicodechar{γ}{\ensuremath{\gamma}}
\newunicodechar{δ}{\ensuremath{\delta}}
\newunicodechar{ε}{\ensuremath{\varepsilon}}
\newunicodechar{θ}{\ensuremath{\theta}}
\newunicodechar{λ}{\ensuremath{\lambda}}
\newunicodechar{μ}{\ensuremath{\mu}}
\newunicodechar{σ}{\ensuremath{\sigma}}
\newunicodechar{τ}{\ensuremath{\tau}}
\newunicodechar{φ}{\ensuremath{\varphi}}
\newunicodechar{ψ}{\ensuremath{\psi}}
\newunicodechar{ω}{\ensuremath{\omega}}
\newunicodechar{Σ}{\ensuremath{\Sigma}}
% majuscules produites par \MakeUppercase dans les titres (α-aminés -> Α)
\newunicodechar{Α}{\ensuremath{\alpha}}
\newunicodechar{Β}{\ensuremath{\beta}}
\newunicodechar{Γ}{\ensuremath{\Gamma}}
\newunicodechar{Δ}{\ensuremath{\Delta}}
\newunicodechar{Ε}{\ensuremath{\varepsilon}}
\newunicodechar{Θ}{\ensuremath{\Theta}}
\newunicodechar{Λ}{\ensuremath{\Lambda}}
\newunicodechar{Μ}{\ensuremath{\mu}}
\newunicodechar{Τ}{\ensuremath{\tau}}
\newunicodechar{Φ}{\ensuremath{\Phi}}
\newunicodechar{Ψ}{\ensuremath{\Psi}}
\newunicodechar{Ω}{\ensuremath{\Omega}}
\newunicodechar{↦}{\ensuremath{\mapsto}}
\newunicodechar{∈}{\ensuremath{\in}}
\newunicodechar{∉}{\ensuremath{\notin}}
\newunicodechar{∧}{\ensuremath{\wedge}}
\newunicodechar{⇔}{\ensuremath{\Leftrightarrow}}
\newunicodechar{⟺}{\ensuremath{\Longleftrightarrow}}
\newunicodechar{⇒}{\ensuremath{\Rightarrow}}
\newunicodechar{⟹}{\ensuremath{\Longrightarrow}}
\newunicodechar{∘}{\ensuremath{\circ}}
\newunicodechar{∪}{\ensuremath{\cup}}
\newunicodechar{∩}{\ensuremath{\cap}}
\newunicodechar{≡}{\ensuremath{\equiv}}
\newunicodechar{⊂}{\ensuremath{\subset}}
\newunicodechar{⊥}{\ensuremath{\perp}}
\newunicodechar{∃}{\ensuremath{\exists}}
\newunicodechar{∀}{\ensuremath{\forall}}
\newunicodechar{∛}{\ensuremath{\sqrt[3]{\,}}}
\newunicodechar{∜}{\ensuremath{\sqrt[4]{\,}}}
\newunicodechar{⃗}{\ensuremath{{}^{\to}}}
\newunicodechar{ⁿ}{\ifmmode^{n}\else\textsuperscript{\ensuremath{n}}\fi}
\newunicodechar{ˣ}{\ifmmode^{x}\else\textsuperscript{\ensuremath{x}}\fi}
\newunicodechar{ᵇ}{\ifmmode^{b}\else\textsuperscript{\ensuremath{b}}\fi}
\newunicodechar{ʳ}{\ifmmode^{r}\else\textsuperscript{\ensuremath{r}}\fi}
\newunicodechar{ᵘ}{\ifmmode^{u}\else\textsuperscript{\ensuremath{u}}\fi}
\newunicodechar{ʸ}{\ifmmode^{y}\else\textsuperscript{\ensuremath{y}}\fi}
\newunicodechar{ᵃ}{\ifmmode^{a}\else\textsuperscript{\ensuremath{a}}\fi}
\newunicodechar{ᵉ}{\ifmmode^{e}\else\textsuperscript{\ensuremath{e}}\fi}
\newunicodechar{ᵏ}{\ifmmode^{k}\else\textsuperscript{\ensuremath{k}}\fi}
\newunicodechar{ᵖ}{\ifmmode^{p}\else\textsuperscript{\ensuremath{p}}\fi}
\newunicodechar{ᴺ}{\ifmmode^{N}\else\textsuperscript{\ensuremath{N}}\fi}
\newunicodechar{ᶜ}{\ifmmode^{c}\else\textsuperscript{\ensuremath{c}}\fi}
\newunicodechar{ʰ}{\ifmmode^{h}\else\textsuperscript{\ensuremath{h}}\fi}
\newunicodechar{ᵠ}{\ifmmode^{q}\else\textsuperscript{\ensuremath{q}}\fi}
\newunicodechar{ₙ}{\ifmmode_{n}\else\textsubscript{\ensuremath{n}}\fi}
\newunicodechar{ₐ}{\ifmmode_{a}\else\textsubscript{\ensuremath{a}}\fi}
\newunicodechar{ᵢ}{\ifmmode_{i}\else\textsubscript{\ensuremath{i}}\fi}
\newunicodechar{ᵣ}{\ifmmode_{r}\else\textsubscript{\ensuremath{r}}\fi}
\newunicodechar{ₖ}{\ifmmode_{k}\else\textsubscript{\ensuremath{k}}\fi}
\newunicodechar{ᵦ}{\ifmmode_{\beta}\else\textsubscript{\ensuremath{\beta}}\fi}
\newunicodechar{ₓ}{\ifmmode_{x}\else\textsubscript{\ensuremath{x}}\fi}
\newfontfamily\popdisplay{ArchivoBlack-Regular.ttf}[Path=../../../fonts/]
\newfontfamily\poplogo{SpaceGrotesk-Bold.ttf}[Path=../../../fonts/]
\newfontfamily\popsemi{PlusJakartaSans-SemiBold.ttf}[Path=../../../fonts/]
\newfontfamily\popxbold{PlusJakartaSans-ExtraBold.ttf}[Path=../../../fonts/]
\newfontfamily\popxboldls{PlusJakartaSans-ExtraBold.ttf}[Path=../../../fonts/,LetterSpace=3.0]

\setlist[enumerate]{leftmargin=1.7em,itemsep=5pt,topsep=4pt}
\setlist[itemize]{leftmargin=1.7em,itemsep=4pt,topsep=4pt,label={\color{poporange}\textbullet}}
\setlength{\parindent}{0pt}
\setlength{\parskip}{5pt}
\renewcommand{\arraystretch}{1.25}

% pgfplots : style Pop par défaut
\pgfplotsset{
  every axis/.append style={axis line style={popink,line width=1pt},tick style={popink},
    grid style={popline,line width=0.5pt},label style={font=\small},tick label style={font=\footnotesize}},
  popcurve/.style={poporange,line width=1.8pt,no markers,smooth},
}
% Même style disponible dans un \draw TikZ simple (hors pgfplots)
\tikzset{popcurve/.style={poporange,line width=1.8pt}}
\tikzset{popgrid/.style={popline,very thin}}
\colorlet{popcurve}{poporange} % le nom du style est parfois utilisé comme couleur

% ============================================================
% Pill / squiggle
% ============================================================
\newcommand{\poppill}[3]{%
  \tikz[baseline=-0.55ex]\node[draw=popink,line width=1.2pt,rounded corners=2.6pt,%
    fill=#2,inner xsep=6.5pt,inner ysep=3pt]{%
    \popxboldls\fontsize{7}{7}\selectfont\color{#3}\MakeUppercase{#1}};%
}
\newcommand{\popsquiggle}[1]{%
  \tikz[baseline=-0.4ex]{
    \draw[#1,line width=2.6pt,line cap=round]
      plot [smooth] coordinates {(0,0.09) (0.34,0.01) (0.68,0.21) (1.02,0.01) (1.36,0.10)};
  }%
}
% Mot-clé surligné (marqueur orange sous le texte)
\newcommand{\cle}[1]{{\popxbold\color{popdark}#1}}

% ============================================================
% En-tête / pied
% ============================================================
\pagestyle{fancy}
\fancyhf{}
\renewcommand{\headrulewidth}{0pt}
\renewcommand{\footrulewidth}{0pt}
\lhead{\raisebox{-0.15ex}{\poplogo\fontsize{13}{13}\selectfont\color{popink}OnBuch\color{poporange}+}}
\rhead{\poppill{\DOCMATIERE\ \textbullet\ \DOCNIVEAU\ \textbullet\ Leçon \DOCLECON}{white}{popink}}
\lfoot{\popxbold\fontsize{7.6}{7.6}\selectfont\color{popmuted}\MakeUppercase{Cours signé OnBuch+}\ \textbullet\ {\popsemi\DOCTITRE}}
\rfoot{\tikz[baseline=-0.6ex]\node[rounded corners=8.5pt,fill=popink,minimum height=17pt,minimum width=17pt,inner xsep=5pt,inner sep=0]{\popxbold\fontsize{8}{8}\selectfont\color{popcream}\thepage\,/\,\pageref*{LastPage}};}

% ============================================================
% Titres
% ============================================================
\newcommand{\chaptitre}[2]{%
  \begin{tcolorbox}[enhanced,colback=poporange,colframe=popink,boxrule=2.6pt,arc=11pt,
      drop shadow=popink,left=15pt,right=15pt,top=12pt,bottom=11pt,before skip=3pt,after skip=18pt]
    \poppill{#2}{popink}{popcream}\par\vspace{9pt}
    {\raggedright\hyphenpenalty=10000\exhyphenpenalty=10000\popdisplay\fontsize{22}{24}\selectfont\color{popcream}\MakeUppercase{#1}\par}
  \end{tcolorbox}
}
% Section numérotée : pastille orange avec le numéro + titre Archivo
\newcounter{popsec}
\newcommand{\coursec}[1]{%
  \stepcounter{popsec}\setcounter{popsub}{0}%
  \par\vspace{16pt}\noindent
  \begin{minipage}{\linewidth}
  \tikz[baseline=-0.7ex]\node[draw=popink,line width=1.6pt,fill=poporange,rounded corners=4pt,minimum size=22pt,inner sep=0]{\popdisplay\fontsize{12}{12}\selectfont\color{popcream}\thepopsec};%
  \hspace{8pt}{\popdisplay\fontsize{15}{17}\selectfont\color{popink}\MakeUppercase{#1}}\par
  \vspace{4pt}\noindent{\color{poporange}\rule{36pt}{3.6pt}}
  \end{minipage}\par\nopagebreak\vspace{8pt}
}
\newcounter{popsub}
\newcommand{\courssub}[1]{%
  \stepcounter{popsub}%
  \par\vspace{9pt}\noindent{\popxbold\fontsize{11.5}{13}\selectfont\color{popink}{\color{poporange}\thepopsec.\thepopsub}\hspace{6pt}#1}\par\nopagebreak\vspace{4pt}
}

% ============================================================
% Boîtes pédagogiques — même recette : fond blanc, bordure épaisse colorée,
% ombre dure encre, badge plein. Titre optionnel [#1].
% ============================================================
\tcbset{popbase/.style={breakable,enhanced,colback=white,boxrule=2.3pt,arc=9pt,
  shadow={3pt}{-3pt}{0pt}{popink},left=14pt,right=14pt,top=11pt,bottom=11pt,before skip=11pt,after skip=15pt,
  fontupper=\popsemi\fontsize{10.6}{14.5}\selectfont\color{popink},
  fontlower=\popsemi\fontsize{10.6}{14.5}\selectfont\color{popink}}}
% Coche dessinée (le glyphe ✓ n'existe pas dans les polices du document)
\newcommand{\popcheck}[1]{\tikz[baseline=-0.5ex]\draw[#1,line width=1.6pt,line cap=round,line join=round](-0.13,0.0)--(-0.04,-0.09)--(0.13,0.1);}
\providecommand{\coche}{\popcheck{popgreen}}
\providecommand{\resultat}[1]{\par\smallskip\noindent\fcolorbox{popgreen}{popgreenL}{\parbox{\dimexpr\linewidth-2\fboxsep-2\fboxrule\relax}{\textbf{Résultat.} #1}}\par\smallskip}
\newcommand{\popboxhead}[3]{\poppill{#1}{#2}{white}\ifx\relax#3\relax\else\hspace{6pt}{\popxbold\fontsize{10.6}{12}\selectfont\color{popink}#3}\fi\par\vspace{7pt}}

\newtcolorbox{aretenir}[1][]{popbase,colframe=popgreen,before upper={\popboxhead{À retenir}{popgreen}{#1}}}
\newtcolorbox{exemplebox}[1][]{popbase,colframe=poporange,before upper={\popboxhead{Exemple}{poporange}{#1}}}
\newtcolorbox{definition}[1][]{popbase,colframe=popblue,before upper={\popboxhead{Définition}{popblue}{#1}}}
\newtcolorbox{propriete}[1][]{popbase,colframe=poppurple,before upper={\popboxhead{Propriété}{poppurple}{#1}}}
\newtcolorbox{methode}[1][]{popbase,colframe=popgold,before upper={\popboxhead{Méthode}{popgold}{#1}}}
\newtcolorbox{attention}[1][]{popbase,colframe=poppink,before upper={\popboxhead{Attention}{poppink}{#1}}}
\newtcolorbox{experience}[1][]{popbase,colframe=popblue,colback=popblueL,before upper={\popboxhead{Expérience}{popblue}{#1}}}
% « Le savais-tu ? » : carte violette pleine (contexte, culture, Cameroun)
\newtcolorbox{savaistu}[1][]{popbase,colback=poppurple,colframe=popink,
  fontupper=\popsemi\fontsize{10.4}{14.2}\selectfont\color{white},
  before upper={\poppill{Le savais-tu\,?}{white}{poppurple}\ifx\relax#1\relax\else\hspace{6pt}{\popxbold\color{white}#1}\fi\par\vspace{7pt}}}
% Exercice résolu : énoncé en haut, \tcblower puis la solution sur fond crème
\newcounter{popexo}\newcounter{popexr}
\newtcolorbox{exoresolu}[1][]{popbase,colframe=poporange,colbacklower=popcream,
  segmentation style={popink,line width=1.2pt,dashed},
  before upper={\stepcounter{popexr}\popboxhead{Exercice résolu \thepopexr}{poporange}{#1}},
  before lower={\poppill{Solution}{popink}{popcream}\par\vspace{6pt}}}
% Exercice (non résolu en place) — la correction est dans la partie Corrigés
\newtcolorbox{exercice}[1][]{popbase,colframe=popink,
  before upper={\stepcounter{popexo}\popboxhead{Exercice \thepopexo}{popink}{#1}}}
% Corrigé
\newtcolorbox{corrige}[1][]{popbase,colframe=popgreen,colback=popgreenL,
  before upper={\popboxhead{Corrigé}{popgreen}{#1}}}
% Objectifs / prérequis / bilan
\newtcolorbox{objectifs}{popbase,colframe=popink,colback=poporangeL,
  before upper={\popboxhead{Objectifs de la leçon}{popink}{}}}
\newtcolorbox{prerequis}{popbase,colframe=popink,colback=popblueL,
  before upper={\popboxhead{Prérequis}{popblue}{}}}
\newtcolorbox{bilan}{popbase,colframe=popink,colback=white,boxrule=2.8pt,
  before upper={\popboxhead{Fiche bilan}{poporange}{L'essentiel à connaître pour l'examen}}}
% Figure encadrée avec légende
\newtcolorbox{popfigure}[1][]{enhanced,breakable=false,colback=white,colframe=popline,boxrule=1.4pt,arc=8pt,
  left=10pt,right=10pt,top=10pt,bottom=8pt,before skip=10pt,after skip=12pt,halign=center,
  lower separated=false,halign lower=center,
  fontlower=\popsemi\fontsize{9}{11.5}\selectfont\color{popmuted},
  #1}
\newcommand{\legende}[1]{\par\vspace{6pt}{\popsemi\fontsize{9}{11.5}\selectfont\color{popmuted}#1}}

% ============================================================
% Signature OnBuch+
% ============================================================
\newcommand{\poplogoinline}[1]{{\poplogo\fontsize{#1}{#1}\selectfont\color{popink}OnBuch\color{poporange}+}}
\newcommand{\signatureonbuch}{%
  \begin{tcolorbox}[enhanced,colback=popink,colframe=popink,boxrule=2.2pt,arc=10pt,
      drop shadow=poporange,left=14pt,right=14pt,top=10pt,bottom=10pt,before skip=0pt,after skip=0pt]
    \tikz[baseline=-0.7ex]\node[circle,fill=popgreen,draw=popcream,line width=1.3pt,minimum size=22pt,inner sep=0]{\popcheck{white}};%
    \hspace{9pt}\begin{minipage}[c]{0.62\linewidth}
      {\popxbold\fontsize{12}{13}\selectfont\color{popcream}Signé\ \textbullet\ {\poplogo OnBuch\color{poporange}+}}\par\vspace{2pt}
      {\raggedright\popsemi\fontsize{8}{10.5}\selectfont\color{popline}\MakeUppercase{Équipe pédagogique OnBuch+}\\\MakeUppercase{Conforme au programme officiel MINESEC}\par}
    \end{minipage}\hfill
    {\poplogo\fontsize{22}{22}\selectfont\color{popcream}OnBuch\color{poporange}+}
  \end{tcolorbox}%
}

% ============================================================
% Page de garde
% #1 titre · #2 matière · #3 niveau · #4 module · #5 n° leçon · #6 durée ·
% #7 description / objectifs (texte ou itemize)
% ============================================================
\newcommand{\pagegarde}[7]{%
  \thispagestyle{empty}
  \newgeometry{a4paper,margin=1.7cm,top=1.6cm,bottom=1.4cm}%
  \begin{tikzpicture}[remember picture,overlay]
    \fill[poporange!18] ([xshift=-3.2cm,yshift=-3.6cm]current page.north east) circle (4.6cm);
    \fill[poppurple!14] ([xshift=2.4cm,yshift=7.6cm]current page.south west) circle (3.8cm);
  \end{tikzpicture}%
  {\poplogo\fontsize{30}{30}\selectfont\color{popink}OnBuch\color{poporange}+}\hfill
  \raisebox{0.6ex}{\poppill{Cours signé \textbullet\ Premium}{popink}{popcream}}\par
  \vspace{1pt}\popsquiggle{poppurple}\par
  \vspace{8pt}
  \begin{tcolorbox}[enhanced,colback=poporange,colframe=popink,boxrule=2.8pt,arc=13pt,
      drop shadow=popink,left=18pt,right=18pt,top=12pt,bottom=13pt,after skip=10pt]
    \poppill{#2 \textbullet\ #3}{popink}{popcream}\hspace{5pt}\poppill{#4}{white}{popink}\par\vspace{8pt}
    {\popdisplay\fontsize{14}{14}\selectfont\color{popink}LEÇON #5}\par\vspace{4pt}
    {\raggedright\hyphenpenalty=10000\exhyphenpenalty=10000\popdisplay\fontsize{25}{27}\selectfont\color{popcream}\MakeUppercase{#1}\par}
    \vspace{8pt}
    {\popxbold\fontsize{9.5}{9.5}\selectfont\color{popink}\MakeUppercase{Durée conseillée : #6}}
  \end{tcolorbox}
  \begin{tcolorbox}[enhanced,colback=white,colframe=popink,boxrule=2.2pt,arc=10pt,
      drop shadow=popink,left=16pt,right=16pt,top=10pt,bottom=10pt,after skip=8pt]
    \poppill{Dans cette leçon}{poppurple}{white}\par\vspace{5pt}
    \popsemi\fontsize{9.3}{12.6}\selectfont\color{popink}#7
  \end{tcolorbox}
  \noindent\begin{minipage}[t]{0.487\linewidth}
    \begin{tcolorbox}[enhanced,colback=popgreen,colframe=popink,boxrule=2.2pt,arc=10pt,
        drop shadow=popink,left=11pt,right=11pt,top=10pt,bottom=10pt,height=3.1cm,valign=center]
      \tcbox[colback=white,colframe=popink,boxrule=1.3pt,arc=5pt,boxsep=0pt,nobeforeafter,
        left=3pt,right=3pt,top=3pt,bottom=3pt,valign=center]{\qrcode[height=1.9cm]{https://play.google.com/store/apps/details?id=com.luvvix.onbuch&hl=fr}}%
      \hspace{8pt}\begin{minipage}[c]{0.5\linewidth}
        {\popdisplay\fontsize{11}{12.5}\selectfont\color{white}\MakeUppercase{Télécharge OnBuch}}\par\vspace{4pt}
        {\raggedright\popsemi\fontsize{8.4}{11}\selectfont\color{white}Gratuit \textbullet\ Google Play\\Tuteur IA, annales, concours, résultats\par}
      \end{minipage}
    \end{tcolorbox}
  \end{minipage}\hfill
  \begin{minipage}[t]{0.487\linewidth}
    \begin{tcolorbox}[enhanced,colback=poppurple,colframe=popink,boxrule=2.2pt,arc=10pt,
        drop shadow=popink,left=11pt,right=11pt,top=10pt,bottom=10pt,height=3.1cm,valign=center]
      \tikz[baseline=-0.7ex]\node[circle,fill=white,draw=popink,line width=1.3pt,minimum size=1.6cm,inner sep=0]{\popdisplay\fontsize{24}{24}\selectfont\color{poppurple}+};%
      \hspace{8pt}\begin{minipage}[c]{0.6\linewidth}
        {\popdisplay\fontsize{11}{12.5}\selectfont\color{white}\MakeUppercase{Passe à OnBuch+}}\par\vspace{4pt}
        {\raggedright\popsemi\fontsize{8.4}{11}\selectfont\color{white}Tous les cours signés, Tuteur IA illimité, fiches PDF — onglet OnBuch+ de l'app\par}
      \end{minipage}
    \end{tcolorbox}
  \end{minipage}
  \vspace{8pt}
  \popcontenu
  \vfill
  \signatureonbuch
  \restoregeometry
  \newpage
}

% Rangée « Ce cours contient » (page de garde)
\newcommand{\poptile}[3]{% #1 couleur #2 gros texte #3 légende
  \begin{tcolorbox}[enhanced,colback=#1,colframe=popink,boxrule=2pt,arc=9pt,shadow={2.5pt}{-2.5pt}{0pt}{popink},
      left=6pt,right=6pt,top=9pt,bottom=9pt,halign=center,valign=center,height=1.75cm,width=\linewidth,nobeforeafter]
    {\popdisplay\fontsize{12.5}{13}\selectfont\color{white}\MakeUppercase{#2}}\par\vspace{3pt}
    {\centering\popsemi\fontsize{7.8}{9.5}\selectfont\color{white}#3\par}
  \end{tcolorbox}}
\newcommand{\popcontenu}{%
  \par\noindent{\popxboldls\fontsize{7.5}{7.5}\selectfont\color{popmuted}CE COURS SIGNÉ CONTIENT}\par\vspace{7pt}
  \noindent\begin{minipage}[t]{0.235\linewidth}\poptile{popblue}{Cours}{complet, illustré, pas à pas}\end{minipage}\hfill
  \begin{minipage}[t]{0.235\linewidth}\poptile{poporange}{Méthodes}{et exercices résolus}\end{minipage}\hfill
  \begin{minipage}[t]{0.235\linewidth}\poptile{poppink}{Exercices}{type examen + corrigés}\end{minipage}\hfill
  \begin{minipage}[t]{0.235\linewidth}\poptile{popgreen}{Fiche bilan}{l'essentiel à retenir}\end{minipage}\par}

% Sommaire
\newtcolorbox{sommaire}{enhanced,colback=white,colframe=popink,boxrule=2.2pt,arc=10pt,
  drop shadow=popink,left=18pt,right=18pt,top=13pt,bottom=13pt,before skip=6pt,after skip=20pt,
  before upper={\poppill{Sommaire}{poppurple}{white}\par\vspace{9pt}\popsemi\fontsize{10.6}{14.5}\selectfont\color{popink}}}

% Signature de fin de cours — poussée tout en bas de la dernière page
\newcommand{\findecours}{%
  \par\vfill\nopagebreak
  \begin{center}\popsquiggle{poporange}\end{center}\vspace{4pt}
  \signatureonbuch
}
\AtBeginDocument{\def\llbracket{\mathopen{[\mkern-3.5mu[}}\def\rrbracket{\mathclose{]\mkern-3.5mu]}}} % intervalles d'entiers (glyphe ⟦ absent de la police)
% Glyphes absents de Fira Math : redéfinis à partir de glyphes disponibles
\AtBeginDocument{%
  \def\setminus{\mathbin{\backslash}}\let\smallsetminus\setminus
  \def\vdots{\mathord{\vbox{\baselineskip3.2pt\lineskiplimit0pt\kern4pt\hbox{.}\hbox{.}\hbox{.}}}}%
  \def\ddots{\mathinner{\mkern1mu\raise6.5pt\hbox{.}\mkern2mu\raise3.6pt\hbox{.}\mkern2mu\raise0.7pt\hbox{.}\mkern1mu}}%
  \def\triangle{\mathord{\scalebox{0.9}{$\symup{\Delta}$}}}%
  \def\nabla{\mathord{\raisebox{\depth}{\scalebox{1}[-1]{$\symup{\Delta}$}}}}%
  \def\checkmark{\popcheck{popdark}}%
  \def\star{\mathbin{\ast}}\def\bigstar{\mathord{\ast}}%
  \def\diamond{\mathbin{\scalebox{0.6}{$\Diamond$}}}%
  \def\bigcup{\mathop{\mathchoice{\raisebox{-0.3em}{\scalebox{1.7}{$\cup$}}}{\scalebox{1.25}{$\cup$}}{\cup}{\cup}}\displaylimits}%
  \def\bigcap{\mathop{\mathchoice{\raisebox{-0.3em}{\scalebox{1.7}{$\cap$}}}{\scalebox{1.25}{$\cap$}}{\cap}{\cap}}\displaylimits}%
  \def\lesseqqgtr{\mathrel{\lessgtr}}\def\bigoplus{\mathop{\oplus}\displaylimits}%
}
\providecommand{\ssi}{si et seulement si} % abréviation fréquente
\AtBeginDocument{\providecommand{\wideparen}{\overparen}} % arc de cercle

%% FIGFIX-BEGIN v5 — correctifs globaux des figures (docs/onbuch-generation/tools/figfix)
\usepackage{adjustbox}\usepackage{etoolbox}\tcbuselibrary{raster}
% toute figure TikZ/pgfplots plus large que la ligne est réduite à la largeur disponible
\BeforeBeginEnvironment{tikzpicture}{\begin{adjustbox}{max width=\linewidth}}
\AfterEndEnvironment{tikzpicture}{\end{adjustbox}}
% légendes pgfplots lisibles par-dessus les courbes
\pgfplotsset{every axis/.append style={legend style={fill=white,fill opacity=0.92,text opacity=1,draw=black!15,inner sep=3pt}}}
% étiquettes d'arêtes (syntaxe edge["..."]) avec fond blanc
\tikzset{every edge quotes/.append style={fill=white,inner sep=1.2pt}}
%% FIGFIX-END
\begin{document}

\pagegarde{Lutte contre les comportements préjudiciables à la santé reproductive}{SVT}{1ère TI}{Module 2 : L'Éducation à la Santé}{NCT05}{3 h}{Cette leçon analyse les comportements à risque pour la santé reproductive des adolescents (rapports non protégés, grossesses précoces, avortements clandestins) et leurs conséquences sanitaires et sociales. Elle permet de maîtriser les moyens de prévention et de contraception pour adopter des comportements responsables.
\vspace{4pt}
\begin{multicols}{2}\small
\begin{itemize}
\item À la fin de cette leçon, je sais décrire les grandes étapes de la physiologie de la reproduction humaine (cycle ovarien, spermatogenèse, fécondation).
\item À la fin de cette leçon, je sais identifier les comportements à risque pour la santé reproductive dans une situation concrète (rapports non protégés, multiplicité des partenaires, consommation de substances).
\item À la fin de cette leçon, je sais expliquer les mécanismes physiologiques conduisant à une grossesse précoce et ses risques obstétricaux spécifiques (fistule, éclampsie, accouchement difficile).
\item À la fin de cette leçon, je sais analyser les conséquences sanitaires (IST, VIH/SIDA, stérilité, séquelles d'avortement) et sociales (décrochage scolaire, stigmatisation, pauvreté) des comportements à risque.
\item À la fin de cette leçon, je sais distinguer les principales méthodes de contraception (naturelles, barrières, hormonales, chirurgicales, d'urgence) selon leur mode d'action, leur efficacité (indice de Pearl) et leur accessibilité au Cameroun.
\item À la fin de cette leçon, je sais argumenter le choix d'une contraception adaptée à un profil d'adolescent(e) en tenant compte du double objectif : prévention grossesse ET prévention IST (double protection).
\end{itemize}\end{multicols}}

\begin{sommaire}
\begin{multicols}{2}
\begin{enumerate}
\item 1. Rappels physiologiques : les bases de la fertilité
\item 2. Comportements à risque : identification et déterminants
\item 3. Grossesses précoces : conséquences sanitaires et sociales
\item 4. Avortements clandestins : un drame de santé publique
\item 5. Moyens de contraception : panorama et choix responsable
\item 6. Stratégies de prévention et alternatives responsables
\item Activité d'intégration
\item Méthodes et exercices résolus
\item Exercices
\item Corrigés des exercices
\item Fiche bilan
\end{enumerate}
\end{multicols}
\end{sommaire}

\begin{objectifs}
\begin{itemize}
\item À la fin de cette leçon, je sais décrire les grandes étapes de la physiologie de la reproduction humaine (cycle ovarien, spermatogenèse, fécondation).
\item À la fin de cette leçon, je sais identifier les comportements à risque pour la santé reproductive dans une situation concrète (rapports non protégés, multiplicité des partenaires, consommation de substances).
\item À la fin de cette leçon, je sais expliquer les mécanismes physiologiques conduisant à une grossesse précoce et ses risques obstétricaux spécifiques (fistule, éclampsie, accouchement difficile).
\item À la fin de cette leçon, je sais analyser les conséquences sanitaires (IST, VIH/SIDA, stérilité, séquelles d'avortement) et sociales (décrochage scolaire, stigmatisation, pauvreté) des comportements à risque.
\item À la fin de cette leçon, je sais distinguer les principales méthodes de contraception (naturelles, barrières, hormonales, chirurgicales, d'urgence) selon leur mode d'action, leur efficacité (indice de Pearl) et leur accessibilité au Cameroun.
\item À la fin de cette leçon, je sais argumenter le choix d'une contraception adaptée à un profil d'adolescent(e) en tenant compte du double objectif : prévention grossesse ET prévention IST (double protection).
\item À la fin de cette leçon, je sais proposer des alternatives responsables (abstinence, fidélité, préservatif, planning familial) face à une situation à risque donnée.
\item À la fin de cette leçon, je sais localiser les structures de prise en charge au Cameroun (Centres de Santé des Jeunes, CEDA, ONG) pour un accompagnement confidentiel et gratuit.
\end{itemize}
\end{objectifs}

\begin{prerequis}
\begin{itemize}
\item Anatomie de l'appareil reproducteur mâle et femelle (Seconde / début 1ère).
\item Notions de base sur les IST/VIH/SIDA (Seconde).
\item Cycle menstruel et rôle des hormones ovariennes (FSH, LH, œstrogènes, progestérone).
\item Définition de la puberté et changements morphologiques associés.
\item Notions de citoyenneté et droits à la santé sexuelle et reproductive (Éducation à la vie).
\end{itemize}
\end{prerequis}

\newpage

\coursec{Rappels physiologiques : les bases de la fertilité}

Au quartier Mokolo à Yaoundé, Amina, 17 ans, élève en 1ère C, découvre qu'elle est enceinte après un rapport non protégé avec son petit ami. Paniquée, elle envisage un avortement clandestin chez une « tante » du quartier, ignorant les risques d'hémorragie, d'infection ou de stérilité définitive. Son amie Binta, mieux informée, tente de la convaincre de consulter le Centre de Santé des Jeunes (CSJ) pour un accompagnement psychologique et médical. Cette situation, hélas fréquente dans nos lycées, illustre l'urgence d'une éducation à la santé reproductive responsable. Pour comprendre comment une grossesse survient et comment l'éviter ou la planifier, nous devons d'abord maîtriser les bases physiologiques de la fertilité humaine.

\courssub{Anatomie fonctionnelle des appareils reproducteurs}

La reproduction humaine met en jeu deux appareils reproducteurs aux organisations distinctes mais complémentaires. Une bonne connaissance de leur anatomie est le prérequis indispensable pour comprendre la fécondation et l'action des moyens de contraception.

\begin{popfigure}
\begin{tikzpicture}[scale=0.7, every node/.style={font=\small}]
% Appareil reproducteur masculin (vue sagittale médiane)
\begin{scope}[xshift=-8cm]
\draw[popdark, thick] (0,0) ellipse (1.5 and 3); % Bassin
% Vessie
\draw[popblue, fill=popblueL, thick] (0,2.5) .. controls (1,3) and (1,3.5) .. (0,4) .. controls (-1,3.5) and (-1,3) .. cycle;
\node[popblue] at (0,3.2) {Vessie};
% Rectum
\draw[poporange, fill=poporangeL, thick] (0,-2.5) .. controls (-1,-3) and (-1,-3.5) .. (0,-4) .. controls (1,-3.5) and (1,-3) .. cycle;
\node[poporange] at (0,-3.2) {Rectum};
% Vésicules séminales
\draw[poppurple, fill=poppurpleL, thick] (-0.8,0.8) ellipse (0.6 and 0.4);
\draw[poppurple, fill=poppurpleL, thick] (0.8,0.8) ellipse (0.6 and 0.4);
\node[poppurple] at (0,0.8) {Vésicules séminales};
% Prostate
\draw[poppurple!70!popdark, fill=poppurpleL!70!popdark, thick] (0,0) ellipse (0.8 and 0.5);
\node[poppurple!70!popdark] at (0,0) {Prostate};
% Épididymes
\draw[popgreen!70!popdark, fill=popgreenL, thick] (-1.2,-1.5) ellipse (0.3 and 0.6);
\draw[popgreen!70!popdark, fill=popgreenL, thick] (1.2,-1.5) ellipse (0.3 and 0.6);
\node[popgreen!70!popdark] at (-1.2,-1.5) {Épididyme};
\node[popgreen!70!popdark] at (1.2,-1.5) {Épididyme};
% Canaux déférents (remontent derrière la vessie)
\draw[popdark, thick] (-1.2,-0.9) .. controls (-1.5,-0.5) and (-1.5,1.5) .. (-1,2.2);
\draw[popdark, thick] (1.2,-0.9) .. controls (1.5,-0.5) and (1.5,1.5) .. (1,2.2);
% Testicules
\draw[popgreen, fill=popgreenL, thick] (-1.2,-2.8) ellipse (0.5 and 0.8);
\draw[popgreen, fill=popgreenL, thick] (1.2,-2.8) ellipse (0.5 and 0.8);
\node[popgreen] at (-1.2,-2.8) {Testicule};
\node[popgreen] at (1.2,-2.8) {Testicule};
% Pénis / Urethrè
\draw[popdark, thick] (0,-0.5) -- (0,-1.5) .. controls (0,-2) and (0.5,-2.5) .. (0.5,-3);
\draw[popdark, thick] (0,-0.5) -- (0,-1.5) .. controls (0,-2) and (-0.5,-2.5) .. (-0.5,-3);
\node[popdark] at (0.8,-2) {Uréthre};
\node[popdark] at (-0.8,-2) {Pénis};
% Légendes fléchées
\draw[->, popline] (-2.8, 2.5) -- (-1.5, 2) node[left] {Vessie};
\draw[->, popline] (-2.8, 0.8) -- (-1.5, 0.8) node[left] {Vésicules séminales};
\draw[->, popline] (-2.8, 0) -- (-1.5, 0) node[left] {Prostate};
\draw[->, popline] (-2.8, -1.5) -- (-1.5, -1.5) node[left] {Épididyme (maturation)}; 
\draw[->, popline] (-2.8, -2.8) -- (-1.5, -2.8) node[left] {Testicule (spermatogenèse)};
\draw[->, popline] (-2.8, -3.8) -- (-1.2, -3.5) node[left] {Canal déférent};
\node[align=center, font=\bfseries\small] at (0, 4.5) {Appareil reproducteur masculin};
\end{scope}

% Appareil reproducteur féminin (vue sagittale médiane)
\begin{scope}[xshift=8cm]
\draw[popdark, thick] (0,0) ellipse (1.5 and 3); % Bassin
% Vessie
\draw[popblue, fill=popblueL, thick] (0,2.5) .. controls (1,3) and (1,3.5) .. (0,4) .. controls (-1,3.5) and (-1,3) .. cycle;
\node[popblue] at (0,3.2) {Vessie};
% Rectum
\draw[poporange, fill=poporangeL, thick] (0,-2.5) .. controls (-1,-3) and (-1,-3.5) .. (0,-4) .. controls (1,-3.5) and (1,-3) .. cycle;
\node[poporange] at (0,-3.2) {Rectum};
% Utérus (corps + col)
\draw[poppink, fill=poppinkL, thick] (-0.5,0) -- (-0.5,1.5) .. controls (-0.5,2) and (0.5,2) .. (0.5,1.5) -- (0.5,0);
\node[poppink] at (0,0.7) {Corps de l'utérus};
% Col de l'utérus
\draw[poppink!70!popdark, fill=poppinkL!70!popdark, thick] (-0.5,0) -- (-0.5,-1) -- (0.5,-1) -- (0.5,0);
\node[poppink!70!popdark] at (0,-0.5) {Col de l'utérus};
% Endomètre (muqueuse utérine)
\draw[poppink!80!popdark, thick, decorate, decoration={snake, amplitude=1pt, segment length=3pt}] (-0.4,0.1) -- (-0.4,1.4);
\draw[poppink!80!popdark, thick, decorate, decoration={snake, amplitude=1pt, segment length=3pt}] (0.4,0.1) -- (0.4,1.4);
\node[poppink!80!popdark] at (1.3,0.7) {Endomètre};
% Trompes
\draw[poppurple, thick] (-0.5,1.5) .. controls (-1.5,1.8) and (-2,1.5) .. (-2.2,1);
\draw[poppurple, thick] (0.5,1.5) .. controls (1.5,1.8) and (2,1.5) .. (2.2,1);
\node[poppurple] at (-2.6,1.2) {Trompe utérine};
\node[poppurple] at (2.6,1.2) {Trompe utérine};
% Ovaires
\draw[popgreen, fill=popgreenL, thick] (-2.2,1) ellipse (0.4 and 0.6);
\draw[popgreen, fill=popgreenL, thick] (2.2,1) ellipse (0.4 and 0.6);
\node[popgreen] at (-2.2,1) {Ovaire};
\node[popgreen] at (2.2,1) {Ovaire};
% Vagin
\draw[popdark, thick] (-0.5,-1) -- (-0.5,-2.5);
\draw[popdark, thick] (0.5,-1) -- (0.5,-2.5);
\node[popdark] at (0,-2.8) {Vagin};
% Légendes fléchées
\draw[->, popline] (2.8, 2.5) -- (1.5, 2) node[right] {Vessie};
\draw[->, popline] (2.8, 1.5) -- (1.5, 1.5) node[right] {Trompe (site fécondation)};
\draw[->, popline] (2.8, 0.7) -- (1.5, 0.7) node[right] {Utérus (nidation)};
\draw[->, popline] (2.8, -0.5) -- (1.5, -0.5) node[right] {Col de l'utérus};
\draw[->, popline] (2.8, -2) -- (1.5, -2) node[right] {Vagin};
\node[align=center, font=\bfseries\small] at (0, 4.5) {Appareil reproducteur féminin};
\end{scope}
\end{tikzpicture}
\legende{Schéma simplifié des appareils reproducteurs masculin (gauche) et féminin (droite) en vue sagittale médiane. Repérer les sites de production des gamètes (testicules, ovaires), de la maturation spermatique (épididyme), de la fécondation (ampoule de la trompe) et de la nidation (endomètre utérin).}
\end{popfigure}

\begin{definition}[Gamète et Gonade]
Une \cle{gonade} est un organe producteur de gamètes et d'hormones sexuelles stéroïdiennes. Chez l'homme, ce sont les \cle{testicules} (produisant des spermatozoïdes et de la testostérone) ; chez la femme, les \cle{ovaires} (produisant des ovocytes et des œstrogènes/progestérone). Le \cle{gamète} mâle est le spermatozoïde (haploïde, 23 chromosomes), le gamète femelle est l'ovocyte II (haploïde, 23 chromosomes, bloqué en métaphase II), souvent appelé « ovule » par abus de langage.
\end{definition}

\courssub{La spermatogenèse : une production continue et massive}

Contrairement à la femme qui naît avec son stock d'ovocytes, l'homme produit ses gamètes \emph{tout au long de sa vie} à partir de la puberté.

\begin{experience}[Observation histologique d'un tubule séminifère (document microscopique)]
\textbf{Matériel :} Lame mince de testicule humain (coloration Hématoxyline-Éosine), microscope optique (objectif ×40).
\textbf{Observation :} On observe des tubules séminifères bordés de cellules de Sertoli (noyau clair, volumineux, nucléole apparent). Vers la lumière du tubule, les stades de maturation s'enchaînent de la périphérie vers le centre : \textbf{spermatogonies} (cellules souches, bord externe) $\rightarrow$ \textbf{spermatocytes I} (gros noyaux, prophase de méiose I, chromosomes visibles) $\rightarrow$ \textbf{spermatocytes II} (rares, petite taille) $\rightarrow$ \textbf{spermatides} (petites, rondes, noyau dense) $\rightarrow$ \textbf{spermatozoïdes} (têtes pointées vers la lumière, flagelles dans la lumière). Entre les tubules, dans le tissu conjonctif, les \textbf{cellules de Leydig} (cytoplasme éosinophile abondant, noyau central) sont visibles.
\textbf{Interprétation :} La spermatogenèse est un processus continu de divisions mitotiques (spermatogonies) et méiotiques (spermatocytes I $\rightarrow$ II $\rightarrow$ spermatides) suivi de la \cle{spermiogenèse} (différenciation cellulaire sans division : condensation nucléaire, formation de l'acrosome et du flagelle). Les cellules de Sertoli assurent la nutrition, la protection (barrière hémato-testiculaire) et la sécrétion de l'hormone anti-müllérienne (AMH) et d'inhibine (rétrocontrôle négatif sur FSH). Les cellules de Leydig sécrètent la testostérone sous l'effet de la LH.
\end{experience}

\begin{propriete}[Caractéristiques de la spermatogenèse]
\begin{itemize}
    \item \textbf{Durée :} Environ \textbf{72 à 74 jours} (2 mois et demi) pour former un spermatozoïde mature à partir d'une spermatogonie souche.
    \item \textbf{Régulation hormonale :} La \textbf{FSH} agit sur les cellules de Sertoli (soutien de la méiose, sécrétion ABP -- protéine liant les androgènes). La \textbf{LH} stimule les cellules de Leydig $\rightarrow$ sécrétion de \textbf{testostérone}. La testostérone (via ABP) maintient la spermatogenèse et exerce un rétrocontrôle négatif sur l'axe hypothalamo-hypophysaire (GnRH, LH). L'inhibine (Sertoli) rétrocontrôle spécifiquement la FSH.
    \item \textbf{Production :} Continue dès la puberté (vers 12-14 ans), sans cyclicité. Un homme adulte produit \textbf{100 à 300 millions de spermatozoïdes par jour} (environ 1000 par battement cardiaque).
    \item \textbf{Maturation finale (épididyme) :} Se termine dans l'\cle{épididyme} (acquisition de la motilité progressive, du pouvoir fécondant et de la capacité de reconnaissance de la zone pellucide) en \textbf{$\sim$12 jours}. Le spermatozoïde mature est stocké dans la queue de l'épididyme et le canal déférent.
\end{itemize}
\end{propriete}

\courssub{L'oogenèse et le cycle ovarien : une cyclicité hormonale précise}

Chez la femme, la production des gamètes (oogenèse) est achevée avant la naissance (stock d'ovogonies $\rightarrow$ ovocytes I bloqués en prophase I). Le cycle ovarien, qui débute à la puberté (ménarche), est une succession cyclique d'événements ovariens et utérins pilotée par l'axe hypothalamo-hypophysaire.

\begin{experience}[Analyse des variations hormonales au cours d'un cycle ovarien de 28 jours (données de référence)]
\textbf{Données :} Dosages sanguins quotidiens de FSH, LH, Œstradiol (E2), Progestérone (P4) chez une femme à cycles réguliers.
\textbf{Observations (résumées) :}
\begin{enumerate}
    \item \textbf{J1-J5 (Règles / Ménstrues) :} Taux bas de toutes les hormones (FSH, LH, E2, P4). L'endomètre fonctionnel se desquame sous l'effet de la chute brutale de la progestérone du cycle précédent.
    \item \textbf{J5-J13 (Phase folliculaire / Pré-ovulatoire) :} \textbf{FSH} augmente légèrement $\rightarrow$ recrute une cohorte de follicules antraux. Un follicule devient \cle{dominant} (récepteurs FSH + LH, vascularisation). Il sécrète massivement des \textbf{œstrogènes (E2)} (via cellules de la thèque $\rightarrow$ androgènes $\rightarrow$ cellules de la granulose $\rightarrow$ aromatase $\rightarrow$ œstrogènes). Le pic d'œstrogènes (J12-J13) exerce un \emph{feed-back positif} sur l'hypophyse antérieure (sensibilisation à la GnRH).
    \item \textbf{J14 (Ovulation) :} \textbf{Pic massif de LH} (et FSH) déclenché par le feed-back positif des œstrogènes. Rupture du follicule de Graaf $\rightarrow$ libération de l'ovocyte II entouré de sa couronne radiée (cellules de la granulose) et de la zone pellucide. L'ovocyte est capté par les franges de la trompe.
    \item \textbf{J15-J28 (Phase lutéale / Post-ovulatoire) :} Le follicule vidé se luteinise $\rightarrow$ \cle{corps jaune}. Il sécrète massivement de la \textbf{progestérone (P4)} et des œstrogènes. Ces hormones inhibent FSH et LH (feed-back négatif fort) $\rightarrow$ atrésie des autres follicules. Si pas de grossesse, le corps jaune régresse (J24-J28) $\rightarrow$ chute brutale de P4 et E2 $\rightarrow$ ischémie endométriale $\rightarrow$ règles (J1 du cycle suivant).
\end{enumerate}
\textbf{Conclusion :} Le cycle est orchestré par une boucle de rétroaction entre ovaire et hypophyse. La phase folliculaire est variable en durée ; la phase lutéale est constante (\textbf{$\sim$14 jours}).
\end{experience}

\begin{popfigure}
\begin{tikzpicture}
\begin{axis}[
    width=\linewidth, height=8cm,
    xmin=0, xmax=29,
    ymin=0, ymax=10,
    axis lines=middle,
    xlabel={Jour du cycle},
    ylabel={Concentration (UI/L ou ng/mL) / Épaisseur (mm)},
    xtick={0,7,14,21,28},
    xticklabels={0, 7, 14, 21, 28},
    ytick=\empty,
    legend style={at={(0.02,0.98)}, anchor=north west, font=\small, fill=white, draw=popline},
    clip=false,
    domain=0:28, samples=300,
    /pgfplots/x coord trafo/.code={\pgfmathparse{#1}\pgfmathresult},
    /pgfplots/y coord trafo/.code={\pgfmathparse{#1}\pgfmathresult},
]
% FSH
\addplot[popblue, thick, name path=FSH] {1.5 + 0.5*sin(deg(x*180/14)) + 0.8*exp(-(x-14)*(x-14)/8)};
\addlegendentry{FSH}
% LH
\addplot[poporange, thick, name path=LH] {1 + 6*exp(-(x-14)*(x-14)/4)};
\addlegendentry{LH}
% Œstrogènes
\addplot[poppink, thick, name path=E2] {0.5 + 3.5*exp(-(x-13)*(x-13)/10) + 1.5*exp(-(x-21)*(x-21)/15)};
\addlegendentry{Œstradiol (E2)}
% Progestérone
\addplot[poppurple, thick, name path=P4] {0.2 + 4*exp(-(x-21)*(x-21)/20)};
\addlegendentry{Progestérone (P4)}
% Endomètre (échelle secondaire approximative, tracé qualitatif)
\addplot[popgreen, thick, dashed, name path=Endo] {1 + 0.25*x - 0.005*x*x + 1.5*exp(-(x-21)*(x-21)/15) - 2*exp(-(x-28)*(x-28)/2)};
\addlegendentry{Endomètre (mm)}

% Annotations événements ovariens
\node[align=center, font=\small, popgreen] at (axis cs:10,9.5) {Follicule\\dominant};
\node[align=center, font=\small\bfseries, poporange] at (axis cs:14,9.5) {OVULATION};
\node[align=center, font=\small, popgold] at (axis cs:21,9.5) {Corps\\jaune};
\node[align=center, font=\small\bfseries, popred] at (axis cs:0,9.5) {RÈGLES};
\node[align=center, font=\small\bfseries, popred] at (axis cs:28,9.5) {RÈGLES};

% Lignes pointillées phases
\draw[dashed, popmuted] (axis cs:14,0) -- (axis cs:14,7);
\node[rotate=90, popmuted, font=\small] at (axis cs:7,8.5) {Phase folliculaire (variable)};
\node[rotate=90, popmuted, font=\small] at (axis cs:21,8.5) {Phase lutéale (constante $\approx$ 14 j)};
\end{axis}
\end{tikzpicture}
\legende{Profil hormonal synchronisé du cycle ovarien de 28 jours (simulation pgfplots). Noter le pic de LH déclencheur de l'ovulation (J14), l'inversion des profils œstrogènes/progestérone entre les deux phases, et la constance de la phase lutéale (14 jours). L'endomètre prolifère sous l'effet des œstrogènes (phase folliculaire), puis se sécrétoire sous l'effet de la progestérone (phase lutéale).}
\end{popfigure}

\begin{definition}[Ovulation]
L'\cle{ovulation} est la libération d'un \cle{ovocyte II} (stade méiotique bloqué en métaphase II) par rupture du \cle{follicule de Graaf} sous l'effet du \cle{pic LH} (surge LH). L'ovocyte est capté par les franges de la trompe (fimbriae). Il ne termine sa méiose II (expulsion du 2ème globule polaire) qu'en cas de fécondation par un spermatozoïde.
\end{definition}

\begin{propriete}[Constance de la phase lutéale]
La date d'ovulation est variable (J14 $\pm$ 2j en cycle 28j, mais peut varier de J10 à J20 selon les femmes et les cycles) ; \textbf{seule la phase lutéale est constante : elle dure 14 jours} (entre l'ovulation et le premier jour des règles suivantes). C'est la base du calcul rétrospectif de la date d'ovulation : \textit{Jour d'ovulation = Durée du cycle - 14}.
\end{propriete}

\courssub{Fécondation, fenêtre de fécondité et nidation}

La fécondation est la fusion d'un spermatozoïde et d'un ovocyte II pour former un zygote (1 cellule, 2n = 46 chromosomes). Elle a lieu dans l'\cle{ampoule de la trompe utérine}.

\begin{methode}[Déroulement de la fécondation et début de grossesse]
\begin{enumerate}
    \item \textbf{Dépôt et capacitation :} Lors de l'éjaculation (volume 3-5 mL, concentration 15-200 M/mL, total 100-300 M spz), les spermatozoïdes subissent la \cle{capacitation} dans l'appareil génital féminin (perte de protéines décapacitantes séminales, modification de la membrane plasmique, acquisition motilité hyperactive) : durée \textbf{$\sim$5-7 h}.
    \item \textbf{Migration et sélection :} Seuls $\sim$1\% franchissent le col de l'utérus (glaires cervicales perméables et filantes sous l'effet des œstrogènes pré-ovulatoires). Migration utéro-tubaire guidée par chimiotactisme (facteurs sécrétés par l'ovocyte/cellules de la granulose) et contractions utérines.
    \item \textbf{Rencontre dans l'ampoule :} $\sim$100-200 spz entourent l'ovocyte. Un seul franchit la \cle{zone pellucide} (réaction acrosomique $\rightarrow$ lyse enzymatique de la ZP3 $\rightarrow$ pénétration) $\rightarrow$ déclenchement de la \cle{réaction corticale} (exocytose des granules corticaux $\rightarrow$ durcissement de la ZP $\rightarrow$ blocage de la polyspermie).
    \item \textbf{Fusion et syngamie :} Fusion des membranes plasmatiques $\rightarrow$ entrée du noyau spermatique + centriole paternel dans l'ooplasme. Fin de méiose II de l'ovocyte $\rightarrow$ expulsion du 2ème globule polaire. Formation des pronucléus mâle et femelle $\rightarrow$ migration, apposition (syngamie) $\rightarrow$ formation du \cle{zygote} (2n, 46 chromosomes, 2 centrosomes paternels).
    \item \textbf{Segmentation et migration tubaire :} Divisions mitotiques successives (clivages) sans croissance pendant la migration vers l'utérus (3-4 jours). Stade \cle{morula} (J3, $\sim$16-32 blastomères compacts) $\rightarrow$ \cle{blastocyste} (J5-J6, $\sim$100-150 cellules : \textbf{trophoblaste} périphérique + \textbf{masse cellulaire interne} / embryoblaste). Zona pellucide lysée (hatching).
    \item \textbf{Nidation (J6-J7 post-fécondation) :} Le blastocyste « éclot » (hatching) et s'implante par le pôle trophoblastique dans l'endomètre \textbf{progestéronisé} (phase lutéale, fenêtre d'implantation). Le trophoblaste sécrète la \cle{$\beta$-hCG} (hormone chorionique gonadotrope, détectable dans le sang maternel $\sim$J8-J9, dans les urines $\sim$J14) qui « sauve » le corps jaune $\rightarrow$ maintien de la sécrétion de progestérone $\rightarrow$ pas de règles (aménorrhée).
\end{enumerate}
\end{methode}

\begin{popfigure}
\begin{tikzpicture}[scale=0.7]
% Chronogramme fenêtre de fécondité
\draw[->, thick] (0,0) -- (14,0) node[right] {Jours relatifs à l'ovulation (J0)};
\foreach \x in {-5,-4,-3,-2,-1,0,1,2} \draw (\x+7,0.1) -- (\x+7,-0.1) node[below] {\x};
% Survie spermatozoïdes (barre horizontale)
\draw[popblue, line width=6pt, opacity=0.6] (2, 2) -- (7, 2); % J-5 à J0
\node[popblue, left, align=right] at (1.5, 2) {Survie spermatozoïdes\\dans glaires cervicales (3-5 j)};
% Survie ovocyte (barre courte)
\draw[poppink, line width=6pt, opacity=0.6] (7, 1) -- (8, 1); % J0 à J+1
\node[poppink, left, align=right] at (1.5, 1) {Survie ovocyte\\après ovulation (12-24 h)};
% Fenêtre fertile (zone ombrée)
\draw[popgreen!50!popdark, fill=popgreenL, opacity=0.3] (2, 0.5) rectangle (8, 2.5);
\node[popgreen!50!popdark, align=center, font=\small\bfseries] at (5, 3.3) {FENÊTRE DE FÉCONDITÉ THÉORIQUE\\J-5 à J+1 (6 jours)};
% Flèches rapports
\draw[->, thick, popdark] (3, 3) -- (3, 2.5) node[above, align=center, font=\small] {Rapport J-4};
\draw[->, thick, popdark] (7.5, 3) -- (7.5, 2.5) node[above, align=center, font=\small] {Rapport J0};
\draw[->, thick, popdark, dashed] (9, 3) -- (9, 2.5) node[above, align=center, font=\small] {Rapport J+2 (trop tard)};
% Légendes
\node[align=left, font=\small] at (11, 2) {
\textbf{Clé :}\\
$\bullet$ Spz déposés J-5 $\rightarrow$ vivants J0 (ovulation).\\
$\bullet$ Spz déposés J0 $\rightarrow$ fécondent ovocyte.\\
$\bullet$ Spz déposés J+1 $\rightarrow$ ovocyte déjà mort.\\
$\bullet$ Spz déposés J-6 $\rightarrow$ morts avant ovulation.
};
\end{tikzpicture}
\legende{Chronogramme de la fenêtre de fécondité. Les spermatozoïdes survivent 3 à 5 jours dans les voies génitales féminines (glaires cervicales favorables) ; l'ovocyte ne survit que 12 à 24 h après l'ovulation. La période fertile s'étend donc de \textbf{5 jours avant} à \textbf{1 jour après} l'ovulation.}
\end{popfigure}

\begin{methode}[Calcul de la période fertile théorique (Méthode Ogino-Knaus)]
Cette méthode statistique, basée sur la constance de la phase lutéale (14 j) et la survie maximale des spermatozoïdes (5 j), permet d'estimer la période fertile \textbf{a posteriori} sur les 6 à 12 derniers cycles observés.
\begin{enumerate}
    \item Noter la durée du \textbf{cycle le plus court} ($D_{min}$) et du \textbf{cycle le plus long} ($D_{max}$) observés récemment.
    \item \textbf{Premier jour fertile} = $D_{min} - 18$ (ex: cycle court 26j $\rightarrow$ J8). \textit{Note : correspond à (ovulation la plus précoce) - 5j survie spz = ($D_{min}-14$)-5 = $D_{min}-19$ ; la formule scolaire retient souvent -18 par marge de sécurité.}
    \item \textbf{Dernier jour fertile} = $D_{max} - 11$ (ex: cycle long 30j $\rightarrow$ J19). \textit{Note : correspond à (ovulation la plus tardive) + 1j survie ovocyte = ($D_{max}-14$)+1 = $D_{max}-13$ ; la formule scolaire retient souvent -11 par marge de sécurité.}
    \item La période fertile théorique (période « à risque ») s'étend de \textbf{J($D_{min}-18$) à J($D_{max}-11$)}.
\end{enumerate}
\end{methode}

\begin{attention}[Fiabilité de la méthode Ogino-Knaus chez l'adolescente]
\textbf{TRÈS FAIBLE (Indice de Pearl > 20).} Chez l'adolescente, les cycles sont souvent \textbf{irréguliers} (anovulatoires, courts, longs) pendant 2 à 3 ans après la ménarche. $D_{min}$ et $D_{max}$ fluctuent fortement d'un cycle à l'autre, rendant le calcul imprécis et la période « à risque » très large (souvent > 15 jours). \textbf{Ne jamais utiliser cette méthode comme unique contraception.} Elle ne protège pas non plus des IST/VIH.
\end{attention}

\begin{attention}[Erreur fréquente : « On ne peut pas tomber enceinte pendant les règles »]
\textbf{FAUX.} C'est possible si :
\begin{itemize}
    \item Cycle court (ex: 21 jours) $\rightarrow$ Ovulation vers J7.
    \item Règles longues (6-7 jours) $\rightarrow$ Rapport en fin de règles (J6-J7).
    \item Survie spermatozoïdes 5 jours $\rightarrow$ Spz vivants au moment de l'ovulation (J7).
\end{itemize}
\textbf{Conclusion :} Aucune période du cycle n'est « sans risque » absolu sans contraception efficace.
\end{attention}

\begin{savaistu}[Probabilité de grossesse par cycle (Fécondabilité)]
Un seul rapport sexuel non protégé pendant la fenêtre de fécondité suffit pour engrosser. La \cle{fécondabilité} mensuelle (probabilité de grossesse par cycle pour un couple fertile ayant des rapports réguliers non protégés) est d'environ \textbf{20 à 25 \%}. Elle diminue avec l'âge (surtout après 35 ans chez la femme) mais reste \textbf{maximale chez l'adolescente} (cycles très fertiles malgré l'irrégularité, qualité ovocytaire optimale).
\end{savaistu}

\courssub{Particularités de la fertilité chez l'adolescente}

La puberté s'installe par l'activation de l'axe gonadotrope (GnRH pulsatile $\rightarrow$ FSH/LH $\rightarrow$ stéroïdes). Chez la jeune fille camerounaise, la ménarche survient en moyenne vers \textbf{12-13 ans} (plus précoce en zone urbaine, liée à l'amélioration de l'état nutritionnel et à l'exposition aux perturbateurs endocriniens).

\begin{propriete}[Fertilité adolescente : forte mais imprévisible]
\begin{itemize}
    \item \textbf{Cycles anovulatoires fréquents :} 50-80\% des cycles les 2 premières années post-ménarche. Absence de pic LH $\rightarrow$ pas de corps jaune $\rightarrow$ pas de progestérone $\rightarrow$ règles irrégulières, souvent abondantes et prolongées (ménorragies par hyperœstrogénie relative).
    \item \textbf{Ovulations précoces possibles :} Même avant la régularisation, une ovulation peut survenir \textbf{avant les premières règles} (première ovulation silencieuse) ou pendant une aménorrhée. Une adolescente peut tomber enceinte \textbf{avant sa première ménstruation apparente}.
    \item \textbf{Fécondité maximale :} Qualité ovocytaire optimale (faible taux d'anéuploïdies), endomètre réceptif. Risque de grossesse par rapport non protégé \textbf{plus élevé} que chez la femme de 25-30 ans.
    \item \textbf{Immunité cervicale immature :} Le col de l'utérus (zone de transformation jonctionnelle ecto-endocervicale) est étendu et exposé (ectopie / « érosion » physiologique), plus vulnérable aux IST (HPV oncogènes, \emph{Chlamydia trachomatis}, \emph{Neisseria gonorrhoeae}, VIH), augmentant le risque de complications ultérieures (salpingites, stérilité tubaire, cancer du col de l'utérus).
\end{itemize}
\end{propriete}

\begin{exoresolu}[Détermination de la période fertile chez une lycéenne]
\textbf{Énoncé :} Nadine, élève en 1ère C à Douala, note ses cycles sur 6 mois : 28 j, 32 j, 26 j, 30 j, 29 j, 31 j. Elle a eu un rapport non protégé le J12 de son cycle actuel (qui a débuté le 1er mars). Est-elle dans sa période fertile théorique ? Que lui conseillez-vous ?
\textbf{Solution détaillée :}
\begin{enumerate}
    \item \textbf{Identifier $D_{min}$ et $D_{max}$ :} $D_{min} = 26$ j ; $D_{max} = 32$ j.
    \item \textbf{Calculer les bornes (Méthode Ogino-Knaus standard) :}
    \begin{align*}
        \text{Premier jour fertile} &= D_{min} - 18 = 26 - 18 = \textbf{J8} \\
        \text{Dernier jour fertile} &= D_{max} - 11 = 32 - 11 = \textbf{J21}
    \end{align*}
    \item \textbf{Période fertile théorique :} \textbf{Du J8 au J21 inclus} (14 jours de « risque »).
    \item \textbf{Situation de Nadine :} Rapport le J12. \textbf{OUI, elle est en pleine période fertile théorique.}
    \item \textbf{Conseil médical urgent :} Risque de grossesse réel et élevé. Orientation immédiate vers le Centre de Santé des Jeunes (CSJ) ou une formation sanitaire pour :
    \begin{itemize}
        \item Contraception d'urgence (CU) : Pilule à la lévonorgestrel (1,5 mg en dose unique, efficace jusqu'à \textbf{72 h}, idéalement < 12 h) ou Acétate d'ulipristal (30 mg, efficace jusqu'à \textbf{120 h} / 5 jours). La CU n'est pas abortive (retarde l'ovulation).
        \item Dépistage IST (VIH, Hépatites B/C, Syphilis, Chlamydiae) et prophylaxie post-exposition (PEP) si risque VIH < 48-72h.
        \item Prescription d'une contraception moderne adaptée (pilule œstroprogestative, implant sous-cutané, DIU au cuivre ou hormonal, préservatif masculin/féminin \textbf{obligatoire} en double protection).
    \end{itemize}
    \item \textbf{Rappel pédagogique :} La méthode Ogino est peu fiable seule chez l'adolescente ; elle ne doit servir qu'à comprendre sa physiologie, pas à se protéger.
\end{enumerate}
\end{exoresolu}

\begin{aretenir}[Bilan : Les bases de la fertilité]
\begin{itemize}
    \item \textbf{Homme :} Spermatogenèse continue (72 j), massive (100-300 M/j), dès puberté, pilotée par Testostérone (LH/Leydig) + FSH (Sertoli/ABP). Maturation épididymaire (12 j).
    \item \textbf{Femme :} Stock ovocytaire fini (ovogenèse fœtale). Cycle ovarien cyclique : FSH $\rightarrow$ Follicule dominant $\rightarrow$ Œstrogènes $\rightarrow$ Pic LH $\rightarrow$ Ovulation (Ovocyte II) $\rightarrow$ Corps jaune $\rightarrow$ Progestérone.
    \item \textbf{Phase lutéale = 14 j constante.} Ovulation = J (Durée cycle - 14).
    \item \textbf{Fécondation :} Ampoule de la trompe. Capacitation (5-7h) $\rightarrow$ Réaction acrosomique $\rightarrow$ Blocage polyspermie (réaction corticale) $\rightarrow$ Syngamie $\rightarrow$ Zygote.
    \item \textbf{Fenêtre fertile = J-5 à J+1} (survie spz 3-5j dans glaires, survie ovocyte 12-24h).
    \item \textbf{Nidation :} J6-J7 post-fécondation (blastocyste). $\beta$-hCG maintient corps jaune $\rightarrow$ test grossesse positif $\sim$J14 post-fécondation (1er jour retard règles).
    \item \textbf{Adolescente :} Fertilité forte, cycles irréguliers/anovulatoires $\rightarrow$ méthode Ogino \textbf{inutilisable} comme contraception. Risque grossesse précoce et IST majoré (col immature). Double protection (contraception + préservatif) impérative.
\end{itemize}
\end{aretenir}

\coursec{Comportements à risque : identification et déterminants}

Après avoir revisité les \cle{bases physiologiques de la fertilité} (cycle ovarien, spermatogenèse, fenêtre de fécondité, rôle des hormones), nous comprenons maintenant qu'un rapport sexuel, même unique, peut engendrer une grossesse si les conditions biologiques sont réunies : présence d'un ovocyte viable, de spermatozoïdes mobiles et d'un terrain utérin favorable. Mais la physiologie ne dit pas tout : c'est le \textbf{comportement} de l'individu face à sa fertilité qui détermine l'issue sanitaire. Pourquoi, malgré la connaissance théorique, tant de jeunes s'exposent-ils encore ? C'est l'objet de cette section : identifier, nommer et comprendre les déterminants des comportements à risque.

\courssub{Définition et cadre conceptuel}

\begin{definition}[Comportement à risque pour la santé reproductive]
Tout acte, toute pratique ou toute abstention délibérée ou subie qui expose un individu — particulièrement un adolescent ou un jeune adulte — à une \cle{grossesse non désirée}, à une \cle{infection sexuellement transmissible (IST)} (dont le \cle{VIH/sida}), ou à une atteinte à son \cle{intégrité physique et psychique} (violences, traumatismes, stigmatisation).
\end{definition}

Cette définition est volontairement large : elle ne se limite pas à l'absence de préservatif. Elle intègre la \cle{contrainte}, la \cle{méconnaissance}, l'\cle{influence de substances} et les \cle{croyances erronées}. Un comportement devient « à risque » dès lors que la \textbf{probabilité d'un dommage évitable} devient significative.

\courssub{Les principaux comportements à risque identifiés}

Nous passons en revue six catégories majeures, documentées par les enquêtes démographiques et de santé (EDS-MICS) au Cameroun et la littérature scientifique internationale.

\courssub{Le rapport sexuel non protégé : la porte d'entrée principale}

Un rapport est dit \cle{non protégé} lorsqu'aucune méthode contraceptive efficace (préservatif masculin ou féminin, contraception hormonale, DIU, implant) n'est utilisée \textbf{correctement} et \textbf{systématiquement} au moment de la pénétration et jusqu'au retrait.

\begin{experience}[Observation : le liquide pré-éjaculatoire et le risque de grossesse]
\textbf{Question de départ :} Le retrait (coït interrompu) avant l'éjaculation protège-t-il de la grossesse ?
\textbf{Hypothèse :} Non, car le liquide pré-éjaculatoire (liquide de Cowper) peut contenir des spermatozoïdes résiduels dans l'urètre.
\textbf{Données expérimentales :} Analyses microscopiques de liquides pré-éjaculatoires chez des volontaires (études de Killick et al., 2011 ; Zukerman et al., 2003). Résultats : présence de spermatozoïdes mobiles chez 16 à 41\% des échantillons, en concentration variable mais suffisante pour féconder.
\textbf{Interprétation :} Le retrait n'empêche pas le dépôt de spermatozoïdes dans le vagin \emph{avant} l'éjaculation principale.
\textbf{Conclusion :} Le retrait n'est \textbf{pas} une méthode contraceptive fiable (indice de Pearl $\approx 22\%$). Il n'offre \textbf{aucune} protection contre les IST/VIH.
\end{experience}

\courssub{La précocité des rapports sexuels : immaturité biologique et psychosociale}

L'Organisation Mondiale de la Santé (OMS) définit l'adolescence de 10 à 19 ans. Au Cameroun, l'entrée dans la vie sexuelle active survient souvent avant 18 ans.

\begin{savaistu}[Âge au premier rapport au Cameroun]
Selon l'EDS-MICS 2018, l'âge médian au premier rapport sexuel est d'environ \textbf{16,6 ans pour les filles} et \textbf{15,8 ans pour les garçons} (15-24 ans). Près de 15\% des filles et 25\% des garçons ont eu un premier rapport avant 15 ans. Cette précocité s'accompagne souvent d'un manque d'information, de pouvoir de négociation et d'accès aux services de santé.
\end{savaistu}

\begin{aretenir}[Conséquences spécifiques de la précocité]
\begin{itemize}
    \item \textbf{Physiologiques :} Immaturité du col utérin (\cle{ectopie cervicale} : zone de transformation exposée sur l'exocol) $\rightarrow$ susceptibilité accrue aux IST (Chlamydia, HPV, VIH). Bassin osseux pas encore totalement ossifié $\rightarrow$ risque de dystocie, fistule obstétricale.
    \item \textbf{Psychosociales :} Abandon scolaire, dépendance économique, exposition aux violences conjugales, stigmatisation familiale.
\end{itemize}
\end{aretenir}

\courssub{La multiplicité des partenaires : réseaux de transmission}

Avoir des partenaires \textbf{concurrents} (plusieurs dans la même période) ou \textbf{successifs} (enchaînement rapide sans dépistage) multiplie exponentiellement le risque d'exposition aux IST/VIH. C'est la dynamique des \cle{réseaux sexuels} : si un maillon est infecté, l'agent pathogène se propage rapidement.

\begin{propriete}[Effet de réseau et risque d'IST]
Dans un réseau où la prévalence du VIH est de $p$, le risque pour un individu ayant $n$ partenaires non protégés successifs (indépendants) est approximativement $1 - (1-p)^n$. La \cle{concurrence} (partenaires simultanés) accélère la transmission car elle crée des « ponts » temporels entre sous-populations.
\end{propriete}

\courssub{Consommation d'alcool et de drogues : le jugement altéré}

L'alcool, le cannabis, le \textbf{Tramadol} (très répandu dans certains milieux jeunes au Cameroun), les amphétamines, etc., altèrent les fonctions préfrontales (prise de décision, inhibition, évaluation des conséquences).

\begin{attention}[Piège fréquent : « J'étais ivre, je n'ai pas pensé au préservatif »]
L'état d'ébriété ou d'intoxication \textbf{ne supprime pas le risque biologique}. Au contraire, il est \textbf{corrélé} à : non-utilisation du préservatif, rapports non consentis, multiplicité des partenaires, oubli de pilule. La prévention passe par la réduction de la consommation \emph{avant} la situation à risque.
\end{attention}

\courssub{Violences sexuelles, rapports transactionnels et absence de consentement}

Ce sont des situations où le pouvoir de négocier la protection est \textbf{retiré} à la victime.
\begin{itemize}
    \item \textbf{Viols et agressions :} absence totale de choix, traumatismes physiques (déchirures facilitant l'entrée du VIH), psychologiques.
    \item \textbf{Rapports transactionnels (\emph{Sex against grades} / \emph{Sex for money}) :} asymétrie de pouvoir (enseignant/élève, employeur/employée, « sugar daddy/mummy » / jeune). Le préservatif est souvent refusé par le dominant ; la contraception d'urgence ou le dépistage post-exposition (PEP) sont rarement accessibles dans l'urgence.
\end{itemize}

\begin{definition}[Consentement libre et éclairé]
Accord volontaire, informé, spécifique, réversible et donné par une personne capable. L'absence de « non » n'est pas un « oui ». La contrainte (physique, morale, économique, hiérarchique) annule le consentement.
\end{definition}

\courssub{Méconnaissance des signes de fertilité et fausses croyances}

Beaucoup de jeunes (et d'adultes) ignorent la \cle{fenêtre de fécondité} (5 jours avant + jour de l'ovulation) et croient en des « méthodes traditionnelles » inefficaces.

\begin{attention}[Erreur fréquente : « Le retrait (coït interrompu) est une méthode contraceptive fiable »]
\textbf{FAUX.}
\begin{itemize}
    \item \textbf{Mécanisme :} Le liquide pré-éjaculatoire (glandes de Cowper) lubrifie l'urètre \emph{avant} l'éjaculation. Il peut entraîner des spermatozoïdes résiduels d'une éjaculation antérieure (non urinée) ou produits précocement.
    \item \textbf{Efficacité :} Indice de Pearl \textbf{$\approx$ 22 grossesses pour 100 femmes/an} (usage courant), contre 2 pour le préservatif masculin, 0,1 pour l'implant.
    \item \textbf{IST/VIH :} Protection \textbf{nulle}.
\end{itemize}
\emph{Autres croyances sans fondement :} douche vaginale après rapport (ne remonte pas les spermatozoïdes déjà dans le col), position debout/sautille (la gravité n'empêche pas la migration spermique), « période sûre » mal calculée (cycles irréguliers à l'adolescence).
\end{attention}

\courssub{Analyse de situations : une démarche méthodique}

Face à un cas concret (récit de vie, article de presse, scénario de TP), utilisez cette grille structurée.

\begin{methode}[Grille d'analyse d'une situation à risque pour la santé reproductive]
\begin{enumerate}
    \item \textbf{Identifier l'acte sexuel :} Type (vaginal, anal, oral), contexte (consenti / forcé / transactionnel), moment dans le cycle (si connu).
    \item \textbf{Identifier la protection (ou son absence) :} Préservatif (masculin/féminin) utilisé correctement ? Contraception hormonale (pilule, injectable, implant) à jour ? Double protection ? Rien ?
    \item \textbf{Lister les risques encourus :}
    \begin{itemize}
        \item \cle{Grossesse non désirée} (probabilité selon phase du cycle, méthode).
        \item \cle{IST bactériennes} (Chlamydia, Gonocoque, Syphilis) — souvent asymptomatiques.
        \item \cle{IST virales} (VIH, HPV, Herpès, Hépatite B) — chroniques ou graves.
        \item \cle{Risques sociaux} : rupture scolaire, rejet familial, mariage forcé, stigmatisation, précarité économique.
        \item \cle{Risques psychiques} : anxiété, dépression, trouble de stress post-traumatique (violences).
    \end{itemize}
    \item \textbf{Proposer des alternatives responsables immédiat et différé :}
    \begin{itemize}
        \item \textbf{Immédiat (si < 72h-120h) :} Contraception d'urgence (Lévonorgestrel ou Ulipristal), Prophylaxie Post-Exposition (PEP) au VIH (si risque avéré, < 48h-72h), Dépistage IST, Prise en charge médicale/légale/psychologique (violences).
        \item \textbf{Durable :} Choix d'une contraception adaptée (conseil au Centre de Santé / CPEAJ), Vaccination HPV (filles 9-14 ans, rattrapage possible), Dépistage régulier, Négociation du préservatif (compétences de vie), Réorientation vers services sociaux/éducatifs.
    \end{itemize}
\end{enumerate}
\end{methode}

\courssub{Données épidémiologiques au Cameroun : regard sur la réalité}

Les chiffres de l'Enquête Démographique et de Santé (EDS-MICS 2018) et des rapports ONUSIDA/Ministère de la Santé Publique dressent un tableau inquiétant pour la tranche 15-19 ans.

\begin{popfigure}
\begin{tikzpicture}
\begin{axis}[
    ybar,
    bar width=0.35cm,
    width=\linewidth,
    height=8cm,
    enlarge x limits=0.25,
    ylabel={Pourcentage (\%)},
    symbolic x coords={Rapport précoce (<15 ans), Rapport non protégé (dernier rapport), Partenaires multiples (12 derniers mois), Connaissance VIH complète},
    xtick=data,
    x tick label style={rotate=30, anchor=east, font=\footnotesize},
    ymin=0, ymax=50,
    ymajorgrids=true,
    grid style=dashed,
    legend style={at={(0.5,-0.25)}, anchor=north, legend columns=-1, font=\footnotesize},
    nodes near coords,
    nodes near coords align={vertical},
    every node near coord/.append style={font=\tiny, rotate=90, anchor=west},
    popblue!80!popdark,
]
\addplot coordinates {
    (Rapport précoce (<15 ans), 14.5)
    (Rapport non protégé (dernier rapport), 38.2)
    (Partenaires multiples (12 derniers mois), 8.7)
    (Connaissance VIH complète, 22.4)
};
\legend{Filles 15-19 ans}
\addplot[poporange!80!popdark] coordinates {
    (Rapport précoce (<15 ans), 24.8)
    (Rapport non protégé (dernier rapport), 45.6)
    (Partenaires multiples (12 derniers mois), 18.3)
    (Connaissance VIH complète, 28.1)
};
\legend{Garçons 15-19 ans}
\end{axis}
\end{tikzpicture}
\legende{Indicateurs de santé reproductive chez les adolescents 15-19 ans au Cameroun (Source : EDS-MICS 2018, données arrondies). On note la précocité plus marquée chez les garçons, l'usage faible du préservatif chez les deux sexes, et une connaissance complète du VIH encore insuffisante.}
\end{popfigure}

\begin{exemplebox}[Exemple chiffré : Grossesses précoces au Cameroun]
D'après l'EDS-MICS 2018, \textbf{28\% des adolescentes âgées de 15 à 19 ans} ont déjà commencé leur vie reproductive : elles sont soit mères (24\%), soit enceintes de leur premier enfant (4\%). Ce taux grimpe à plus de 40\% dans les régions de l'Extrême-Nord, du Nord et de l'Adamaoua, corrélé au faible niveau d'instruction, à la pauvreté et au mariage d'enfants. Chaque grossesse précoce porte en elle le risque de mortalité maternelle (hémorragie, éclampsie, sepsis, avortement clandestin) et de mortalité néonatale.
\end{exemplebox}

\courssub{Arbre de décision : du choix à la conséquence}

Modélisons la bifurcation critique au moment du rapport sexuel. Cet arbre illustre que la \textbf{protection efficace} est le nœud déterminant.

\begin{popfigure}
\begin{tikzpicture}[
    node distance=1.5cm and 2.5cm,
    decision/.style={diamond, draw=popdark, thick, fill=popblueL, text width=2.8cm, text centered, font=\small, inner sep=2pt},
    chance/.style={rectangle, rounded corners, draw=popdark, thick, fill=poporangeL, text width=3.2cm, text centered, font=\small, inner sep=3pt},
    consequence/.style={rectangle, rounded corners, draw=popdark, thick, fill=popgreenL, text width=3.5cm, text centered, font=\small, inner sep=3pt},
    arrow/.style={-Stealth, thick, popdark},
    label/.style={font=\footnotesize, text width=2.5cm, text centered, midway, fill=popcream, rounded corners, inner sep=1pt}
]
\node[chance] (start) {Situation : Rapport sexuel imminent};
\node[decision, below of=start, align=center] (protect){Protection efficace utilisée ?\\ (Préservatif + Contraception fiable)};
\node[consequence, below left=1.5cm and 3cm of protect, align=center] (prot_yes){Conséquences immédiates :\\ Plaisir partagé, sérénité.\\ Conséquences différées :\\ Pas de grossesse non désirée,\\ Pas d'IST/VIH, Poursuite scolarité/projets.};
\node[consequence, below right=1.5cm and 3cm of protect, align=center] (prot_no){Conséquences immédiates :\\ Risque grossesse (selon cycle),\\ Risque IST/VIH (selon statut partenaire),\\ Anxiété post-coïtale.\\ Conséquences différées :\\ Grossesse précoce / Avortement clandestin,\\ IST chroniques / Sida,\\ Décrochage scolaire, Stigmatisation,\\ Fistule obstétricale, Mortalité maternelle.};

\draw[arrow] (start) -- (protect);
\draw[arrow] (protect) -- node[label, left] {OUI} (prot_yes);
\draw[arrow] (protect) -- node[label, right] {NON} (prot_no);
\end{tikzpicture}
\legende{Arbre de décision simplifié : l'absence de protection efficace ouvre la voie à une cascade de risques sanitaires et sociaux évitables. La « double protection » (préservatif + contraception hormonale/DIU/implant) est la stratégie optimale pour les jeunes actifs sexuellement.}
\end{popfigure}

\courssub{Exemple d'application : étude de cas résolue}

\begin{exoresolu}[Cas pratique : Amina, 16 ans, élève en Première C]
\textbf{Énoncé :} Amina, 16 ans, élève sérieuse, entretient une relation avec Paul, 19 ans, apprenti mécanicien. Paul refuse le préservatif (« ça gâte le plaisir », « si tu m'aimes, tu fais confiance »). Amina ne prend pas de pilule (peur des effets secondaires, pas d'argent pour l'implant). Lors d'une fête, après deux bières, ils ont un rapport non protégé. Amina est en milieu de cycle (J14). Elle a peur. Que faire ? Analysez avec la grille.

\textbf{Solution détaillée :}

\textbf{1. Acte :} Rapport vaginal consenti mais sous influence alcoolique, non protégé, en période fertile estimée (J14 $\approx$ ovulation). Risque grossesse \textbf{maximal} (fenêtre de fécondité).

\textbf{2. Protection :} \textbf{Aucune.} Retrait non pratiqué (Paul éjacule en elle). Pilule non prise. Pas de DIU/implant.

\textbf{3. Risques encourus :}
\begin{itemize}
    \item \textbf{Grossesse :} Probabilité $\approx$ 20-30\% pour un rapport unique en J14 sans protection. Conséquence : grossesse précoce, risque abandon scolaire, avortement clandestin (illégal, dangereux), mariage forcé.
    \item \textbf{IST/VIH :} Statut sérologique de Paul inconnu. Risque Chlamydia, Gonocoque, HPV, VIH. L'alcool a pu masquer des symptômes ou favoriser des micro-lésions.
    \item \textbf{Social/Psychique :} Anxiété majeure, perte de confiance, isolement, chantage potentiel.
\end{itemize}

\textbf{4. Alternatives responsables :}
\begin{itemize}
    \item \textbf{URGENT (J0 à J3-J5) :} Se rendre \textbf{immédiatement} au Centre de Santé le plus proche ou au \textbf{CPEAJ} (Centre de Promotion de l'Éducation et de l'Action des Jeunes).
    \begin{itemize}
        \item Demander la \textbf{contraception d'urgence} (Lévonorgestrel 1,5 mg en dose unique, efficace jusqu'à 72h, mieux jusqu'à 120h avec Ulipristal). Gratuite ou faible coût en secteur public.
        \item Demander la \textbf{Prophylaxie Post-Exposition (PEP) au VIH} (Truvada + Isentress ou équivalent) \textbf{si risque VIH estimé élevé} (partenaire à risque, rapport anal, violences). À initier \textbf{< 48h}, idéalement < 4h. Prescrite par médecin.
        \item Dépistage IST de base (syphilis, hépatite B) + traitement présomptif si indiqué (ex: Azithromycine + Cefixime).
        \item Conseil psychologique, dépistage grossesse à J21 si règles absentes.
    \end{itemize}
    \item \textbf{COURT TERME :} Choisir une contraception durable (Implant sous-cutané 3-5 ans = meilleur rapport efficacité/coût/observance pour adolescente ; ou DIU Cuivre ; ou Injectable trimestriel). Négocier le préservatif \textbf{systématique} (double protection) avec Paul. Si refus persistant $\rightarrow$ questionner la relation (respect, projet de vie).
    \item \textbf{LONG TERME :} Vaccination HPV (si non faite), éducation à la vie affective et sexuelle (compétences de négociation, estime de soi), maintien à l'école.
\end{itemize}

\textbf{Conclusion pédagogique :} Amina n'est pas « coupable », elle est \textbf{en situation de vulnérabilité} (asymétrie de pouvoir, manque d'accès contraception, pression sociale, alcool). La réponse doit être \textbf{médicale, humaine, non jugeante et rapide}. Le droit à la santé reproductive inclut l'accès à la contraception d'urgence et à la PEP sans obstacle.
\end{exoresolu}

\begin{aretenir}[Synthèse des déterminants des comportements à risque]
\begin{center}
\begin{tabularx}{\linewidth}{|l|X|X|}\hline
\textbf{Dimension} & \textbf{Exemples de déterminants} & \textbf{Leviers d'action} \\ \hline
\textbf{Individuelle} & Méconnaissance, fausses croyances, faible estime de soi, consommation substances & Éducation complète (EVC), compétences de vie, estime de soi \\ \hline
\textbf{Relationnelle} & Pression partenaire, violences, transactionnalité, absence communication & Éducation au consentement, genre, droits, médiation \\ \hline
\textbf{Communautaire} & Normes sociales (mariage précoce, silence sur sexualité), stigmatisation services & Mobilisation leaders, parents, pairs-éducateurs, médias \\ \hline
\textbf{Structurelle} & Pauvreté, accès limité services (ruptures stock, accueil non adapté), lois restrictives (avortement), éducation sexuelle absente/insuffisante & Plaidoyer, financement, formation prestataires, CPEAJ fonctionnels, accès à l'avortement sécurisé \\ \hline
\end{tabularx}
\end{center}
\end{aretenir}

Cette analyse montre que la lutte ne se résume pas à « distribuer des préservatifs ». Elle exige une \textbf{approche multisectorielle} (Santé, Éducation, Affaires sociales, Justice, Jeunesse) centrée sur les \textbf{droits} et le \textbf{pouvoir d'agir} des jeunes. La section suivante examinera en détail les conséquences dramatiques des grossesses précoces et des avortements clandestins.

\coursec{Grossesses précoces : conséquences sanitaires et sociales}

\noindent
Dans la section précédente, nous avons identifié les comportements à risque — rapports non protégés, multiplicité des partenaires, consommation d'alcool et de drogues — qui exposent les jeunes à des grossesses non désirées. Nous allons maintenant analyser en profondeur ce qui se passe lorsqu'une adolescente, dont le corps n'a pas achevé sa croissance, devient mère. C'est un drame sanitaire et social qui brise des destins, et que la science, la loi et l'éducation peuvent prévenir.

\courssub{Définition et cadre physiologique : un corps pas encore prêt}

\begin{definition}[Grossesse précoce]
On appelle \cle{grossesse précoce} toute grossesse survenant chez une fille âgée de moins de 18 ans, voire 20 ans. À cet âge, la maturation physique n'est pas achevée : le \cle{bassin} (pelvis) n'a pas atteint ses dimensions définitives, la croissance staturale peut se poursuivre, et les réserves nutritionnelles (fer, acide folique, calcium) sont souvent insuffisantes pour soutenir le développement fœtal sans puiser dans le capital de la mère.
\end{definition}

\begin{experience}[Observation clinique : le bassin immature]
\textbf{Matériel :} Radiographies de bassin (ou schémas anatomiques) d'une adolescente de 15 ans et d'une femme de 25 ans.\\
\textbf{Observation :} Chez l'adolescente, les ailes iliaques sont moins évasées, le détroit supérieur est plus étroit (forme android ou anthropoïde), la conjuguée vraie est plus courte. La symphyse pubienne et les ligaments sacro-iliaques sont moins souples.\\
\textbf{Interprétation :} Ce bassin « immature » ne permet pas toujours le passage de la tête fœtale (diamètre bipariétal moyen 9,5 cm) sans dystocie. Le travail se prolonge, augmentant les risques de souffrance fœtale, de rupture utérine et de fistule obstétricale.\\
\textbf{Conclusion :} L'immaturité pelvienne est la cause anatomique majeure de la morbidité maternelle spécifique aux adolescentes.
\end{experience}

\courssub{Risques maternels spécifiques : des complications évitables}

Les données épidémiologiques mondiales (OMS, UNFPA) et les registres des maternités camerounaises (CHU Yaoundé, Laquintinie Douala, hôpital régional de Garoua) convergent : l'adolescente n'est pas une « petite adulte » obstétricale.

\begin{propriete}[Risque maternel majoré chez l'adolescente (Données OMS / Enquête DHS Cameroun 2018)]
Le risque de \cle{mortalité maternelle} est \textbf{2 à 5 fois plus élevé} chez les filles de 15--19 ans que chez les femmes de 20--24 ans. Les principales causes de décès et de morbidité sévère sont :
\begin{itemize}
    \item \textbf{Hypertension artérielle gravidique (HTA gravidique) et \cle{éclampsie} :} L'immaturité vasculaire et endothéliale favorise une placentation défectueuse. L'éclampsie (convulsions) est une urgence vitale.
    \item \textbf{Anémie sévère :} Les besoins en fer doublent pendant la grossesse. L'adolescente, souvent déjà anémique (règles abondantes, paludisme, alimentation déséquilibrée), ne peut pas couvrir les besoins fœto-placentaires. Hb < 7 g/dL fréquent.
    \item \textbf{Fistule obstétricale :} Conséquence directe du travail prolongé sans césarienne à temps (voir encadré).
    \item \textbf{Hémorragies de la délivrance :} Atonie utérine par surdistension (macrosomie fœtale, polyhydramnios, travail prolongé) et immaturité des fibres musculaires.
    \item \textbf{Césarienne fréquente :} Indiquée pour disproportion fœto-pelvienne, présentation du siège, souffrance fœtale aigüe. Elle expose aux risques infectieux, anesthésiques et aux complications des grossesses ultérieures (rupture utérine sur cicatrice).
\end{itemize}
\end{propriete}

\begin{definition}[Fistule obstétricale]
La \cle{fistule obstétricale} est une communication anormale entre la \cle{vessie} et le \cle{vagin} (fistule vésico-vaginale) ou entre le \cle{rectum} et le \cle{vagin} (fistule recto-vaginale). Elle survient après un \cle{travail prolongé} (souvent > 24--48 h) sans césarienne : la tête fœtale comprime les tissus mères contre le bassin, provoquant une ischémie, une nécrose et une lyse tissulaire. \textbf{Conséquence :} incontinence urinaire et/ou fécale permanente, infections récidivantes, infertilité secondaire, exclusion sociale totale (répudiation, isolement). Au Cameroun, on estime à plusieurs milliers le nombre de femmes vivant avec une fistule non réparée, majoritairement issues de grossesses précoces en zone rurale.
\end{definition}

\begin{popfigure}
\begin{tikzpicture}[scale=0.85, every node/.style={font=\small}]
% Axes
\draw[->, thick] (0,0) -- (11,0) node[right] {Ratio de risque (Adolescente < 18 ans / Femme 20-30 ans)};
\draw[->, thick] (0,0) -- (0,5.5) node[above] {Fréquence relative};
% Barres
\fill[poporange!80] (1,0) rectangle (2.5,4.2) node[midway, above, black] {4,2};
\fill[popblue!80] (3.5,0) rectangle (5,3.5) node[midway, above, black] {3,5};
\fill[popgreen!80] (6,0) rectangle (7.5,2.8) node[midway, above, black] {2,8};
\fill[poppurple!80] (8.5,0) rectangle (10,2.2) node[midway, above, black] {2,2};
% Étiquettes
\node[align=center] at (1.75,-0.5) {Éclampsie};
\node[align=center] at (4.25,-0.5) {Fistule\\obstétricale};
\node[align=center] at (6.75,-0.5) {Anémie\\sévère};
\node[align=center] at (9.25,-0.5) {Césarienne\\d'urgence};
% Ligne de référence 1
\draw[dashed, thick, popdark] (0,1) -- (11,1) node[right] {Référence (Femme 20-30 ans = 1)};
% Légende simplifiée sans tabular
\node[draw, fill=popcream, rounded corners, anchor=north east, align=left, text width=4cm] at (10.5,5) {
    \textbf{Source :} OMS, UNFPA, DHS Cameroun 2018.\par
    Ratios approximatifs, ajustés sur l'âge.\par
    \textit{Interprétation :} Une adolescente a 4,2 fois plus de risque d'éclampsie.
};
\end{tikzpicture}
\legende{Infographie : Risques obstétricaux majeurs multipliés chez l'adolescente de moins de 18 ans par rapport à la femme de 20--30 ans. L'éclampsie et la fistule sont les complications les plus spécifiquement majorées par l'immaturité physiologique.}
\end{popfigure}

\courssub{Risques fœto-néonataux : le prix payé par l'enfant}

La grossesse précoce est une \cle{grossesse à haut risque} pour le fœtus et le nouveau-né. La compétition nutritionnelle mère-fœtus (le « vol » de nutriments par le fœtus chez une mère en croissance) et l'immaturité utéro-placentaire expliquent ces issues défavorables.

\begin{experience}[Analyse de données : Poids de naissance selon l'âge maternel (Données DHS Cameroun / Registres maternités)]
\textbf{Question :} L'âge maternel influence-t-il le poids de naissance et la prématurité ?\\
\textbf{Données (extrait simulé représentatif) :}
\begin{center}
\begin{tabularx}{\linewidth}{|l|X|X|X|X|}\hline
\rowcolor{popblueL}
\textbf{Âge maternel} & \textbf{\% Poids < 2500 g} & \textbf{\% Prématurité (< 37 SA)} & \textbf{Mortinatalité (‰)} & \textbf{Mortalité néonatale (‰)} \\ \hline
< 16 ans & 28\% & 22\% & 45 & 38 \\ \hline
16--19 ans & 18\% & 15\% & 32 & 28 \\ \hline
20--24 ans & 11\% & 9\% & 22 & 18 \\ \hline
25--29 ans & 9\% & 7\% & 18 & 14 \\ \hline
\end{tabularx}
\end{center}
\textbf{Interprétation :} Plus la mère est jeune, plus les indicateurs se dégradent. Le pic de risque est avant 16 ans.\\
\textbf{Conclusion :} La grossesse précoce est un facteur indépendant de faible poids de naissance, de prématurité et de mortalité périnatale.
\end{experience}

\begin{propriete}[Conséquences fœto-néonatales majeures]
\begin{itemize}
    \item \textbf{Prématurité :} Accouchement avant 37 SA. Poumons immatures (syndrome de détresse respiratoire), thermorégulation défaillante, risque hémorragie intraventriculaire, entérocolite nécrosante.
    \item \textbf{Faible poids de naissance (< 2500 g) :} Retard de croissance intra-utérin (RCIU) asymétrique (cerveau épargné, abdomen réduit) ou symétrique (infection, anomalie génétique). Conséquences à long terme : retard staturo-pondéral, troubles neurodéveloppementaux, maladies chroniques à l'âge adulte (hypothèse DOHaD : \emph{Developmental Origins of Health and Disease}).
    \item \textbf{Mortinatalité et mortalité néonatale :} Liées à la prématurité, l'asphyxie périnatale (travail prolongé), les infections néonatales (rupture prématurée des membranes, hygiène précaire).
\end{itemize}
\end{propriete}

\courssub{Conséquences sociales : le cercle vicieux de la pauvreté}

Au-delà du corps, c'est tout l'avenir de la jeune fille qui bascule. Au Cameroun, comme dans beaucoup de pays d'Afrique subsaharienne, la grossesse précoce est la première cause d'exclusion scolaire des filles.

\begin{popfigure}
\begin{tikzpicture}[node distance=1.2cm and 1.5cm, >=Stealth, font=\small]
% Nœuds
\node[draw, rounded corners, fill=poporange!20, thick, align=center] (gp){Grossesse précoce\\ (non désirée)};
\node[draw, rounded corners, fill=poppink!20, thick, right=of gp, align=center] (as){Arrêt scolaire\\ (souvent définitif)};
\node[draw, rounded corners, fill=poppurple!20, thick, right=of as, align=center] (fq){Faible qualification\\ (pas de diplôme)};
\node[draw, rounded corners, fill=popblue!20, thick, right=of fq, align=center] (ep){Emploi précaire\\ / secteur informel};
\node[draw, rounded corners, fill=popgreen!20, thick, right=of ep, align=center] (pauv){Pauvreté\\ économique};
\node[draw, rounded corners, fill=popred!20, thick, below=of pauv, align=center] (srf){Santé reproductive\\ fragile};
% Flèches principales
\draw[->, thick, poporange] (gp) -- (as);
\draw[->, thick, poppink] (as) -- (fq);
\draw[->, thick, poppurple] (fq) -- (ep);
\draw[->, thick, popblue] (ep) -- (pauv);
\draw[->, thick, popgreen] (pauv) -- (srf);
% Flèche de retour (cercle vicieux)
\draw[->, thick, popred, bend left=40] (srf) to node[left, align=center, text width=2.5cm] {\textbf{Cercle vicieux :}\\ nouvelles grossesses précoces non planifiées, absence de soins} (gp);
% Annotations
\node[draw, fill=popcream, rounded corners, anchor=north west, text width=5cm] at (gp.north west) {
    \textbf{Point de rupture :} L'école.\par Maintien scolarité = autonomie future.
};
\end{tikzpicture}
\legende{Fluxogramme : Le cercle vicieux de la grossesse précoce. L'arrêt scolaire initie une cascade aboutissant à la pauvreté et à une santé reproductive fragile, qui favorise de nouvelles grossesses précoces. L'éducation est le levier principal pour briser ce cycle.}
\end{popfigure}

\begin{aretenir}[Conséquences sociales structurantes]
\begin{enumerate}
    \item \textbf{Exclusion / Abandon scolaire :} Stigmatisation par les pairs et parfois le corps enseignant, fatigue physique, absence de garde d'enfant, pression familiale pour le mariage. Au Cameroun, moins de 10\% des filles enceintes reprennent le chemin de l'école après l'accouchement.
    \item \textbf{Stigmatisation familiale et communautaire :} Honte, « déshonneur », perte de la dot (bride price), rejet par le partenaire. La fille devient une « charge ».
    \item \textbf{Précarité économique :} Sans diplôme, elle accède seulement à l'économie informelle (vente au marché, travail domestique), revenus instables, dépendance vis-à-vis de la famille ou du mari.
    \item \textbf{Mariage forcé / précoce :} Souvent présenté comme « solution » pour légitimer l'enfant. Perpétue le cycle : épouse enfant, mère enfant, sans autonomie décisionnelle (contraception, espacement des naissances).
    \item \textbf{Perte d'autonomie et pauvreté intergénérationnelle :} Ses enfants auront moins accès à l'éducation, à la santé, à la nutrition. Le cycle se répète.
\end{enumerate}
\end{aretenir}

\courssub{Cadre légal camerounais : droits et protections}

La loi camerounaise protège l'élève enceinte, mais l'ignorance des textes et les pesanteurs socioculturelles entravent leur application.

\begin{attention}[Droit à l'éducation de l'élève enceinte -- Piège fréquent]
\textbf{Une élève enceinte a le DROIT ABSOLU de continuer sa scolarité.}
\begin{itemize}
    \item \textbf{Loi n° 98/004 du 14 avril 1998 d'orientation de l'éducation :} L'éducation est un droit pour tout enfant. L'exclusion pour cause de grossesse est une discrimination fondée sur le sexe.
    \item \textbf{Circulaire MINESEC n° 00000... (annuelle, rentrée scolaire) :} Rappelle aux chefs d'établissement l'interdiction d'exclure une élève enceinte. Elle doit être autorisée à passer les examens (Probatoire, Baccalauréat). Des aménagements (horaires, infirmerie) sont prévus.
    \item \textbf{Code pénal (Art. 337--339) :} L'avortement est illégal (sauf danger vital pour la mère, certifié par deux médecins). L'avortement clandestin est un délit pénal pour la femme et le praticien. \textbf{Attention :} La peur de la prison pousse vers les avortements clandestins mortels.
\end{itemize}
\textbf{Erreur fréquente :} Croire que l'école a le droit de « renvoyer » l'élève pour « mauvais exemple ». C'est \textbf{illégal} et constitue une \textbf{violence institutionnelle}. Tout élève, parent ou enseignant doit le signaler (Délégué régional MINESEC, Commission des droits de l'enfant).
\end{attention}

\begin{savaistu}[Le Cameroun et la fécondité des adolescentes]
Le Cameroun affiche l'un des \textbf{taux de fécondité des adolescentes (15--19 ans) les plus élevés de la zone CEMAC} : \textbf{> 100‰} (plus de 100 naissances pour 1000 filles de cet âge), selon les Enquêtes Démographiques et de Santé (EDS) successives. Les régions de l'Extrême-Nord, du Nord et de l'Adamaoua enregistrent les taux les plus forts (mariages précoces, faible scolarisation des filles). La région du Centre (Yaoundé) voit une augmentation des grossesses précoces en milieu urbain scolaire, liées à la précocité des rapports et au faible usage contraceptif.
\end{savaistu}

\courssub{Prévention et alternatives responsables : briser le cycle}

Face à ce constat, l'approche « abstinence seulement » a montré ses limites. Une stratégie globale, fondée sur les droits et les preuves scientifiques, s'impose.

\begin{methode}[Argumentaire pour convaincre un pair ou un parent : « L'école, meilleure contraception sociale »]
\textbf{Étape 1 -- Écoute et empathie :} « Je comprends tes craintes / tes valeurs. Parlons sans jugement. »\\
\textbf{Étape 2 -- Faits scientifiques :} « Le corps d'une fille de 15 ans n'est pas prêt. Le risque de mourir en accouchant est 4 fois plus grand. Le risque de fistule (incontinence à vie) est réel. »\\
\textbf{Étape 3 -- Projet de vie :} « Tu veux devenir infirmière / enseignante / entrepreneure ? La grossesse maintenant, c'est l'arrêt de l'école. Sans diplôme, c'est la précarité. »\\
\textbf{Étape 4 -- Solutions concrètes :} « Il existe des méthodes pour ne pas tomber enceinte : préservatif (protège aussi du VIH), pilule, implant, DIU. Elles sont gratuites ou peu chères au Centre de Santé / Planning Familial. »\\
\textbf{Étape 5 -- Appui :} « Je t'accompagne au Centre de Santé des Jeunes (CSJ) / à l'infirmerie scolaire. C'est confidentiel. »
\end{methode}

\begin{exemplebox}[Cas concret : Aïcha, 16 ans, classe de Première C, Maroua]
\textbf{Situation :} Aïcha, excellente élève, tombe enceinte de son petit ami (19 ans, étudiant). Ses parents veulent la marier immédiatement au père de l'enfant (coutume). Le proviseur menace de l'exclure « pour préserver l'image de l'école ».
\textbf{Analyse :}
\begin{itemize}
    \item \textbf{Risques sanitaires :} Bassin immature (région sahélienne, malnutrition chronique possible), anémie, éclampsie.
    \item \textbf{Droits violés :} Exclusion illégale (Circulaire MINESEC). Mariage forcé (Code civil : âge minimum 18 ans, consentement libre).
    \item \textbf{Alternatives :} Maintien à l'école avec aménagements (infirmerie, horaires). Accès à la contraception post-partum (implant) pour espacer. Soutien psychosocial (assistante sociale, ONG locale).
\end{itemize}
\textbf{Issues possibles :} Si Aïcha est informée de ses droits et accompagnée, elle peut obtenir son Probatoire, poursuivre des études d'infirmière, et devenir un modèle pour ses sœurs. Sans accompagnement : mariage précoce, multiparité rapide, pauvreté, fistule potentielle.
\end{exemplebox}

\begin{exoresolu}[Analyse d'une situation : Identifier les risques et proposer des alternatives]
\textbf{Énoncé :} Fatou, 15 ans, élève en Première TI à Douala, vit avec sa tante. Elle a des rapports non protégés avec son copain de 17 ans. Elle rate ses règles depuis 6 semaines. Elle a peur d'en parler à sa tante (« elle va me tuer ») et à l'infirmière de l'école (« elle va le dire au proviseur »). Elle pense à acheter un « produit » chez le vendeur ambulant du carrefour pour « faire partir la grossesse ».
\textbf{1. Identifiez au moins 3 comportements à risque dans cette situation.}\\
\textbf{2. Expliquez les conséquences sanitaires immédiates et différées pour Fatou et le fœtus si elle mène la grossesse à terme sans suivi.}\\
\textbf{3. Proposez un plan d'action concret, légal et réaliste pour Fatou.}

\tcblower
\textbf{Solution détaillée :}
\textbf{1. Comportements à risque :}
\begin{itemize}
    \item Rapports sexuels non protégés (risque grossesse + IST/VIH).
    \item Absence de dialogue avec un adulte de confiance / professionnel de santé (retard diagnostic, isolement).
    \item Intention de recourir à un avortement clandestin (produit non stérile, dosage inconnu, geste non médicalisé) $\rightarrow$ risque hémorragie, infection, perforation utérine, mort, stérilité, prison.
\end{itemize}
\textbf{2. Conséquences si grossesse menée à terme sans suivi :}
\begin{itemize}
    \item \textbf{Mère :} Anémie sévère (pas de supplémentation fer/acide folique), éclampsie non dépistée (pas de prise de tension), travail prolongé $\rightarrow$ fistule / rupture utérine / hémorragie délivrance, césarienne d'urgence dans de mauvaises conditions.
    \item \textbf{Fœtus / Nouveau-né :} RCIU, prématurité, faible poids, asphyxie néonatale, mortinatalité, mortalité néonatale (absence de réanimation).
\end{itemize}
\textbf{3. Plan d'action pour Fatou :}
\begin{enumerate}
    \item \textbf{Urgence :} Se rendre \textbf{aujourd'hui} à l'infirmerie scolaire ou au Centre de Santé des Jeunes (CSJ) le plus proche (ex: Centre de Santé de District de son arrondissement). \textbf{Confidentialité garantie} (secret médical, déontologie).
    \item \textbf{Confirmation :} Test de grossesse (urines/sang) + Échographie datation (gratuite dans certains CSJ / ONG comme CAMNAFAW, UNFPA).
    \item \textbf{Prise en charge prénatale :} Consultations prénatales (CPN) mensuelles : surveillance tension, poids, Hb, prophylaxie paludisme (MILDA + IPTp), supplémentation Fer/Acide folique, préparation à l'accouchement.
    \item \textbf{Droits scolaires :} Informer le proviseur (ou un enseignant de confiance) de sa grossesse \textbf{après} confirmation médicale. Exiger l'application de la Circulaire MINESEC : aménagement d'horaires, autorisation d'absence pour CPN, passage des examens. Refuser toute exclusion.
    \item \textbf{Soutien psychosocial :} Assistante sociale de l'école / ONG (ALVF, RENATA) pour médiation familiale (tante), prévention mariage forcé, orientation vers formation professionnelle si arrêt scolaire définitif.
    \item \textbf{Contraception post-partum :} Conseil et pose d'une méthode longue durée (Implant 3 ans ou DIU 5--10 ans) \textbf{avant} la sortie de maternité pour éviter grossesse rapprochée.
\end{enumerate}
\end{exoresolu}

\begin{exercice}[Mise en situation : Plaidoyer]
Vous êtes délégué(e) de la classe de Première C. Une camarade, Marie, 17 ans, vient d'apprendre sa grossesse. Elle est paniquée : « Je vais arrêter l'école, mes parents vont me chasser, mon copain nie la paternité. » En vous appuyant sur les connaissances de cette leçon (données sanitaires, droits légaux, ressources locales), rédigez un \textbf{message de soutien et d'orientation} (150--200 mots) que vous pourriez lui envoyer par WhatsApp ou lui dire en personne. Votre message doit être empathique, factuel, juridiquement exact et orienté vers l'action.
\end{exercice}

\begin{corrige}[Exercice : Plaidoyer pour Marie]
\textbf{Éléments attendus dans la réponse :}
\begin{itemize}
    \item \textbf{Empathie :} « Chère Marie, je suis là pour toi, pas de jugement. »
    \item \textbf{Faits sanitaires :} « Ton corps est encore en croissance, les risques (éclampsie, anémie, fistule) sont réels mais \textbf{évitables} par un suivi prénatal précoce et gratuit au Centre de Santé. »
    \item \textbf{Droits :} « Tu as le DROIT de rester à l'école (Circulaire MINESEC). Le proviseur ne peut pas t'exclure. Tu passeras ton Probatoire. »
    \item \textbf{Action immédiate :} « On y va ensemble \textbf{ce midi} à l'infirmerie / au CSJ de [Nom du quartier]. C'est confidentiel. Ils feront le test, l'écho, te donneront le fer/acide folique et la MILDA. »
    \item \textbf{Soutien social :} « Si tes parents réagissent mal, l'assistante sociale de l'école / l'ONG RENATA / ALVF peuvent médier. Le mariage forcé est illégal (18 ans minimum). »
    \item \textbf{Avenir :} « Ton diplôme, c'est ton autonomie. On t'aide à t'organiser (horaires, garde future). »
    \item \textbf{Contraception future :} « Après l'accouchement, on choisira une méthode (implant/DIU) pour que tu décides de la suite. »
\end{itemize}
\textbf{Note :} La qualité du message réside dans l'équilibre entre chaleur humaine et précision technique/juridique.
\end{corrige}

\coursec{Avortements clandestins : un drame de santé publique}

\medskip
\noindent Dans la section précédente, nous avons vu qu'une grossesse précoce, qu'elle soit désirée ou non, expose la jeune fille à des risques obstétricaux majeurs (fistules, éclampsie, accouchement difficile) et à une rupture scolaire souvent définitive. Face à cette situation, certaines jeunes filles, par peur du jugement familial, de l'exclusion scolaire ou par manque d'information sur la contraception d'urgence, prennent la décision dramatique de recourir à un avortement clandestin. Observons les conséquences de ce choix désespéré.

\courssub{Définition et cadre légal au Cameroun}

\begin{definition}[Avortement clandestin]
On appelle \textbf{avortement clandestin} toute interruption volontaire de grossesse (IVG) réalisée :
\begin{itemize}
    \item En dehors du cadre légal prévu par la loi ;
    \item Par une personne non habilitée (non personnel de santé qualifié) ;
    \item Ou dans un environnement non médicalisé (domicile, arrière-boutique, « cabinet » illégal) ne garantissant pas les conditions d'asepsie, de surveillance hémodynamique et de prise en charge des urgences.
\end{itemize}
\cle{Avortement clandestin} \cle{IVG}
\end{definition}

\noindent Au Cameroun, le cadre légal est strictement défini par le \textbf{Code Pénal (Articles 337 à 339)}.

\begin{propriete}[Cadre légal de l'IVG au Cameroun (Code Pénal, Art. 337-339)]
L'avortement est \textbf{autorisé} uniquement dans deux cas stricts, et doit être pratiqué par un médecin :
\begin{enumerate}
    \item \textbf{Péril maternel :} La grossesse met en danger la vie de la mère (ex : grossesse extra-utérine, cardiopathie sévère décompensée, éclampsie grave avant viabilité fœtale). L'attestation de deux médecins est requise (sauf urgence vitale immédiate).
    \item \textbf{Viol :} La grossesse est issue d'un viol ou d'un attentat à la pudeur. Une attestation du Procureur de la République est obligatoire avant l'acte.
\end{enumerate}
Tout avortement pratiqué en dehors de ces deux cas est un \textbf{délit pénal} passible d'emprisonnement (1 à 5 ans) et d'amende pour la femme, le praticien et les complices.
\cle{Cadre légal IVG Cameroun}
\end{propriete}

\begin{attention}[Piège fréquent : Confusion « IVG » et « Avortement spontané »]
Ne confondez pas :
\begin{itemize}
    \item \textbf{Avortement spontané (fausse couche) :} Arrêt naturel de la grossesse avant 22 SA (semaines d'aménorrhée) ou fœtus < 500 g. Ce n'est \textbf{pas} un délit, c'est un accident obstétrical fréquent (10-15\% des grossesses cliniques).
    \item \textbf{IVG (Interruption Volontaire de Grossesse) :} Acte médical intentionnel. Elle est \textbf{légale} si elle respecte les conditions de l'Art. 337, \textbf{clandestine} sinon.
\end{itemize}
Au Probatoire, précisez toujours le cadre légal camerounais, ne recopiez pas la loi française (loi Veil 1975) qui est plus permissive (IVG sur demande jusqu'à 14 SA).
\end{attention}

\courssub{Techniques dangereuses : l'horreur du « fait maison »}

\noindent Faute d'accès à l'IVG légale, les jeunes filles recourent à des méthodes traditionnelles ou détournées, transmises par le bouche-à-oreille ou les réseaux sociaux. Analysons ces pratiques à la lumière de l'anatomie et de la physiologie utérine.

\begin{experience}[Analyse des risques des méthodes abortives traditionnelles (Démarche d'investigation)]
\textbf{Observation :} Dans les services de gynécologie-obstétrique des hôpitaux de référence (Laquintinie Douala, Général Yaoundé, Hôpitaux régionaux), on reçoit régulièrement des jeunes filles en urgence vitale après « tentative d'avortement ».

\textbf{Hypothèses :} Les substances ou objets introduits dans l'utérus provoquent des lésions directes (perforation) ou des réactions toxiques/infectieuses systémiques.

\textbf{Données (Techniques recensées au Cameroun) :}
\begin{center}
\begin{tabularx}{\linewidth}{|l|X|X|}\hline
\rowcolor{popblueL}
\textbf{Catégorie} & \textbf{Exemples concrets (Contexte camerounais)} & \textbf{Mécanisme d'action supposé / Danger principal} \\ \hline
\textbf{Corps étrangers} & Cintre métallique déplié, tige de stylo, sonde urinaire, tige de manioc, osselets. & Traumatisme direct : perforation utérine, lacération col, hémorragie par section vaisseaux utérins. \\ \hline
\textbf{Plantes abortives toxiques} & Décoctions de feuilles de manioc (cyanure), « ndolé sauvage » (\emph{Vernonia} spp. non identifiées), écorces de \emph{Alstonia boonei}, graines de \emph{Ricinus communis} (ricin). & Toxicité systémique (cyanure $\to$ anoxie cellulaire, ricine $\to$ nécrose viscérale) + stimulation utérine anarchique $\to$ rupture utérine. \\ \hline
\textbf{Médicaments détournés} & \textbf{Misoprostol (Cytotec\textsuperscript{\textregistered})} seul, par voie orale ou vaginale, \textbf{sans suivi médical}, doses aléatoires (souvent surdoses). & Hyperstimulation utérine $\to$ tétanie utérine, rupture utérine (surtout si utérus cicatriciel), hémorragie post-partum par atonie secondaire, avortement incomplet. \\ \hline
\textbf{Curetage « sauvage »} & Pratiqué par « infirmier de quartier », matrone non formée, dans l'arrière-boutique, sans anesthésie, sans asepsie, matériel non stérile (curettes rouillées, spéculums sales). & Perforation utérine, lésions intestinales/vésicales, infection massive (septicémie), syndrome d'Asherman (synéchies). \\ \hline
\end{tabularx}
\end{center}

\textbf{Interprétation :} Toutes ces techniques partagent un point commun : elles agressent l'utérus (organe musculaire vascularisé) sans contrôle de l'hémostase, sans antibioprophylaxie, sans surveillance des signes vitaux. Elles transforment un acte médical codifié en un traumatisme grave.

\textbf{Conclusion :} L'avortement clandestin n'est pas un « acte médical raté », c'est un \textbf{acte de violence physique} infligé à l'appareil reproducteur, aux conséquences imprévisibles et souvent irréversibles.
\end{experience}

\begin{popfigure}
\begin{tikzpicture}[scale=0.9, every node/.style={font=\small}]
    % Uterus Sain (Gauche)
    \begin{scope}[xshift=-4cm]
        \draw[popdark, line width=1.5pt, fill=popgreenL!30] (0,0) .. controls (0,1.5) and (1.5,2.5) .. (1.5,3.5) .. controls (1.5,4.5) and (0,5.5) .. (0,6) .. controls (0,5.5) and (-1.5,4.5) .. (-1.5,3.5) .. controls (-1.5,2.5) and (0,1.5) .. (0,0);
        % Endometre
        \draw[poppink, line width=1pt, fill=poppink!20] (0,0.5) .. controls (0,1.5) and (1.2,2.5) .. (1.2,3.5) .. controls (1.2,4.5) and (0,5.5) .. (0,5.5) .. controls (0,5.5) and (-1.2,4.5) .. (-1.2,3.5) .. controls (-1.2,2.5) and (0,1.5) .. (0,0.5);
        % Trompes
        \draw[popblue, line width=1.5pt] (-1.5,5.5) .. controls (-2.5,5.5) and (-2.5,6.5) .. (-1.5,6.5);
        \draw[popblue, line width=1.5pt] (1.5,5.5) .. controls (2.5,5.5) and (2.5,6.5) .. (1.5,6.5);
        % Ovaires
        \draw[poporange, fill=poporangeL!30] (-2.5,6.5) ellipse (0.4 and 0.3);
        \draw[poporange, fill=poporangeL!30] (2.5,6.5) ellipse (0.4 and 0.3);
        % Vagin
        \draw[popdark, line width=1.5pt, fill=popgreenL!30] (-0.8,0) -- (-0.8,-1.5) -- (0.8,-1.5) -- (0.8,0) -- cycle;
        % Col
        \draw[popdark, line width=1.5pt] (-0.8,0) -- (-0.8,0.5);
        \draw[popdark, line width=1.5pt] (0.8,0) -- (0.8,0.5);
        
        \node[align=center, font=\bfseries\small, popdark] at (0, -2.2) {Utérus sain (repos)};
        \node[align=center, font=\tiny, popgreen] at (0, 3) {Endomètre\\intact};
        \node[align=center, font=\tiny, popblue] at (-2.5, 7.2) {Trompe\\perméable};
        \node[align=center, font=\tiny, poporange] at (-2.5, 6.5) {Ovaire};
    \end{scope}

    % Uterus Perfore / Infecte (Droite)
    \begin{scope}[xshift=4cm]
        % Utérus déformé
        \draw[popdark, line width=1.5pt, fill=poporangeL!20] (0,0) .. controls (0,1.5) and (1.5,2.5) .. (1.5,3.5) .. controls (1.5,4.5) and (0,5.5) .. (0,6) .. controls (0,5.5) and (-1.5,4.5) .. (-1.5,3.5) .. controls (-1.5,2.5) and (0,1.5) .. (0,0);
        % Perforation fond utérin
        \draw[popink, line width=2pt, decorate, decoration={snake, amplitude=1pt, segment length=4pt}] (-0.2, 5.2) -- (-1.5, 5.8);
        \fill[popink] (-0.2, 5.2) circle (2pt);
        \node[popink, font=\tiny, align=center] at (-1.8, 6.0) {Perforation\\par corps étranger};
        
        % Infection propagation
        \draw[popink, ->, line width=1.5pt, opacity=0.8] (-0.2, 5.2) .. controls (-1, 5.5) and (-1, 4.5) .. (-1.5, 4.0);
        \draw[popink, ->, line width=1.5pt, opacity=0.8] (-0.2, 5.2) .. controls (0.5, 5.0) and (1.0, 4.0) .. (1.5, 3.5);
        \draw[popink, ->, line width=1.5pt, opacity=0.8] (-0.2, 5.2) .. controls (0, 4.5) and (0, 3.5) .. (0, 2.5);
        
        % Bactéries symboliques
        \foreach \x/\y in {-1.2/4.5, 0.8/3.8, -0.5/3.0, 1.0/2.5} {
            \draw[popink, fill=popink!20] (\x,\y) circle (1.5pt);
            \draw[popink] (\x-0.5pt,\y) -- (\x+0.5pt,\y);
            \draw[popink] (\x,\y-0.5pt) -- (\x,\y+0.5pt);
        }
        
        % Péritoine / Septicémie
        \draw[popink, dashed, line width=1pt] (-2.5, 3) rectangle (2.5, 6.5);
        \node[popink, font=\tiny, align=center] at (0, 6.8) {Péritoine $\to$ Péritonite $\to$ Septicémie};
        
        % Endomètre détruit
        \draw[popink, line width=1pt, pattern=north east lines, pattern color=popink] (0,0.5) .. controls (0,1.5) and (1.2,2.5) .. (1.2,3.5) .. controls (1.2,4.5) and (0,5.5) .. (0,5.5) .. controls (0,5.5) and (-1.2,4.5) .. (-1.2,3.5) .. controls (-1.2,2.5) and (0,1.5) .. (0,0.5);
        \node[popink, font=\tiny, align=center] at (0, 3) {Endomètre\\nécrosé / synéchies};
        
        % Trompes bouchées
        \draw[popdark, line width=1.5pt, draw=popink] (-1.5,5.5) .. controls (-2.5,5.5) and (-2.5,6.5) .. (-1.5,6.5);
        \draw[popink, line width=3pt] (-2.0, 6.0) -- (-2.0, 6.0); % Bouchon
        \node[popink, font=\tiny] at (-2.8, 6.5) {Trompe bouchée};
        
        % Vagin
        \draw[popdark, line width=1.5pt, fill=poporangeL!20] (-0.8,0) -- (-0.8,-1.5) -- (0.8,-1.5) -- (0.8,0) -- cycle;
        % Lacérations col
        \draw[popink, line width=2pt] (-0.8,0.2) -- (-0.5,0.5);
        \draw[popink, line width=2pt] (0.8,0.2) -- (0.5,0.5);
        \node[popink, font=\tiny, align=center] at (0, -0.5){Lacérations\\col / vagin};
        
        \node[align=center, font=\bfseries\small, popdark] at (0, -2.2) {Utérus après avortement clandestin};
    \end{scope}
    
    % Flèche centrale comparatif
    \draw[popblue, very thick, ->] (-1.5, -2.5) -- (1.5, -2.5) node[midway, above, font=\small, popblue] {Conséquences immédiates \& séquelles};
\end{tikzpicture}
\legende{Comparaison schématique : Utérus sain au repos (gauche) vs Utérus ravagé par un avortement clandestin (droite). Perforation du fond utérin par un corps étranger (cintre), propagation bactérienne vers le péritoine (péritonite, septicémie), nécrose endométriale source de synéchies (stérilité), lacérations cervicales et vaginales.}
\end{popfigure}

\courssub{Complications immédiates : l'urgence vitale}

\noindent Les complications surviennent dans les heures ou jours qui suivent l'acte. Elles sont la première cause de consultation en urgence gynécologique chez l'adolescente.

\begin{propriete}[Complications immédiates majeures (Triade mortelle)]
\begin{enumerate}
    \item \textbf{Hémorragie massive (Choc hémorragique) :} L'utérus est un organe très vascularisé (artères utérines = branches de l'iliaque interne). Une perforation ou une atonie utérine (épuisement musculaire après hyperstimulation par misoprostol) empêche la rétraction vasculaire physiologique. Perte de \SI{> 500}{mL} en quelques minutes $\to$ tachycardie, hypotension, sueurs froides, confusion $\to$ \textbf{choc hypovolémique}. \textbf{Urgence absolue :} Transfusion sanguine (souvent introuvable en zone rurale), hémostase chirurgicale (curetage de révision, ligature artères utérines, \textbf{hystérectomie} salvatrice).
    \item \textbf{Infection sévère (Septicémie / Péritonite) :} Rupture de la barrière stérile utérine + germes vaginaux (anaérobies, \emph{E. coli}, \emph{Streptocoques}) + matériel sale $\to$ endométrite $\to$ péritonite diffuse $\to$ septicémie (choc septique). Fièvre $> \SI{39}{\celsius}$, douleurs abdominales en « planche de bois », iléus paralytique, troubles de la coagulation (CIVD). \textbf{Mortalité élevée sans réanimation.}
    \item \textbf{Perforation utérine et lésions viscérales voisines :} Le corps étranger (cintre, curette) traverse la paroi utérine (souvent au fond, zone fine) et blesse : \textbf{vésicale} (fistule vésico-utérine/vaginale $\to$ incontinence urinaire), \textbf{rectale} (fistule recto-vaginale $\to$ incontinence fécale), \textbf{iléale/colonique} (péritonite fécale). Nécessite laparotomie d'urgence, colostomie de dérivation.
\end{enumerate}
\cle{Choc hypovolémique} \cle{Septicémie} \cle{Péritonite} \cle{Fistule}
\end{propriete}

\begin{exemplebox}[Cas clinique : « Merveille », 16 ans, Douala]
Merveille, élève en 1ère C, découvre une grossesse de 14 SA. Effrayée par la réaction de ses parents (père strict, mère absente), elle achète 4 comprimés de Cytotec\textsuperscript{\textregistered} (Misoprostol 200 µg) auprès d'un vendeur ambulant au marché. Elle les insère dans le vagin le soir. 3h plus tard : douleurs atroces, saignements abondants (imbibe 3 serviettes/heure), fièvre \SI{39,5}{\celsius}. Transportée à l'hôpital par une amie. Diagnostic : \textbf{Avortement incomplet compliqué d'hémorragie et début de septicémie}. Prise en charge : Réanimation (remplissage volémique, antibiotiques IV large spectre : Amoxicilline/Acide clavulanique + Gentamicine + Métronidazole), curetage d'évacuation sous anesthésie générale. Sortie à J5 avec prescription contraceptive (Implant). \textbf{Bilan :} Vie sauvée, mais risque de séquelles tubaires (stérilité future) et traumatisme psychologique majeur.
\end{exemplebox}

\courssub{Séquelles à long terme : l'hypothèque sur l'avenir}

\noindent Si la jeune fille survit à la phase aiguë, elle n'est pas « guérie ». Les séquelles s'installent insidieusement.

\begin{aretenir}[Séquelles majeures de l'avortement clandestin]
\begin{itemize}
    \item \textbf{Stérilité secondaire (Facteur tubaire \& utérin) :}
    \begin{itemize}
        \item \emph{Tubaire :} Salpingites séquellaires $\to$ trompes bouchées (hydrosalpinx) ou adhérences péritonéales $\to$ impossibilité rencontre spermatozoïde/ovocyte ou grossesse extra-utérine (GEU) récidivante.
        \item \emph{Utérin :} \textbf{Syndrome d'Asherman} (synéchies intra-utérines) : cicatrices fibreuses collant les parois utérines suite à curetage agressif sur endomètre basal $\to$ aménorrhée secondaire, fausses couches à répétition, placenta praevia/accreta si grossesse ultérieure.
    \end{itemize}
    \item \textbf{Douleurs pelviennes chroniques (DPC) :} Adhérences péritonéales, névralgies pudendales post-traumatiques, syndrome de congestion pelvienne. Impact sur la vie sexuelle (dyspareunie) et la qualité de vie.
    \item \textbf{Troubles psychologiques :} Syndrome de stress post-traumatique (SSPT) : cauchemars, reviviscences, évitement (gynécologue, sexualité), dépression, risque suicidaire. Stigmatisation sociale (rumeurs, rejet).
    \item \textbf{Décès maternel :} L'OMS estime que les avortements non sûrs représentent \textbf{15 à 20\% de la mortalité maternelle} au Cameroun (ratio ~400-500 décès / 100 000 NV). La plupart sont des adolescentes.
\end{itemize}
\cle{Syndrome d'Asherman} \cle{Stérilité secondaire} \cle{Mortalité maternelle}
\end{aretenir}

\begin{popfigure}
\begin{tikzpicture}
\begin{axis}[
    width=\linewidth,
    height=8cm,
    ybar, % Corrigé : ybar au lieu de ybar stacked pour pourcentages indépendants
    bar width=25pt,
    ymin=0, ymax=100,
    ylabel={Pourcentage (\% des décès maternels)},
    xlabel={Causes de mortalité maternelle (Estimation OMS Cameroun)},
    xtick={1,2,3,4,5},
    xticklabels={Hémorragie du post-partum, Infections (puerpérales/septicémie), Hypertension gravidique (Éclampsie), \textbf{Avortements non sûrs}, Autres causes (anémie, rupture utérine)},
    legend style={at={(0.5,-0.25)}, anchor=north, legend columns=-1, font=\small},
    ymajorgrids=true,
    grid style=dashed,
    nodes near coords,
    nodes near coords align={vertical},
    every node near coord/.append style={font=\tiny, rotate=90, anchor=west, yshift=2pt},
    enlarge x limits=0.2,
]
\addplot[fill=popink] coordinates {(1,35) (2,15) (3,12) (4,18) (5,20)};
\addlegendentry{Part estimée}
\end{axis}
\end{tikzpicture}
\legende{Part estimée des avortements non sûrs (clandestins) dans la mortalité maternelle au Cameroun (~18\%), 3ème cause après l'hémorragie du post-partum et les infections. Données inspirées des estimations OMS/UNFPA pour l'Afrique subsaharienne francophone.}
\end{popfigure}

\courssub{Prise en charge des complications : les Soins Post-Avortement (SPA)}

\noindent C'est un point éthique et juridique crucial : \textbf{la vie prime sur la loi}.

\begin{definition}[Soins Post-Avortement (SPA) - Standards OMS / Ministère de la Santé Publique Cameroun]
Ensemble des interventions médicales d'urgence et de suivi pour traiter les complications d'un avortement (spontané ou provoqué), comprenant :
\begin{enumerate}
    \item \textbf{Stabilisation d'urgence :} Réanimation hémodynamique (remplissage, transfusion), antibiothérapie probabiliste IV large spectre, analgésie.
    \item \textbf{Evacuation utérine complète :} Aspiration manuelle intra-utérine (AMIU) ou curetage instrumental sous contrôle visuel (hystéroscopie si dispo) ou guidage échographique. \textbf{Objectif :} Retirer tout tissu résiduel (source hémorragie/infection).
    \item \textbf{Prise en charge des lésions associées :} Réparation cervicale/vaginale, laparotomie si perforation/viscères lésés.
    \item \textbf{Conseil et contraception post-événement :} \textbf{Indispensable.} La fertilité revient dès l'ovulation suivante (J14-J21 post-avortement). Proposition de toutes les méthodes (Implant, DIU, Pilule, Préservatif) \textbf{avant la sortie}. Dépistage IST/VIH.
    \item \textbf{Soutien psychologique et social :} Écoute, orientation vers structures d'appui (ONG, services sociaux).
\end{enumerate}
\cle{SPA (Soins Post-Avortement)} \cle{AMIU}
\end{definition}

\begin{propriete}[Droit aux soins sans condition (Éthique médicale \& Droit camerounais)]
\begin{itemize}
    \item \textbf{Secret professionnel (Code de déontologie médicale, Serment d'Hippocrate) :} Le soignant \textbf{ne doit pas dénoncer} la patiente pour avortement clandestin. La priorité est de sauver la vie.
    \item \textbf{Non-discrimination :} Accès aux soins d'urgence garanti quel que soit l'âge, le statut marital, l'origine de la grossesse.
    \item \textbf{Consentement :} Mineure $\to$ consentement des parents/tuteurs recherché pour l'anesthésie/chirurgie, mais \textbf{l'urgence vitale prime} (art. 37 Code Pénal : « n'est pas punissable celui qui accomplit un acte commandé par la nécessité... pour sauver une vie »).
\end{itemize}
\end{propriete}

\begin{methode}[Conduite à tenir face à une complication d'avortement (Témoin / Proche / Personnel non médical)]
\textbf{Situation :} Jeune fille saignant abondamment, fièvre, douleurs abdominales intenses, choc (pâle, sueur, pouls filant), après tentative d'avortement.
\begin{enumerate}
    \item \textbf{APPELER LES SECOURS / AMENER AU CENTRE DE SANTÉ LE PLUS PROCHE} (Centre de Santé Intégré, Hôpital de District, Hôpital Régional). Ne pas attendre « que ça passe ». \textbf{Le temps = la vie.}
    \item \textbf{NE RIEN INTRODUIRE DANS LE VAGIN} (pas d'eau, pas de tisane, pas de gant, pas de tissu). Risque d'aggraver l'infection, la perforation, de pousser des débris vers le haut.
    \item \textbf{NE PAS DONNER À MANGER NI À BOIRE} (estomac vide obligatoire pour une éventuelle anesthésie générale en urgence).
    \item \textbf{ALLONGER LA PATIENTE}, jambes surélevées (position anti-choc), la couvrir (lutte hypothermie).
    \item \textbf{RASSURER :} « Tu es entre de bonnes mains, le secret médical est total, on va te soigner. » Lutter contre la peur de la police/parents.
    \item \textbf{SI SAIGNEMENT EXTERNE MASSIF :} Compression vulvaire propre (serviette propre) en attendant le transport.
\end{enumerate}
\end{methode}

\begin{attention}[Mythe dangereux : « Le Cytotec (Misoprostol) seul à la maison, c'est sans danger »]
\textbf{FAUX. ARCHI-FAUX.}
\begin{itemize}
    \item \textbf{Risque 1 : Hémorragie incontrôlée.} Sans échographie de datation préalable, on ignore l'âge gestationnel. $> 12$ SA $\to$ risque hémorragie massive multiplié par 10. Sans oxytociques injectables (ocytocine, ergométrique) ni possibilité de curetage immédiat $\to$ décès à domicile.
    \item \textbf{Risque 2 : Avortement incomplet.} Restes ovulaires (trophoblaste) $\to$ hémorragies répétées, infection chronique, stérilité (synéchies).
    \item \textbf{Risque 3 : Grossesse extra-utérine (GEU) méconnue.} Le Misoprostol \textbf{ne vide pas une GEU}. La patiente croit avoir avorté, la trompe éclate 2 semaines plus tard $\to$ hémorragie intra-péritonéale mortelle.
    \item \textbf{Risque 4 : Absence de contraception post-avortement.} Ovulation $\approx$ J14 $\to$ nouvelle grossesse immédiate (cercle vicieux).
    \item \textbf{Seule voie sûre :} Protocole OMS (Mifépristone 200 mg PO + Misoprostol 800 µg VO/SL/Vaginal 24-48h après) \textbf{DANS UN CADRE MÉDICALISÉ} avec échographie de contrôle J7-J14.
\end{itemize}
\end{attention}

\courssub{Prévention : briser le cycle par l'information et l'accès}

\noindent La lutte contre l'avortement clandestin ne se fait pas par la répression, mais par la \textbf{prévention primaire} (éviter la grossesse non désirée) et \textbf{secondaire} (accès à la contraception d'urgence et aux SPA).

\begin{savaistu}[L'avortement non sûr dans le monde : Chiffres clés OMS 2023]
\begin{itemize}
    \item \textbf{45\%} des avortements dans le monde sont « non sûrs » (non médicalisés / personnel non formé).
    \item \textbf{97\%} de ces avortements non sûrs ont lieu dans les \textbf{pays en développement}.
    \item \textbf{39 000 femmes} meurent chaque année des suites d'un avortement non sûr (presque toutes évitables).
    \item \textbf{7 millions} de femmes sont hospitalisées pour complications.
    \item Coût économique énorme pour les systèmes de santé (soins d'urgence, séquelles à vie).
\end{itemize}
Au Cameroun, la \textbf{Stratégie Nationale de Santé de la Reproduction (2020-2024)} intègre le renforcement de l'offre de SPA dans tous les CSI et Hôpitaux de District, et la formation des prestataires à l'AMIU (Aspiration Manuelle Intra-Utérine), technique simple, peu coûteuse, salvatrice.
\end{savaistu}

\begin{exoresolu}[Analyse d'une situation : « Le dilemme de Aïcha »]
\textbf{Énoncé :} Aïcha, 17 ans, élève en 1ère TI à Maroua, est enceinte de 10 SA (confirmée par test urinaire + écho). Son copain a fui. Ses parents menacent de la chasser si ils apprennent la grossesse. Elle entend parler d'une « tante » qui fait « le nettoyage » avec des feuilles et un cintre pour 5000 FCFA. Elle hésite entre cette option et aller au Centre de Santé Urbain (CSU) le plus proche.
\textbf{Questions :}
\begin{enumerate}
    \item Identifiez le comportement à risque envisagé par Aïcha.
    \item En vous appuyant sur vos connaissances anatomiques et physiologiques, expliquez \textbf{trois} complications immédiates probables de cette technique.
    \item Quelle est l'alternative responsable que vous lui conseillez ? Justifiez par deux arguments médicaux et un argument légal/éthique.
\end{enumerate}

\tcblower
\textbf{Solution détaillée :}
\begin{enumerate}
    \item \textbf{Comportement à risque :} Recours à un \textbf{avortement clandestin} par méthode traditionnelle traumatisante (corps étranger + plantes toxiques) pratiqué par une personne non qualifiée dans un environnement non médicalisé.
    \item \textbf{Complications immédiates probables :}
    \begin{enumerate}
        \item \textbf{Perforation utérine + lésion vasculaire :} Le cintre métallique, rigide et non stérile, traverse la paroi utérine fine (10 SA = utérus taille orange, paroi ~1cm). Il sectionne les artères utérines $\to$ \textbf{hémorragie intra-péritonéale massive, choc hypovolémique}.
        \item \textbf{Péritonite septique / Septicémie :} Introduction de germes vaginaux + germes telluriques (terre sur cintre/feuilles) + débris nécrosés dans la cavité péritonéale stérile $\to$ inflammation aiguë du péritoine $\to$ passage bactéries dans le sang $\to$ \textbf{choc septique}.
        \item \textbf{Lésions vésicales/rectales :} L'utérus antéversé repose sur la vessie ; le fond utérin est proche du rectum. Perforation postérieure $\to$ fistule recto-vaginale (incontinence fécale) ; antérieure $\to$ fistule vésico-vaginale (incontinence urinaire). Séquelles définitives sans chirurgie réparatrice complexe.
    \end{enumerate}
    \item \textbf{Alternative responsable :} \textbf{Se rendre immédiatement au CSU / Hôpital de District pour une prise en charge médicale.}
    \begin{itemize}
        \item \textit{Argument médical 1 :} À 10 SA, une \textbf{IVG médicamenteuse légale} (si cas prévu par la loi : péril maternel psychologique/physique avéré par médecin, ou viol) ou une \textbf{prise en charge de la grossesse} avec suivi prénatal adapté (adolescente = grossesse à risque) est possible. Surtout, si elle refuse la grossesse, le personnel de santé peut orienter vers un service autorisé (commission médicale) ou assurer un \textbf{SPA de qualité} si avortement déjà débuté.
        \item \textit{Argument médical 2 :} Au CSU, elle bénéficiera d'une \textbf{échographie de datation précise}, d'un \textbf{dépistage IST/VIH}, et surtout d'une \textbf{contraception efficace (Implant/DIU)} posée immédiatement après l'évacuation (si IVG légale) ou après l'accouchement, prévenant une grossesse immédiate.
        \item \textit{Argument légal/éthique :} Le \textbf{secret professionnel} lie le soignant. Elle sera soignée \textbf{sans dénonciation}. La loi (Art 37 CP) protège l'acte médical d'urgence pour sauver une vie. La « tante » n'offre aucune garantie de confidentialité, de compétence, ni de prise en charge des complications.
    \end{itemize}
\end{enumerate}
\end{exoresolu}

\begin{exercice}[Entraînement Probatoire : QCM \& Rédaction]
\textbf{1. QCM :} Une jeune fille de 16 ans consulte pour fièvre \SI{39,5}{\celsius}, douleurs abdominales diffuses en « planche de bois », saignements vaginaux modérés, 48h après avoir bu une décoction de feuilles de manioc pour « faire partir une grossesse ». L'examen retrouve une défense abdominale généralisée. Le diagnostic le plus probable est :
\begin{itemize}
    \item A. Appendicite aigüe.
    \item B. Salpingite aigüe.
    \item C. \textbf{Péritonite diffuse sur avortement septique clandestin.}
    \item D. Grossesse extra-utérine rompue.
\end{itemize}

\textbf{2. Rédaction (5 pts) :} En vous appuyant sur l'anatomie de l'appareil reproducteur féminin et la physiologie de la coagulation, expliquez pourquoi la perforation utérine par un corps étranger (ex: cintre) expose à un risque hémorragique majeur difficile à contrôler sans chirurgie. Citez les vaisseaux concernés et le mécanisme physiologique défaillant.

\textbf{3. Argumentation (5 pts) :} « La légalisation totale de l'avortement sur demande ferait chuter la mortalité maternelle. » Discutez cette affirmation en distinguant les arguments médicaux (sécurité de l'acte), sociaux (accès, stigmatisation) et éthiques (droit à la vie, autonomie), en vous référant au contexte camerounais.
\end{exercice}

\begin{corrige}[Exercice 1 - Éléments de correction]
\textbf{1. QCM :} \textbf{C.} Le tableau associe : notion d'avortement clandestin (décoction toxique + corps étranger potentiel), signes de péritonite (défense, fièvre, douleurs), saignements. C'est la définition clinique de l'avortement septique compliqué de péritonite.
\textbf{2. Rédaction :} L'utérus est vascularisé par les \textbf{artères utérines} (branches de l'artère iliaque interne) qui s'anastomosent avec les artères ovariennes. Ces artères s'enroulent en hélice dans le myomètre. \textbf{Mécanisme hémostatique physiologique :} Après expulsion fœtale/placentaire, la rétraction des fibres musculaires lisses (myomètre) comprime les sinus veineux et les artères en spire $\to$ hémostase mécanique (ligatures vivantes de Pinard). \textbf{Perforation par corps étranger :} Le cintre déchire le myomètre de manière anarchique (plaie irrégulière, étoilée), sectionnant les artères radiales et utérines \textbf{sans permettre la rétraction musculaire} (perte de continuité du muscle). Le saignement est artériel (jet, rouge vif), intra-péritonéal (invisible au premier abord) et ne s'arrête pas spontanément. Seule une ligature chirurgicale (ou embolisation) l'arrête.
\textbf{3. Argumentation :} \textit{Médical :} Légalisation $\to$ accès IVG sûre (Mifépristone/Misoprostol ou AMIU) dans structures habilitées $\to$ disparition avortements « cintre/plantes » $\to$ chute drastique complications (hémorragie, infection, perforation) $\to$ baisse mortalité maternelle (ex: Afrique du Sud post-1996 : -90\% mortalité liée avortement). \textit{Social :} Réduction stigmatisation $\to$ consultation précoce $\to$ meilleure contraception post-IVG. \textit{Éthique/Juridique Cameroun :} Débat vif. Droit à la vie (fœtus) vs Autonomie corporelle / Droit à la santé (mère). Loi actuelle restrictive (Art 337) = compromis. Argument « sur demande » : réduit avortements tardifs (plus risqués). Argument « restriction » : protection vie fœtale, risque banalisation. \textbf{Synthèse :} La preuve épidémiologique mondiale (OMS) montre que la restriction légale n'empêche pas l'avortement, elle le rend juste dangereux. La mortalité maternelle par avortement est quasi nulle là où il est légal, sûr et accessible.
\end{corrige}

\coursec{Moyens de contraception : panorama et choix responsable}

Après avoir mesuré la gravité des avortements clandestins — complications infectieuses, hémorragies, stérilité définitive, décès maternels — une évidence s'impose : la meilleure lutte contre l'avortement à risque, c'est la \textbf{prévention de la grossesse non désirée}. Au Cameroun, comme ailleurs, l'accès à une contraception adaptée, efficace et acceptée est un droit fondamental en santé reproductive. Pourtant, entre les rumeurs du quartier (« la pilule rend stérile », « l'implant fait grossir »), la peur des effets secondaires et la méconnaissance des mécanismes, beaucoup de jeunes (et d'adultes) n'utilisent aucune méthode ou une méthode inadaptée. Comment s'y retrouver parmi la dizaine de moyens disponibles ? Comment choisir \emph{sa} méthode, celle qui protège vraiment, sans nuire à sa santé ni à son avenir ? C'est l'objet de cette section : dresser un panorama rigoureux, comparer l'efficacité réelle, et donner les clés d'un \textbf{choix libre et éclairé}.

\courssub{Classification des méthodes et mécanismes d'action}

Observons une situation concrète : \emph{Aïcha, 17 ans, lycéenne à Douala, a un petit ami. Elle ne veut pas tomber enceinte maintenant, mais elle a peur des hormones. Son amie lui parle du « collier » (méthode des jours fixes), son grand frère lui glisse des préservatifs, l'infirmière du CSB lui propose l'implant. Quelle méthode choisir ?} Pour répondre, classons les moyens selon leur \textbf{principe d'action physiologique} : empêcher la rencontre spermatozoïde–ovocyte, bloquer l'ovulation, ou empêcher la nidation.

\begin{definition}[Moyen de contraception]
Un moyen de contraception est un dispositif, un médicament, une méthode chirurgicale ou comportementale qui vise à empêcher la survenue d'une grossesse lors de rapports sexuels potentiellement fécondants. On distingue les méthodes \textbf{réversibles} (arrêt = retour de fertilité) des méthodes \textbf{définitives} (stérilisation chirurgicale).
\end{definition}

\begin{popfigure}
\centering
\begin{tikzpicture}[
    cell/.style={rectangle, draw=popline, minimum height=0.9cm, align=center, font=\small},
    header/.style={cell, fill=popblue, font=\bfseries\small, text=white},
    rowodd/.style={cell, fill=popblueL!30},
    roweven/.style={cell, fill=white},
    title/.style={font=\bfseries\large, text=popdark, node distance=0.3cm}
]
% En-têtes
\node[header, minimum width=2.2cm] (h1) {Méthode};
\node[header, minimum width=2.0cm, right=0pt of h1] (h2) {Mode d'action principal};
\node[header, minimum width=1.6cm, right=0pt of h2, align=center] (h3){Indice Pearl\\Typique (\%/an)};
\node[header, minimum width=1.3cm, right=0pt of h3, align=center] (h4){Protection\\IST/VIH};
\node[header, minimum width=1.4cm, right=0pt of h4] (h5) {Durée d'action};
\node[header, minimum width=1.8cm, right=0pt of h5, align=center] (h6){Accessibilité\\Cameroun};
\node[header, minimum width=1.7cm, right=0pt of h6] (h7) {Adapté Ado ?};

% Données - Ligne 1 : Naturelles
\node[rowodd, minimum width=2.2cm, below=0pt of h1, align=center] (r1c1){Naturelles\\ (Abstinence, MAMA,\\ Température, Billings)};
\node[rowodd, minimum width=2.0cm, right=0pt of r1c1, align=center] (r1c2){Éviter rapport\\dans période fertile\\(observation cycles)};
\node[rowodd, minimum width=1.6cm, right=0pt of r1c2] (r1c3) {15 -- 25};
\node[rowodd, minimum width=1.3cm, right=0pt of r1c3] (r1c4) {Non};
\node[rowodd, minimum width=1.4cm, right=0pt of r1c4] (r1c5) {Chaque cycle};
\node[rowodd, minimum width=1.8cm, right=0pt of r1c5, align=center] (r1c6){Gratuit,\\nécessite formation};
\node[rowodd, minimum width=1.7cm, right=0pt of r1c6, align=center] (r1c7){\textcolor{poporange}{Attention}\\(cycles irréguliers ado)};

% Ligne 2 : Barrières
\node[roweven, minimum width=2.2cm, below=0pt of r1c1, align=center] (r2c1){Barrières\\ Préservatif M / F,\\Diaphragme + spermicide};
\node[roweven, minimum width=2.0cm, right=0pt of r2c1, align=center] (r2c2){Barrière physique\\empêchant rencontre\\sperme--ovocyte};
\node[roweven, minimum width=1.6cm, right=0pt of r2c2] (r2c3) {2 -- 18};
\node[roweven, minimum width=1.3cm, right=0pt of r2c3, align=center] (r2c4){\textbf{OUI}\\(Seul moyen)};
\node[roweven, minimum width=1.4cm, right=0pt of r2c4] (r2c5) {Par rapport};
\node[roweven, minimum width=1.8cm, right=0pt of r2c5, align=center] (r2c6){Gratuit (CSB, ONG),\\Pharmacie (peu cher)};
\node[roweven, minimum width=1.7cm, right=0pt of r2c6, align=center] (r2c7){\textbf{OUI}\\(Essentiel double protection)};

% Ligne 3 : Hormonales Oestro-progestatives
\node[rowodd, minimum width=2.2cm, below=0pt of r2c1, align=center] (r3c1){Hormonales\\ combinées (Œstro-\\progestatives)\\Pilule, Patch, Anneau};
\node[rowodd, minimum width=2.0cm, right=0pt of r3c1, align=center] (r3c2){Blocage ovulation\\(inhibition pic LH)\\+ épaississement glaire};
\node[rowodd, minimum width=1.6cm, right=0pt of r3c2] (r3c3) {0,3 -- 6};
\node[rowodd, minimum width=1.3cm, right=0pt of r3c3] (r3c4) {Non};
\node[rowodd, minimum width=1.4cm, right=0pt of r3c4, align=center] (r3c5){Quotidien / Hebdo /\\Mensuel};
\node[rowodd, minimum width=1.8cm, right=0pt of r3c5, align=center] (r3c6){Peu cher\\(300--500 FCFA/mois)};
\node[rowodd, minimum width=1.7cm, right=0pt of r3c6, align=center] (r3c7){Oui si pas\\contre-indication\\œstrogènes};

% Ligne 4 : Hormonales Progestatives seules
\node[roweven, minimum width=2.2cm, below=0pt of r3c1, align=center] (r4c1){Hormonales\\ progestatives seules\\ Micro-pilule, Injectable,\\Implant, DIU Hormonal};
\node[roweven, minimum width=2.0cm, right=0pt of r4c1, align=center] (r4c2){Épaississement glaire\\cervicale (blocus)\\+ Atrophie endomètre\\+ (Parfois blocage ovulation)};
\node[roweven, minimum width=1.6cm, right=0pt of r4c2] (r4c3) {0,05 -- 3};
\node[roweven, minimum width=1.3cm, right=0pt of r4c3] (r4c4) {Non};
\node[roweven, minimum width=1.4cm, right=0pt of r4c4, align=center] (r4c5){Quotidien / 3 mois /\\3--5 ans / 5 ans};
\node[roweven, minimum width=1.8cm, right=0pt of r4c5, align=center] (r4c6){Injectable peu cher,\\Implant/DIU Gratuit\\ou subventionné (PF)};
\node[roweven, minimum width=1.7cm, right=0pt of r4c6, align=center] (r4c7){\textbf{OUI}\\(Idéal : observance facile)};

% Ligne 5 : Mécanique (DIU Cuivre)
\node[rowodd, minimum width=2.2cm, below=0pt of r4c1, align=center] (r5c1){Mécanique\\ DIU Cuivre};
\node[rowodd, minimum width=2.0cm, right=0pt of r5c1, align=center] (r5c2){Effet spermicide\\(inflammation locale)\\+ Empêche nidation};
\node[rowodd, minimum width=1.6cm, right=0pt of r5c2] (r5c3) {< 0,8};
\node[rowodd, minimum width=1.3cm, right=0pt of r5c3] (r5c4) {Non};
\node[rowodd, minimum width=1.4cm, right=0pt of r5c4] (r5c5) {5 -- 10 ans};
\node[rowodd, minimum width=1.8cm, right=0pt of r5c5, align=center] (r5c6){Gratuit / 1500 FCFA\\(CSB, Campagnes PF)};
\node[rowodd, minimum width=1.7cm, right=0pt of r5c6, align=center] (r5c7){\textbf{OUI}\\(Longue durée, pas\\d'hormones)};

% Ligne 6 : Chirurgicales
\node[roweven, minimum width=2.2cm, below=0pt of r5c1, align=center] (r6c1){Chirurgicales\\ (Définitives)\\Vasectomie (H),\\Ligature trompes (F)};
\node[roweven, minimum width=2.0cm, right=0pt of r6c1, align=center] (r6c2){Section canaux\\déferents / Trompes\\= Stérilisation};
\node[roweven, minimum width=1.6cm, right=0pt of r6c2] (r6c3) {< 0,1};
\node[roweven, minimum width=1.3cm, right=0pt of r6c3] (r6c4) {Non};
\node[roweven, minimum width=1.4cm, right=0pt of r6c4] (r6c5) {Définitive};
\node[roweven, minimum width=1.8cm, right=0pt of r6c5, align=center] (r6c6){Hôpital (chirurgie),\\coût variable};
\node[roweven, minimum width=1.7cm, right=0pt of r6c6, align=center] (r6c7){\textcolor{popred}{NON}\\(Adolescent = projet\\de vie non figé)};

% Ligne 7 : Urgence
\node[rowodd, minimum width=2.2cm, below=0pt of r6c1, align=center] (r7c1){Urgence\\ Pilule LNG 1,5 mg,\\DIU Cuivre};
\node[rowodd, minimum width=2.0cm, right=0pt of r7c1, align=center] (r7c2){Retarde ovulation\\(LNG) ou effet\\spermicide/nidation (DIU)};
\node[rowodd, minimum width=1.6cm, right=0pt of r7c2, align=center] (r7c3){Variable\\(selon délai)};
\node[rowodd, minimum width=1.3cm, right=0pt of r7c3] (r7c4) {Non};
\node[rowodd, minimum width=1.4cm, right=0pt of r7c4, align=center] (r7c5){Ponctuelle\\(1 rapport)};
\node[rowodd, minimum width=1.8cm, right=0pt of r7c5, align=center] (r7c6){Pharmacie, CSB,\\Gratuit en campagne};
\node[rowodd, minimum width=1.7cm, right=0pt of r7c6, align=center] (r7c7){Dépannage seul\\(Pas méthode régulière)};
\end{tikzpicture}
\legende{Tableau comparatif des principales méthodes contraceptives : classification, mécanisme d'action, efficacité (Indice de Pearl usage typique), protection IST, durée, accessibilité au Cameroun et pertinence pour l'adolescent. Source : OMS, MINESEC, PNUD Cameroun.}
\end{popfigure}

\courssub{Efficacité : l'Indice de Pearl, référence mondiale}

Comment comparer objectivement l'efficacité d'un préservatif et d'un implant ? On utilise l'\textbf{Indice de Pearl} (IP).

\begin{definition}[Indice de Pearl]
L'Indice de Pearl (IP) est le nombre de grossesses survenant chez 100 femmes utilisant une méthode donnée pendant \textbf{un an} (1200 mois-femme d'exposition).
\[
IP = \frac{\text{Nombre de grossesses} \times 1200}{\text{Nombre de femmes} \times \text{Mois d'exposition}}
\]
\emph{Plus l'indice est bas, plus la méthode est efficace.}
\end{definition}

\begin{attention}[Usage parfait vs Usage typique : le piège de l'observance]
L'indice de Pearl \textbf{théorique} (usage parfait) suppose une utilisation rigoureuse à chaque rapport ou prise quotidienne exacte. L'indice \textbf{réel} (usage typique) intègre les oublis, erreurs de pose, ruptures, arrêts intempestifs.
\begin{itemize}
    \item \textbf{Pilule combinée :} IP parfait = 0,3\% ; IP typique = 6 à 9\% (oubli fréquent chez l'ado).
    \item \textbf{Préservatif masculin :} IP parfait = 2\% ; IP typique = 13 à 18\% (mise en place tardive, rupture, glissement).
    \item \textbf{Implant / DIU :} IP parfait $\approx$ IP typique $\approx$ 0,05--0,8\% (\textbf{méthodes « oubli-proof »} : pas d'action quotidienne requise).
\end{itemize}
\textbf{Conclusion pour l'ado :} Privilégier les méthodes à \textbf{longue durée d'action (LARC)} — Implant, DIU — dont l'efficacité réelle ne dépend pas de la mémoire ou de la discipline quotidienne.
\end{attention}

\courssub{La double protection : stratégie recommandée pour les jeunes}

Un constat épidémiologique s'impose : au Cameroun, les IST (Chlamydia, Gonococcie, Syphilis, VIH) circulent activement chez les 15--24 ans. \textbf{Seul le préservatif (masculin ou féminin) protège simultanément contre les grossesses non désirées ET les IST/VIH.}

\begin{propriete}[Double protection]
L'association systématique \textbf{Préservatif (M ou F) + Autre méthode contraceptive efficace (Implant, DIU, Pilule, Injectable)} constitue la \textbf{double protection}. Elle est la stratégie de référence pour les adolescents et jeunes adultes sexuellement actifs.
\begin{itemize}
    \item \textbf{Préservatif} : Barrière contre VIH/IST + contraception de « secours » (efficacité modérée seule).
    \item \textbf{Autre méthode} : Contraception hautement efficace (IP < 1\%) qui prend le relais si le préservatif lâche.
\end{itemize}
\end{propriete}

\begin{exemplebox}[Cas concret : Stéphane et Mireille, 17 ans, lycéens à Yaoundé]
Stéphane utilise des préservatifs (gratuit au CSB), mais un soir, le préservatif se déchire. Mireille a un implant sous-cutané (pose gratuite lors d'une campagne « Planning Familial » au lycée).
\begin{itemize}
    \item \textbf{Risque IST/VIH :} Présent (rupture) $\rightarrow$ Dépistage conseillé à J+15 et J+45 au CSB.
    \item \textbf{Risque grossesse :} Quasi nul grâce à l'implant (IP 0,05\%).
    \item \textbf{Bilan :} La double protection a fonctionné. Sans implant, Mireille aurait eu besoin de la pilule du lendemain (stress, coût, délai 72h).
\end{itemize}
\end{exemplebox}

\courssub{Accessibilité et coûts réels au Cameroun : lever les freins financiers}

\begin{exemplebox}[Coûts approximatifs au public au Cameroun (2024--2025, hors consultation)]
\begin{center}
\begin{tabularx}{\linewidth}{|l|X|X|}\hline
\rowcolor{popblue} \textbf{Méthode} & \textbf{Coût public habituel (FCFA)} & \textbf{Points d'accès} \\ \hline
Préservatif Masculin & \textbf{Gratuit} (CSB, ONG, Campagnes) -- 100 (Pharmacie/Épicerie) & CSB, ONG (CAMNAFAW, ACMS), Pharmacies, Distributeurs automatiques \\ \hline
Préservatif Féminin & 100 -- 300 (souvent gratuit ONG) & ONG, Certains CSB, Pharmacies \\ \hline
Pilule combinée (1 plaquette) & 300 -- 500 & Pharmacies, CSB (dépôt légal) \\ \hline
Micro-pilule (1 plaquette) & 300 -- 600 & Pharmacies, CSB \\ \hline
Injectable (Dépo-Provera / Sayana Press, 3 mois) & 500 -- 1000 & CSB, Hôpitaux, ONG \\ \hline
Implant (Jadelle 5 ans / Implanon 3 ans) & \textbf{Gratuit -- 2000} (Campagnes PF, Fonds Mondial, Gouvernement) & CSB (Infirmiers/Sages-femmes formés), Hôpitaux, Cliniques mobiles \\ \hline
DIU Cuivre (10 ans) & \textbf{Gratuit -- 1500} (Campagnes PF) & CSB (Sages-femmes/Médecins), Hôpitaux \\ \hline
DIU Hormonal (Mirena 5 ans) & 15 000 -- 25 000 & Hôpitaux, Cliniques privées (moins accessible) \\ \hline
Pilule urgence (Lévonorgestrel 1,5 mg) & 500 -- 1500 (Gratuit certains CSB/ONG) & Pharmacies (sans ordonnance), CSB, ONG \\ \hline
\end{tabularx}
\end{center}
\emph{Note : La gratuité de l'implant et du DIU Cuivre est effective dans la quasi-totalité des CSB publics et lors des campagnes « Journées de Planification Familiale » organisées par le MINSANTÉ et ses partenaires (UNFPA, Fonds Mondial). L'argument « c'est trop cher » ne tient plus pour les méthodes les plus efficaces.}
\end{exemplebox}

\courssub{Critères de choix pour l'adolescent : l'approche « Choix Libre Éclairé »}

Choisir, ce n'est pas subir. L'OMS et le MINSANTÉ Cameroun promeuvent le \textbf{Conseil en Planification Familiale (CPF)} basé sur les droits : information complète, non-directivité, confidentialité, consentement libre.

\begin{methode}[Algorithme de conseil contraception (Approche « Choix Libre Éclairé »)]
\emph{À réaliser par un personnel de santé formé (infirmier, sage-femme, médecin) au CSB ou en consultation jeune.}
\begin{enumerate}
    \item \textbf{Exclure une grossesse en cours} (test urinaire $\beta$-hCG si doute, rappel date dernières règles).
    \item \textbf{Dépister les IST} (recherche symptômes, proposition dépistage VIH/syphilis/chlamydia/gonocoque) : traiter avant ou pendant la pose si nécessaire.
    \item \textbf{Évaluer les antécédents médicaux} (voir \textbf{Attention} ci-dessous : contre-indications aux œstrogènes).
    \item \textbf{Présenter les options} adaptées au profil (tableau comparatif) : efficacité, durée, effets secondaires prévisibles (règles irrégulières, aménorrhée, prise de poids modérée, maux de tête), mode d'administration.
    \item \textbf{Laisser la personne choisir} sans pression (même si le soignant a une préférence technique).
    \item \textbf{Expliquer l'utilisation pratique} (pose, retrait, signaux d'alerte, rendez-vous de contrôle).
    \item \textbf{Planifier le suivi} (revoir à 1 mois, puis 3--6 mois ; numéro d'urgence si effets secondaires graves).
\end{enumerate}
\end{methode}

\begin{propriete}[L'implant sous-cutané (Jadelle\textsuperscript{\textregistered}, Implanon NXT\textsuperscript{\textregistered}) : la méthode « idéale » pour l'adolescente]
\begin{itemize}
    \item \textbf{Efficacité maximale :} IP $\approx$ 0,05\% (la plus efficace des méthodes réversibles, équivalente à la stérilisation).
    \item \textbf{Durée :} 3 ans (Implanon) à 5 ans (Jadelle) $\rightarrow$ couvre toute la scolarité lycée/universitaire.
    \item \textbf{Observance nulle :} Pas de pilule quotidienne, pas de rendez-vous trimestriel $\rightarrow$ élimine le risque d'oubli, premier échec de la pilule chez l'ado.
    \item \textbf{Progestatif seul (Lévonorgestrel / Étongestrel) :} Pas d'œstrogène $\rightarrow$ \textbf{pas de risque thrombo-embolique}, utilisable même avec contre-indications aux œstrogènes (HTA, migraine avec aura, drépanocytose, tabagisme).
    \item \textbf{Réversibilité immédiate :} Retour de la fertilité dès le retrait (souvent dès le cycle suivant).
    \item \textbf{Gratuit au Cameroun} dans le cadre des programmes publics/ONG.
\end{itemize}
\textbf{Exigible au Probatoire :} Citer l'implant comme méthode de choix pour l'adolescente (efficacité, observance, sécurité, coût) et expliquer son mécanisme (progestatif $\rightarrow$ épaississement glaire cervicale + atrophie endomètre + inhibition ovulation variable).
\end{propriete}

\begin{attention}[Contre-indications majeures aux ŒSTROGÈNES (Pilule combinée, Patch, Anneau) -- \textbf{Critère d'exclusion absolu} chez l'adolescente]
La présence d'UN SEUL de ces facteurs oriente \textbf{obligatoirement} vers une méthode \textbf{sans œstrogène} (Progestatif seul : Implant, Injectable, Micro-pilule, DIU Hormonal/Cuivre) ou Barrière :
\begin{itemize}
    \item Antécédent personnel ou familial (1er degré) de \textbf{thrombose veineuse} (phlébite, embolie pulmonaire) ou \textbf{accident vasculaire cérébral}.
    \item \textbf{Hypertension artérielle} non contrôlée (HTA $\ge$ 140/90 mmHg).
    \item \textbf{Diabète} avec complications vasculaires (néphropathie, rétinopathie).
    \item \textbf{Migraine avec aura} (scotomes, troubles visuels/sensitifs précurseurs) -- \emph{très fréquent chez les jeunes filles}.
    \item \textbf{Cardiopathie valvulaire} (risque endocardite, thrombose).
    \item \textbf{Cancer hormono-dépendant} (sein, endomètre) en cours ou récent.
    \item \textbf{Tabagisme actif} + Âge $>$ 35 ans (rare en 1ère, mais à connaître pour l'examen).
    \item \textbf{Chirurgie majeure} avec immobilisation prolongée (risque thrombo-embolique majoré par les œstrogènes).
    \item \textbf{Allaitement} $<$ 6 semaines post-partum (œstrogènes diminuent la lactation).
\end{itemize}
\textbf{Piège Probatoire :} Ne pas confondre « contre-indication aux œstrogènes » et « contre-indication à toute contraception hormonale ». Les progestatifs seuls (Implant, Injectable, Micro-pilule, DIU hormonal) sont \textbf{autorisés} dans presque tous ces cas (sauf cancer du sein actif, maladie hépatique sévère aiguë, saignements génitaux inexpliqués).
\end{attention}

\courssub{Focus : Les DIU (Dispositifs Intra-Utérins) -- Mécanismes d'action comparés}

Le DIU Cuivre et le DIU Hormonal (LNG-IUS) sont tous deux des « méthodes à longue durée d'action » (LARC), très efficaces, mais leur mécanisme diffère. Comprendre cette différence aide à choisir et à expliquer les effets secondaires (règles abondantes vs aménorrhée).

\begin{popfigure}
\centering
\begin{tikzpicture}[
    scale=0.9,
    every node/.style={font=\small},
    uterus/.style={draw=popdark, thick, fill=poppinkL!30},
    tube/.style={draw=popdark, thick, fill=poppinkL!30},
    ovary/.style={draw=popdark, thick, fill=poporangeL!30},
    diu/.style={draw=popblue, very thick, fill=popblueL!50},
    arrow/.style={->, >=Stealth, thick, popdark},
    label/.style={font=\small, text=popdark, align=center},
    mech/.style={font=\scriptsize, text=poppurple, align=left}
]
% --- DIU CUIVRE (Gauche) ---
\begin{scope}[xshift=-5.5cm]
% Anatomie
\draw[uterus] (0,0) ellipse (1.2cm and 2.5cm); % Corps utérin
\draw[tube] (-1.2, 2.5) .. controls (-2, 3) and (-2, 3.5) .. (-1.5, 4); % Trompe G
\draw[tube] (1.2, 2.5) .. controls (2, 3) and (2, 3.5) .. (1.5, 4);   % Trompe D
\draw[ovary] (-2.2, 4) ellipse (0.5cm and 0.7cm); % Ovaire G
\draw[ovary] (2.2, 4) ellipse (0.5cm and 0.7cm);  % Ovaire D
% DIU Cuivre
\draw[diu] (0, -0.5) -- (0, 1.5); % Tige
\draw[diu] (-0.6, 0.5) -- (0.6, 0.5); % Bras transversaux
\draw[diu] (-0.1, -0.5) -- (0.1, -0.5); % Fil
% Filaments cuivre (symbolisés)
\foreach \y in {-0.3, 0.1, 0.5, 0.9, 1.3} \draw[poporange, thick] (-0.15,\y) -- (0.15,\y);
% Spermatozoïdes (flèches bloquées)
\draw[arrow, popred] (-2.5, 1) -- (-1.5, 1) node[left, label] {Spermatozoïdes};
\draw[popred, decorate, decoration={snake, amplitude=1pt, segment length=4pt}, thick] (-1.3, 0.8) -- (-0.5, 0.8);
\draw[popred, decorate, decoration={snake, amplitude=1pt, segment length=4pt}, thick] (-1.3, 1.2) -- (-0.5, 1.2);
% Légendes mécanismes
\node[label, anchor=north west] at (-3.5, 3.5) {\textbf{DIU Cuivre (ex: TCu-380A)}};
\node[mech, anchor=north west, align=center] at (-3.5, 3.0){
    \textbf{Mécanisme principal :} Effet spermicide local.\\
    Cu$^{2+}$ libérés $\rightarrow$ réaction inflammatoire stérile\\
    (leucocytes, prostaglandines, enzymes) $\rightarrow$ toxique pour\\
    spermatozoïdes (motilité, viabilité, capacitation) ET ovocyte.\\
    \textbf{Secondaire :} Empêche nidation (endomètre inflammatoire).\\
    \textbf{Pas d'effet hormonal systémique.}
};
\end{scope}

% --- DIU HORMONAL (Droite) ---
\begin{scope}[xshift=5.5cm]
% Anatomie
\draw[uterus] (0,0) ellipse (1.2cm and 2.5cm);
\draw[tube] (-1.2, 2.5) .. controls (-2, 3) and (-2, 3.5) .. (-1.5, 4);
\draw[tube] (1.2, 2.5) .. controls (2, 3) and (2, 3.5) .. (1.5, 4);
\draw[ovary] (-2.2, 4) ellipse (0.5cm and 0.7cm);
\draw[ovary] (2.2, 4) ellipse (0.5cm and 0.7cm);
% DIU Hormonal (T avec réservoir)
\draw[diu] (0, -0.5) -- (0, 1.5);
\draw[diu] (-0.6, 0.5) -- (0.6, 0.5);
\draw[poppurple, thick, fill=poppurpleL!50] (-0.1, 0.2) rectangle (0.1, 1.0); % Réservoir LNG
% Flèches diffusion hormone
\draw[arrow, poppurple] (0.3, 0.8) .. controls (0.8, 1) and (0.8, 0) .. (0.3, -0.2) node[right, mech] {Diffusion LNG};
\draw[arrow, poppurple] (-0.3, 0.8) .. controls (-0.8, 1) and (-0.8, 0) .. (-0.3, -0.2);
% Légendes mécanismes
\node[label, anchor=north east] at (3.5, 3.5) {\textbf{DIU Hormonal (ex: Mirena\textsuperscript{\textregistered} 52 mg LNG)}};
\node[mech, anchor=north east, align=center] at (3.5, 3.0){
    \textbf{Mécanisme principal :} Libération locale continue\\
    de Lévonorgestrel (LNG) $\approx$ 20 $\mu$g/j $\rightarrow$ 5 ans.\\
    1. \textbf{Épaississement glaire cervicale} : Blocus infranchissable\\
    pour spermatozoïdes (effet barrière chimique).\\
    2. \textbf{Atrophie endométriale} : Endomètre fin, inactif,\\   inhospitalier $\rightarrow$ empêche nidation.\\
    3. \textbf{Inhibition ovulation partielle} ($\sim$ 15--50\% cycles).\\
    \textbf{Effet systémique très faible} (taux plasmatiques bas).
};
\end{scope}
\end{tikzpicture}
\legende{Schémas comparés des mécanismes d'action : DIU Cuivre (effet inflammatoire spermicide local, sans hormone) vs DIU Hormonal (Lévonorgestrel local : blocus glaireux + atrophie endométriale). Le DIU Cuivre conserve des règles cycliques (souvent plus abondantes/douloureuses) ; le DIU Hormonal induit souvent une aménorrhée (règles absentes) thérapeutique pour les ménorragies.}
\end{popfigure}

\courssub{Contraception d'urgence : le filet de sécurité, pas la méthode habituelle}

\begin{savaistu}[Contraception d'urgence (CU) : Lévonorgestrel 1,5 mg (Norlevo\textsuperscript{\textregistered}, Optinor\textsuperscript{\textregistered}...)]
\begin{itemize}
    \item \textbf{Indication :} Rapport non ou mal protégé (oubli pilule $>$ 12h, rupture préservatif, viol, rapport non consenti). \textbf{PAS pour contraception régulière}.
    \item \textbf{Mécanisme :} \textbf{Retarde l'ovulation} (inhibe le pic LH pré-ovulatoire) si prise \textbf{avant le pic LH}. \textbf{N'est pas abortive} : n'interrompt pas une grossesse déjà implantée (n'agît pas sur l'endomètre si ovulation déjà passée).
    \item \textbf{Délai d'efficacité :} \textbf{Au plus vite} (idéalement $<$ 12h, efficacité $\approx$ 95\%). Efficacité décroît avec le temps : jusqu'à \textbf{72h (3 jours)} après le rapport (efficacité $\approx$ 58\% à 72h). Au-delà : DIU Cuivre possible jusqu'à 5 jours (120h) -- méthode d'urgence la plus efficace.
    \item \textbf{Effet « parapluie » ? \textbf{NON}.} La CU ne protège \textbf{PAS} pour les rapports sexuels ultérieurs (même dans les 24h qui suivent la prise). Il faut utiliser un préservatif ou reprendre/recommencer sa méthode habituelle avec précaution (barrière 7 à 14 jours selon méthode).
    \item \textbf{Accessibilité Cameroun :} Vente libre en pharmacie (500--1500 FCFA), Gratuit dans les CSB/ONG (viol, adolescents), souvent disponible chez les « Agents de Santé Communautaires » formés.
    \item \textbf{Effets secondaires :} Nausées (fréquentes), règles décalées (avance ou retard), saignements intermenstruels. Si vomissements $<$ 2h après prise $\rightarrow$ reprendre une dose.
    \item \textbf{Suivi :} Test de grossesse ($\beta$-hCG urinaire) si règles absentes à J+21 ou anormales. Dépistage IST proposé.
\end{itemize}
\end{savaistu}

\courssub{Exercice résolu : Choisir une méthode adaptée}

\begin{exoresolu}[Profil : Nadine, 18 ans, étudiante en 1ère C à Bafoussam]
\textbf{Énoncé :} Nadine a un petit ami régulier depuis 6 mois. Elle a des règles irrégulières (cycles 28--45 jours), des migraines avec aura (scotomes scintillants) depuis l'adolescence, et fume 3 cigarettes/jour. Elle oublie souvent de prendre ses médicaments. Elle ne veut pas d'enfants avant 5 ans. Elle consulte au CSB pour une contraception. En tant qu'infirmier-conseiller, quelle méthode proposez-vous en priorité ? Justifiez en citant 3 critères médicaux et 2 critères pratiques.

\textbf{Solution détaillée :}
\begin{enumerate}
    \item \textbf{Analyse des contre-indications (Critères médicaux) :}
    \begin{itemize}
        \item \textbf{Migraine avec aura :} \textbf{Contre-indication ABSOLUE aux œstrogènes} (risque AVC multiplié). $\rightarrow$ \textbf{Exclure :} Pilule combinée, Patch, Anneau.
        \item \textbf{Tabagisme (3 cig/jour) + Âge 18 ans :} Bien que le risque thrombotique majeur soit > 35 ans, le tabagisme potentialise le risque vasculaire des œstrogènes. $\rightarrow$ \textbf{Renforce l'exclusion des œstrogènes}.
        \item \textbf{Règles irrégulières :} Contre-indique les méthodes naturelles (Billings, Température, MAMA) qui nécessitent des cycles réguliers pour identifier la période fertile. $\rightarrow$ \textbf{Exclure :} Méthodes naturelles.
    \end{itemize}
    \item \textbf{Analyse du mode de vie (Critères pratiques) :}
    \begin{itemize}
        \item \textbf{Mauvaise observance (oubli fréquent) :} Exclut la pilule (quotidienne), le patch (hebdo), l'anneau (mensuel), l'injectable (rdv trimestriel strict). $\rightarrow$ \textbf{Privilégier méthode « oubli-proof » (LARC)}.
        \item \textbf{Projet de grossesse à 5 ans :} Méthode longue durée (3--5 ans) idéale. Réversibilité immédiate requise. $\rightarrow$ \textbf{Exclure :} Stérilisation chirurgicale.
    \end{itemize}
    \item \textbf{Choix optimal : \textbf{Implant sous-cutané (Jadelle 5 ans ou Implanon 3 ans)}.}
    \begin{itemize}
        \item Progestatif seul (LNG/ETG) : \textbf{Autorisé} migraine avec aura, tabagisme, règles irrégulières.
        \item Efficacité maximale (IP 0,05\%), durée 3--5 ans = couvre projet.
        \item Observance nulle (pose unique).
        \item Gratuit au CSB (campagne PF).
        \item Réversible immédiatement à l'arrêt.
    \end{itemize}
    \item \textbf{Double protection obligatoire :} Prescrire \textbf{Préservatifs masculins gratuits} en plus de l'implant (protection VIH/IST, car partenaire unique mais statut sérologique non prouvé récent).
    \item \textbf{Suivi :} Revoir à 1 mois (contrôle saignements, tolérance), puis annuel. Expliquer saignements irréguliers/aménorrhée fréquents et bénins.
\end{enumerate}
\end{exoresolu}

\begin{aretenir}[Synthèse : Choisir sa contraception en 1ère C / TI]
\begin{itemize}
    \item \textbf{Efficacité réelle (Usage typique) :} Implant / DIU ($<$ 1 grossesse/100 femmes/an) $\gg$ Pilule / Injectable (6--9) $\gg$ Préservatif (13--18) $\gg$ Naturelles (15--25).
    \item \textbf{Double protection = Préservatif + Méthode LARC (Implant/DIU) ou Pilule/Injectable.} Seule combinaison protégeant contre Grossesse + IST/VIH.
    \item \textbf{Contre-indications œstrogènes :} Migraine avec aura, HTA, Thrombose, Tabagisme $>$ 35 ans, Cancer hormono-dépendant $\rightarrow$ Progestatif seul ou Cuivre.
    \item \textbf{Accès Cameroun :} Implant et DIU Cuivre \textbf{gratuits} en CSB/Campagnes. Préservatifs gratuits. Pilule/Injectable très accessibles.
    \item \textbf{Choix Libre Éclairé :} Droit à l'information complète, au choix sans contrainte, à la confidentialité. Le/la soignant(e) conseille, l'adolescent(e) décide.
    \item \textbf{Urgence :} LNG 1,5 mg $\le$ 72h (mieux $<$ 24h). Pas abortive. Pas d'effet parapluie. DIU Cuivre $\le$ 120h (plus efficace).
\end{itemize}
\end{aretenir}

\begin{exercice}[Application : Cas cliniques à résoudre]
\textbf{Cas 1 :} Bruno, 17 ans, vient chercher des préservatifs au CSB. Il dit : « Je n'ai pas besoin d'autre chose, le préservatif protège de tout. » Que lui répondez-vous sur l'efficacité contraceptive réelle (usage typique) et l'intérêt de la double protection pour sa copine qui ne veut pas d'enfant ?
\vspace{0.5cm}
\textbf{Cas 2 :} Fatou, 16 ans, drépanocytaire (SS), règles abondantes et douloureuses. Elle est en couple stable. Son médecin lui déconseille la pilule combinée. Pourquoi ? Quelle méthode LARC serait \emph{thérapeutique} pour ses règles (aménorrhée) et contraceptive ? Justifiez.
\vspace{0.5cm}
\textbf{Cas 3 :} Calcul d'Indice de Pearl. Dans une cohorte de 250 adolescentes utilisant la pilule combinée pendant 8 mois, on observe 3 grossesses. Calculer l'IP. Interpréter par rapport à l'IP théorique (0,3\%). Que concluiez-vous sur l'observance ?
\end{exercice}

\begin{corrige}[Exercice : Cas cliniques]
\textbf{Cas 1 :} Bruno a raison sur les IST/VIH (préservatif = seul protecteur), mais \textbf{tort sur la contraception}. IP typique préservatif = 13--18\%. Sur 1 an, ~15 couples sur 100 auront une grossesse. Pour sa copine (enjeu scolarité), risque inacceptable. \textbf{Conseil :} Double protection = Préservatif (IST) + Implant/DIU (Contraception top efficacité).
\vspace{0.3cm}
\textbf{Cas 2 :} Drépanocytose (SS) = risque thrombose veineuse majoré + risque crise vaso-occlusive. \textbf{Œstrogènes contre-indiqués} (pro-thrombotiques). \textbf{Méthode idéale : DIU Hormonal (LNG-IUS / Mirena\textsuperscript{\textregistered})}. Mécanisme : LNG local $\rightarrow$ atrophie endométriale $\rightarrow$ \textbf{aménorrhée thérapeutique} (arrêt règles = arrêt pertes de fer, amélioration anémie, moins de crises douloureuses). Contraception IP $<$ 0,1\%. Durée 5 ans. Alternative : Implant (aménorrhée fréquente mais moins prévisible que DIU hormonal).
\vspace{0.3cm}
\textbf{Cas 3 :} $IP = \frac{3 \times 1200}{250 \times 8} = \frac{3600}{2000} = 1,8\%$. IP théorique = 0,3\%. IP réel (1,8\%) = 6 fois l'IP théorique. \textbf{Conclusion :} Observance imparfaite (oubli pilules, vomissements, interactions médicamenteuses non gérées). Confirme l'intérêt des LARC (Implant/DIU) dont IP réel $\approx$ IP théorique.
\end{corrige}

\coursec{Stratégies de prévention et alternatives responsables}

Après avoir passé en revue les moyens de contraception disponibles (section 5), il est essentiel de comprendre que la simple connaissance technique ne suffit pas. Au Cameroun, comme le montrent les données de l'EDS-MICS 2018, le taux d'utilisation de la contraception moderne chez les filles de 15 à 19 ans en union n'atteint qu'environ \textbf{20\%}, pour un besoin non satisfait demeurant très élevé (supérieur à 30\%). Ce fossé entre \emph{savoir} et \emph{faire} s'explique par des barrières sociales, culturelles, économiques et psychologiques. Cette section propose une approche globale, ancrée dans les droits et la réalité du terrain camerounais, pour transformer les connaissances en comportements protecteurs durables.

\courssub{L'approche ABCD : une stratégie graduée et réaliste}

L'approche \cle{ABCD}, promue par l'OMS et adaptée au contexte camerounais par le Ministère de la Santé Publique (MINSANTÉ), offre un cadre pragmatique qui respecte la diversité des parcours des jeunes. Elle ne présente pas l'abstinence comme unique option morale, mais comme le premier niveau de protection, complété par d'autres stratégies selon la situation réelle de chaque jeune.

\begin{definition}[Approche ABCD]
L'approche \textbf{ABCD} structure la prévention du VIH, des IST et des grossesses non désirées en quatre piliers complémentaires :
\begin{itemize}
    \item \textbf{A -- Abstinence / Report du premier rapport :} Choix de ne pas avoir de rapports sexuels (pénétratifs) pour éliminer totalement le risque. Pour les jeunes déjà sexuellement actifs, cela peut signifier une période d'abstinence secondaire (repos sexuel).
    \item \textbf{B -- Be faithful (Fidélité / Réduction du nombre de partenaires) :} Engagement dans une relation mutuellement monogame avec un partenaire de statut sérologique connu et négatif, ou réduction drastique du nombre de partenaires pour diminuer l'exposition au risque.
    \item \textbf{C -- Condom (Préservatif masculin ou féminin à chaque rapport) :} Utilisation correcte et systématique du préservatif, seule méthode offrant une \textbf{double protection} (contraception + prévention VIH/IST).
    \item \textbf{D -- Dépistage (VIH/IST régulier) :} Connaissance de son statut sérologique et dépistage des IST (syphilis, chlamydia, gonococcie, etc.) au moins une fois par an, ou après toute prise de risque, pour accéder précocement au traitement et rompre les chaînes de transmission.
\end{itemize}
\end{definition}

\begin{experience}[Évaluation de l'approche ABCD dans un lycée de Yaoundé]
\textbf{Problématique :} Une enquête KAP (Knowledge, Attitudes, Practices) réalisée en 2022 dans un établissement public de Yaoundé auprès de 400 élèves de Première montre que 68\% connaissent l'approche ABCD, mais seulement 22\% déclarent l'appliquer de façon cohérente.
\textbf{Hypothèse :} La connaissance théorique ne se traduit pas en pratique faute de compétences psychosociales (négociation, estime de soi) et d'accès facile aux services.
\textbf{Protocole :} Comparaison de deux groupes : groupe témoin (cours magistral seul) vs groupe intervention (cours + ateliers de \emph{life skills} + orientation vers SAJ).
\textbf{Résultats à 6 mois :} Le groupe intervention montre une augmentation de l'utilisation du préservatif au dernier rapport (45\% vs 28\%, $p<0,01$) et une hausse du dépistage VIH volontaire (38\% vs 15\%).
\textbf{Interprétation :} L'information seule est insuffisante ; le développement de compétences de vie et la liaison effective avec des services adaptés sont déterminants.
\textbf{Conclusion :} L'approche ABCD doit être enseignée comme un \emph{menu d'options} à activer selon le contexte personnel, et non comme une injonction morale unique.
\end{experience}

\begin{popfigure}
\begin{tikzpicture}[scale=0.9, every node/.style={font=\small}]
    % Centre
    \node[circle, draw=popblue, thick, fill=popblue!10, minimum width=3.2cm, align=center] (centre) {Santé\\Reproductive\\des Jeunes};
    
    % 4 rayons principaux
    \foreach \angle/\label/\desc in {
        90/Information (CSE)/{Programmes scolaires\\on jugeants, basés\\sur les droits},
        0/Services (SAJ)/{CSB, CEDA, ONG\\Confidentialité,\\gratuité, horaires},
        -90/Compétences (Vie)/{Négociation, refus,\\estime de soi,\\consentement},
        180/Environnement/{Parents, communauté,\\lois, médias,\\leaders d'opinion}
    }{
        \draw[popblue, thick, ->] (centre) -- ++(\angle:3.5cm) node[anchor=center] (n\angle) {};
        \node[align=center, text width=3.2cm] at (\angle:4.5cm) {\textbf{\label}\\\desc};
    }
    
    % Flèches circulaires entre les piliers (synergie)
    \foreach \a/\b in {90/0, 0/-90, -90/180, 180/90} {
        \draw[popgreen, thick, ->, bend left=15] (\a:3.8cm) to (\b:3.8cm);
    }
    
    % Légende des couleurs
    \node[anchor=north west, align=left, fill=white, draw=popline, rounded corners] at (-7,-4.5) {
        \textcolor{popblue}{$\blacktriangleright$} Piliers structurels \quad
        \textcolor{popgreen}{$\blacktriangleright$} Synergie / Renforcement mutuel
    };
\end{tikzpicture}
\legende{Roue de la prévention : la santé reproductive des jeunes repose sur l'articulation de quatre piliers indissociables. L'information (CSE) éclaire les choix, les services (SAJ) rendent l'action possible, les compétences de vie permettent de passer à l'acte, et l'environnement (famille, communauté, cadre légal) soutient ou freine la démarche. Les flèches vertes illustrent la synergie : aucun pilier ne suffit seul.}
\end{popfigure}

\courssub{Éducation sexuelle complète (CSE) : fondement du pouvoir d'agir}

L'Éducation Sexuelle Complète (\cle{CSE}, \emph{Comprehensive Sexuality Education}) n'est pas un simple cours de biologie. C'est un processus d'apprentissage tout au long de la vie, basé sur un curriculum structuré, scientifiquement exact, adapté à l'âge, non jugeant, et ancré dans les \textbf{droits humains} (droit à la santé, à l'information, à la non-discrimination, à l'intégrité corporelle).

\begin{propriete}[Caractéristiques d'une CSE de qualité (standards UNESCO/OMS)]
Une CSE efficace doit être :
\begin{enumerate}
    \item \textbf{Scientifiquement exacte :} Données validées sur la physiologie, la contraception, les IST, le VIH.
    \item \textbf{Basée sur les droits :} Promotion de l'égalité de genre, du consentement, de la non-violence, de la diversité.
    \item \textbf{Non jugeante :} Ne stigmatise pas la sexualité, l'orientation sexuelle, le statut sérologique, le handicap.
    \item \textbf{Adaptée à l'âge et au développement :} Contenus progressifs de la primaire au secondaire.
    \item \textbf{Centrée sur les compétences de vie (Life Skills) :} Pensée critique, prise de décision, communication, négociation, gestion des émotions, recherche d'aide.
    \item \textbf{Sensible au genre :} Déconstruit les normes de genre nuisibles (masculinité toxique, soumission féminine attendue).
    \item \textbf{Contextualisée :} Intègre les réalités culturelles, religieuses, légales du Cameroun (ex. : âge du consentement, loi sur l'avortement, mariage précoce).
\end{enumerate}
\end{propriete}

\begin{attention}[Piège fréquent : Confondre CSE et « encouragement à la sexualité précoce »]
\emph{Idée reçue :} « Parler de sexe aux jeunes les pousse à en faire plus tôt. »
\emph{Preuve scientifique :} Les méta-analyses (UNESCO 2018, OMS) montrent que la CSE \textbf{retarde} l'âge du premier rapport, \textbf{réduit} le nombre de partenaires, \textbf{augmente} l'utilisation du préservatif et de la contraception, et \textbf{diminue} les grossesses non désirées et les IST. Le silence, lui, laisse place aux rumeurs, à la pornographie comme unique éducation, et aux risques accrus.
\end{attention}

\courssub{Communication Parents-Adolescents et Pairs Éducateurs : briser le tabou}

Au Cameroun, la sexualité reste souvent un sujet tabou en famille. Pourtant, les études (EDS, MICS) montrent que les adolescents qui discutent de sexualité avec leurs parents (ou adultes de référence) ont des comportements plus protecteurs.

\begin{methode}[Comment engager le dialogue Parents-Ados (guide pratique)]
\begin{enumerate}
    \item \textbf{Créer un climat de confiance :} Disponibilité, écoute active sans interruption, suspension du jugement moral immédiat.
    \item \textbf{Saisir les occasions naturelles :} Scène de film, actualité, grossesse dans l'entourage, question de l'enfant. Éviter la « grande discussion » unique et anxiogène.
    \item \textbf{Utiliser le vocabulaire anatomique correct :} Pénis, vulve, vagin, utérus, spermatozoïde, ovule. Cela normalise le corps et protège contre les abus (l'enfant peut nommer ce qui lui arrive).
    \item \textbf{Valider les émotions :} « C'est normal d'avoir des questions / d'être curieux / d'avoir peur. »
    \item \textbf{Donner des informations factuelles :} Si on ne sait pas, on dit : « Bonne question, on va chercher ensemble sur un site fiable (ex. : \texttt{www.sante-adolescent.cm}, \texttt{www.onusida.org}) ou on ira au CSB. »
    \item \textbf{Aborder le consentement et les limites :} « Ton corps t'appartient. Personne n'a le droit de te toucher sans ton accord. Toi non plus. »
    \item \textbf{Orienter vers les services :} Connaître le CSB, le CEDA, le SAJ le plus proche, les numéros verts.
\end{enumerate}
\end{methode}

\begin{savaistu}[Le rôle clé des Pairs Éducateurs (PE) au Cameroun]
Les \textbf{Pairs Éducateurs}, formés par le MINSANTÉ, le MINJEC, le MINESEC et des ONG (CAMNAFAW, ACMS, CARE, Plan International), sont des jeunes de 15 à 24 ans qui mènent des activités de sensibilisation, de distribution de préservatifs (dépôt communautaire), d'orientation vers les SAJ et de suivi de l'observance (ARV, contraception) auprès de leurs pairs. Leur langage, leurs codes culturels et leur accessibilité (même hors école) en font des relais \textbf{incontournables} pour atteindre les jeunes non scolarisés, les filles mariées précocement, les jeunes des rues, les UDPS (Usagers de Drogues Par Injection). En 2023, le réseau national comptait plus de 5 000 PE actifs.
\end{savaistu}

\courssub{Services « Amis des Jeunes » (SAJ) : lever les barrières d'accès}

Un jeune qui a décidé de se protéger (préservatif, contraception, dépistage) se heurte souvent à des obstacles concrets : peur d'être jugé par le personnel, horaires incompatibles avec l'école, coût, absence de confidentialité (dossier visible par les parents), distance.

\begin{definition}[Services Amis des Jeunes (SAJ) -- Standards MINSANTÉ/OMS]
Un SAJ est un service de santé (public, confessionnel, associatif, privé) qui garantit :
\begin{itemize}
    \item \textbf{Confidentialité absolue :} Dossier médical inaccessible aux tiers (parents, enseignants), entretien en privé, pas de notification parentale obligatoire pour les mineurs (sauf danger vital immédiat).
    \item \textbf{Gratuité ou coût très faible :} Préservatifs gratuits, contraceptifs subventionnés (pilule, injectable, implant), dépistage VIH/IST gratuit, contraception d'urgence gratuite.
    \item \textbf{Accueil non jugeant :} Personnel formé à l'approche « youth-friendly » (pas de sermon, pas de refus au motif de l'âge ou du statut marital).
    \item \textbf{Horaires adaptés :} Ouverture le soir (après les cours), le week-end, pendant les vacances.
    \item \textbf{Package de services intégrés :} Contraception, dépistage/traitement IST, conseil VIH, prise en charge violences sexuelles, contraception d'urgence, conseil prénatal, PMTCT, santé mentale, nutrition.
    \item \textbf{Participation des jeunes :} Jeunes dans le comité de gestion, boîtes à suggestions, évaluation mystère.
\end{itemize}
\end{definition}

\begin{exemplebox}[Principaux réseaux SAJ au Cameroun]
\begin{center}
\begin{tabularx}{\linewidth}{|l|X|X|}\hline
\textbf{Structure} & \textbf{Spécificités} & \textbf{Localisation type} \\ \hline
\textbf{CSB (Centres de Santé de Base) publics} & Intégrés au système de santé, gratuité ciblée, personnel d'État & Partout (zones urbaines, périurbaines, rurales) \\ \hline
\textbf{CEDA (Centres d'Écoute et de Dépistage Anonymes)} & Spécialisés VIH/IST, anonymat total, conseil pré/post-test, ARV & Chefs-lieux de région/département \\ \hline
\textbf{CAMNAFAW (Cameroun National Association for Family Welfare)} & Réseau historique IPPF, cliniques mobiles, forte composante jeunes, avortement sécurisé (selon loi) & Grandes villes, zones rurales via mobiles \\ \hline
\textbf{ACMS (Association Camerounaise pour le Marketing Social)} & Marketing social : préservatifs \textit{Prudence}, \textit{Love}, pilule \textit{Confiance}, injectable \textit{Depo-Provera} à prix social, distribution communautaire & Pharmacies, dépôts, boutiques, PE, kiosques \\ \hline
\textbf{ONG spécialisées (CARE, Plan, MSF, etc.)} & Projets ciblés : filles hors école, UDPS, jeunes réfugiés, zones de conflit (NOSO) & Zones d'intervention projets \\ \hline
\end{tabularx}
\end{center}
\end{exemplebox}

\begin{attention}[Droit d'accès des mineurs aux SAJ -- Cadre légal camerounais]
La \textbf{Loi n°2018/012 du 11 juillet 2018} relative à la santé reproductive (Art. 14) stipule : « Tout adolescent a droit à l'information, à l'éducation et aux services de santé reproductive adaptés à son âge. » La \textbf{Circulaire n°00000664/MINSANTÉ/SG/DSS du 14/02/2019} précise que les mineurs de 15 ans et plus peuvent consentir seuls aux actes de santé reproductive (contraception, dépistage VIH, traitement IST, IVG si indication légale). \textbf{Refuser un service à un mineur de 15 ans au motif de l'absence de parent est une violation de la loi.} Pour les moins de 15 ans, l'appréciation de la maturité (« mineur mature ») par le soignant prime, dans l'intérêt supérieur de l'enfant.
\end{attention}

\courssub{Gestion concrète d'une situation à risque : la « check-list » post-exposition}

Malgré la prévention, une situation à risque peut survenir (oubli de pilule, rupture de préservatif, rapport non protégé sous emprise d'alcool, viol). La rapidité et la séquence des actions sont vitales.

\begin{methode}[Conduite à tenir après un rapport sexuel non protégé ou à risque (J0 = jour du rapport)]
\begin{enumerate}
    \item \textbf{J0 à J3 (idéalement < 12h, max 72h / 120h selon produit) : Contraception d'urgence (CU).}
    \begin{itemize}
        \item Pilule LNG 1,5 mg (dose unique) : efficace jusqu'à 72h (3 jours), efficacité décroissante.
        \item Pilule UPA 30 mg (ellaOne\textsuperscript{\textregistered}) : efficace jusqu'à 120h (5 jours), plus efficace que LNG après 72h et chez les femmes > 70 kg.
        \item \textbf{DIU au cuivre (urgence) :} Pose possible jusqu'à 5 jours (120h), efficacité \textasciitilde 99\%, offre une contraception durable immédiate. \textbf{Meilleure option si délai > 72h ou IMC > 25.}
        \item \textbf{Disponibilité :} Gratuite en SAJ, CSB, CEDA ; en pharmacie (OTC) \textasciitilde 500--1000 FCFA (LNG), \textasciitilde 3000 FCFA (UPA).
    \end{itemize}
    \item \textbf{J0 à J7 (si symptômes) / J0 à J3 (prophylaxie post-exposition VIH - PPE) : Dépistage / Traitement IST \& PPE VIH.}
    \begin{itemize}
        \item \textbf{IST symptomatiques (écoulement, ulcération, brûlures mictionnelles, douleurs pelviennes) :} Consultation immédiate pour traitement syndromique (OMS) ou étiologique. \textbf{Traiter le partenaire(s) simultanément} (notification assistée).
        \item \textbf{PPE VIH (Truvada\textsuperscript{\textregistered} ou équivalent 28 jours) :} Indiquée si risque avéré (partenaire séropositif non supprimé, viol, rupture préservatif avec partenaire statut inconnu/positif non traité). \textbf{Délai maximal : 72h, idéalement < 4h.} Disponible gratuitement aux CEDA, hôpitaux de référence, certains SAJ.
    \end{itemize}
    \item \textbf{J21 à J45 (6 semaines) : Dépistage VIH de contrôle (test 4ème génération Ag/Ac).}
    \begin{itemize}
        \item Fenêtre sérologique : 2 à 6 semaines. Test à 6 semaines (J42) = fiable à > 99\%. Test à 3 mois = certitude totale.
        \item Conseil pré et post-test obligatoire.
    \end{itemize}
    \item \textbf{J30 (1 mois) : Test de grossesse (urinaire ou $\beta$-hCG sanguine) si règles en retard.}
    \item \textbf{Parallèle : Démarche contraception durable (implant, DIU, injectable, pilule) + Double protection (Préservatif + Contraception hormonale/DIU).}
    \item \textbf{Soutien psychologique / social :} Si viol, contrainte, anxiété majeure, culpabilité $\rightarrow$ Orientation vers psychologue clinicien, assistant social, association d'aide aux victimes (ex. : ALVF, RENATA), numéro vert.
\end{enumerate}
\end{methode}

\begin{popfigure}
\begin{tikzpicture}[scale=0.85, every node/.style={font=\small}]
    % Nœud central
    \node[ellipse, draw=poporange, thick, fill=poporange!10, minimum width=4cm, minimum height=1.8cm, align=center] (centre) {RAPPORT NON PROTÉGÉ\\OU ACCIDENT (J0)};
    
    % 4 branches principales
    % Branche 1 : Contraception urgence
    \node[rectangle, draw=popblue, thick, fill=popblue!10, rounded corners, align=center, text width=3.5cm] (cu) at (-6, 2.5) 
        {\textbf{1. CONTRACEPTION D'URGENCE}\\\textbf{J0 -- J3 (max J5)}\\\textbullet\ LNG 1,5 mg (J0--J3)\\\textbullet\ UPA 30 mg (J0--J5)\\\textbullet\ DIU Cu (J0--J5) \textbf{++}\\\textit{Gratuit en SAJ/CSB/CEDA}};
    \draw[popblue, thick, ->] (centre.west) -- ++(-1,0) |- (cu.south east);
    
    % Branche 2 : Dépistage / PPE
    \node[rectangle, draw=popgreen, thick, fill=popgreen!10, rounded corners, align=center, text width=3.5cm] (dep) at (-6, -2.5) 
        {\textbf{2. DÉPISTAGE / PPE VIH}\\\textbf{IST : J0--J7 si symptômes}\\\textbullet\ Traitement syndromique OMS\\\textbullet\ Notification partenaire\\\textbf{PPE VIH : J0--J3 (max 72h)}\\\textbullet\ Truvada 28 j\\\textit{Gratuit CEDA/Hôpital}};
    \draw[popgreen, thick, ->] (centre.west) -- ++(-1,0) |- (dep.north east);
    
    % Branche 3 : Contraception durable
    \node[rectangle, draw=poppurple, thick, fill=poppurple!10, rounded corners, align=center, text width=3.5cm] (cd) at (6, 2.5) 
        {\textbf{3. CONTRACEPTION DURABLE}\\\textbf{Dès J0 / J1}\\\textbullet\ Implant (3--5 ans)\\\textbullet\ DIU Cu / Hormonal (5--10 ans)\\\textbullet\ Injectable (3 mois)\\\textbullet\ Pilule (quotidien)\\\textbf{+ PRÉSERVATIF = DOUBLE PROTECTION}};
    \draw[poppurple, thick, ->] (centre.east) -- ++(1,0) |- (cd.south west);
    
    % Branche 4 : Soutien psycho-social
    \node[rectangle, draw=poppink, thick, fill=poppink!10, rounded corners, align=center, text width=3.5cm] (psy) at (6, -2.5) 
        {\textbf{4. SOUTIEN PSYCHO-SOCIAL}\\\textbf{Si besoin (viol, anxiété...)}\\\textbullet\ Écoute / Conseil (SAJ, PE)\\\textbullet\ Psychologue clinicien\\\textbullet\ Assistant social\\\textbullet\ Associations (ALVF, RENATA)\\\textbullet\ \textbf{1510 (Num. vert)}};
    \draw[poppink, thick, ->] (centre.east) -- ++(1,0) |- (psy.north west);
    
    % Flèches de liaison entre branches (continuité)
    \draw[popmuted, dashed, ->] (cu.south) -- ++(0,-0.5) -| (dep.north);
    \draw[popmuted, dashed, ->] (dep.east) -- ++(1,0) |- (cd.west);
    \draw[popmuted, dashed, ->] (cd.south) -- ++(0,-0.5) -| (psy.north);
    
    % Légende
    \node[anchor=north, fill=white, draw=popline, rounded corners, align=center] at (0, -5.2) {
        \textbf{Chronologie critique :} \textcolor{popblue}{CU} \textbullet\ \textcolor{popgreen}{PPE/IST} \textbullet\ \textcolor{poppurple}{Contraception durable} \textbullet\ \textcolor{poppink}{Soutien} \quad
        \textcolor{popmuted}{$\dashrightarrow$} = Enchaînement recommandé
    };
\end{tikzpicture}
\legende{Carte mentale : conduite à tenir face à un rapport non protégé. Chaque branche est une urgence à déclencher sans attendre les autres. La double protection (préservatif + contraception durable) est la norme de sortie de crise. Le numéro vert 1510 (ou 116 pour les violences) offre une orientation anonyme 24h/24.}
\end{popfigure}

\courssub{La négociation du préservatif : la méthode DESC en pratique}

Savoir \emph{qu'il faut} mettre un préservatif ne suffit pas ; il faut savoir \emph{le faire accepter} ou \emph{refuser le rapport} sans le mettre. La méthode \cle{DESC} (Décrire, Exprimer, Spécifier, Conséquences), issue de la communication assertive, est un outil éprouvé pour les adolescents.

\begin{methode}[Technique de négociation DESC -- Exemple appliqué]
\textbf{Situation :} Ton/ta partenaire refuse le préservatif (« Ça gâche le plaisir / On se fait confiance / Je n'en ai pas »).
\begin{enumerate}
    \item \textbf{D -- Décrire le fait (objectif, sans accusation) :} \emph{« On n'a pas de préservatif. / Le préservatif est resté dans le sac. / On n'en a pas acheté. »}
    \item \textbf{E -- Exprimer son sentiment / sa peur (subjectif, « je ») :} \emph{« J'ai peur de tomber enceinte. / J'ai peur d'attraper le VIH ou une IST. / Je ne me sens pas en sécurité. »}
    \item \textbf{S -- Spécifier la solution concrète (proposition claire) :} \emph{« On s'arrête là pour ce soir. / On va en acheter au kiosque d'à côté (ouvert 24h). / On fait autre chose (caresses, masturbation mutuelle) sans pénétration. »}
    \item \textbf{C -- Conséquences positives (gain mutuel) :} \emph{« Comme ça, on profite à fond sans stress, on se respecte, et on se revoit sereinement. / Si on va en acheter, on rigolera du détour. »}
\end{enumerate}
\textbf{Si le refus persiste :} \emph{« Sans préservatif, pas de rapport. C'est non négociable pour ma santé. »} $\rightarrow$ \textbf{Refus du rapport non désiré = Droit absolu.}
\end{methode}

\begin{exoresolu}[Mise en situation : Négociation dans un contexte de pression]
\textbf{Énoncé :} Marie, 17 ans, élève en 1ère C, sort avec Paul, 19 ans, apprenti. Ce samedi, chez Paul (parents absents), l'ambiance devient intime. Paul sort un préservatif, mais il est périmé (date dépassée de 8 mois). Paul dit : « C'est bon, c'est encore solide, on n'a pas le temps d'aller en acheter, les boutiques sont fermées. Fais-moi confiance, je me retire avant. » Marie a très envie de faire l'amour, mais elle a peur. Elle n'a pas de contraception hormonale. Que doit-elle faire ? Justifier chaque étape.

\textbf{Solution détaillée :}
\begin{enumerate}
    \item \textbf{Analyse du risque :} Préservatif périmé = risque de rupture élevé (latex fragilisé). Retrait (coït interrompu) = \textbf{pas une méthode contraceptive fiable} (pré-éjaculation = spermatozoïdes, IST, VIH). Risque grossesse \textbf{réel et élevé} + risque IST/VIH.
    \item \textbf{Application DESC :}
    \begin{itemize}
        \item \textbf{D :} « Ce préservatif est périmé depuis 8 mois. Le retrait n'est pas une protection. »
        \item \textbf{E :} « J'ai vraiment peur de tomber enceinte ce soir, et je tiens à ma santé. »
        \item \textbf{S :} « On s'arrête pour la pénétration ce soir. On peut faire des câlins, se toucher autrement. Demain matin, on va ensemble au SAJ / à la pharmacie de garde pour avoir des préservatifs valides et je me fais poser un implant / DIU. »
        \item \textbf{C :} « On profitera pleinement la prochaine fois, sans peur, et je serai protégée pour longtemps. »
    \end{itemize}
    \item \textbf{Si Paul insiste / se fâche / menace de rompre :} Marie affirme : « Non, c'est non. Ma santé et mon avenir passent avant ce rapport. Si tu m'aimes, tu respectes mon choix. » Elle sort de la situation (appelle une amie, un taxi, rentre chez elle).
    \item \textbf{Suivi (J0) :} Marie se rend au SAJ le plus proche (ouvert week-end) pour : CU (LNG ou UPA si < 72h/120h) + Dépistage IST/VIH de base + Pose implant/DIU (contraception durable) + Conseil + Préservatifs gratuits.
    \item \textbf{Bilan :} Marie a exercé son \textbf{consentement libre, éclairé, réversible}. Elle a utilisé ses \textbf{compétences de vie} (assertivité, connaissance des risques, connaissance des services). Elle a transformé une situation à risque en prise en charge globale de sa santé reproductive.
\end{enumerate}
\end{exoresolu}

\courssub{La double protection : la norme recommandée pour les adolescents sexuellement actifs}

\begin{propriete}[Double protection : Définition et justification]
La \textbf{double protection} consiste en l'utilisation \textbf{simultanée et systématique} de :
\begin{itemize}
    \item \textbf{Un préservatif (masculin ou féminin)} $\rightarrow$ Protection contre le VIH, les IST, et contraception (efficacité typique \textasciitilde 85\%).
    \item \textbf{Une méthode contraceptive hautement efficace (LARC de préférence : Implant, DIU Cu ; ou Injectable, Pilule)} $\rightarrow$ Contraception de secours si échec préservatif (efficacité \textasciitilde 99\%+).
\end{itemize}
\textbf{Pourquoi est-ce la NORME pour les adolescents ?}
\begin{enumerate}
    \item \textbf{Fertilité maximale :} Pic de fécondité 15--24 ans. Échec contraceptif = grossesse à très fort impact (scolarité, santé, économie).
    \item \textbf{Exposition IST/VIH élevée :} Nouveaux partenaires, relations multiples, pouvoir de négociation faible, biologie (col utérin immature = porte d'entrée IST).
    \item \textbf{Observance imparfaite :} Oubli pilule, retard injectable, rupture préservatif. La redondance compense les failles de chaque méthode.
    \item \textbf{Recommandation OMS / MINSANTÉ / MINESEC :} Tout conseil contraception chez l'adolescent doit inclure l'offre systématique de préservatifs et l'explication de la double protection. \textbf{Ne pas la proposer est une faute de conseil professionnel.}
\end{enumerate}
\end{propriete}

\begin{attention}[Erreurs fréquentes à ne pas commettre en conseil]
\begin{itemize}
    \item \textbf{Dire : « Le préservatif suffit, pas besoin d'autre chose. »} $\rightarrow$ Faux : efficacité typique 85\% = 15 grossesses / 100 femmes/an. Inacceptable pour une lycéenne.
    \item \textbf{Dire : « Prends la pilule / l'implant, comme ça pas besoin de capote. »} $\rightarrow$ Faux : Zéro protection VIH/IST. Responsabilité partagée ignorée.
    \item \textbf{Ne pas proposer le DIU ou l'implant aux nullipares (jamais enceintes).} $\rightarrow$ Faux : OMS Catégorie 1 (aucune restriction) pour implant et DIU Cu chez les adolescentes nullipares. Avantages : efficacité maximale, oubli impossible, discrétion, réversibilité immédiate.
    \item \textbf{Oublier la contraception d'urgence comme « filet de sécurité » connu à l'avance.} $\rightarrow$ Chaque jeune sexuellement actif doit savoir où et comment l'obtenir en < 12h.
\end{itemize}
\end{attention}

\courssub{Le rôle pivot de l'école : un écosystème protecteur}

L'école n'est pas seulement un lieu d'enseignement ; c'est un \textbf{lieu de vie} où se jouent des dynamiques de pouvoir, de séduction, de violence, mais aussi de solidarité. Le programme SVTEEHB de Première est la colonne vertébrale cognitive, mais il doit s'appuyer sur un dispositif institutionnel cohérent.

\begin{definition}[Dispositif « École promotrice de santé reproductive » (Cadre MINESEC/MINSANTÉ)]
\begin{itemize}
    \item \textbf{Club Santé / Club Sida / Club Jeunes :} Espace d'expression par les pairs, animations (théâtre-forum, causeries-éducatives, projections-débats), distribution de préservatifs (dépôt communautaire géré par PE formés), orientation vers SAJ.
    \item \textbf{Infirmerie scolaire / Personnel de santé scolaire :} Premier recours confidentiel, distribution préservatifs / CU, orientation vers SAJ/CSB/CEDA, gestion des urgences (viol, malaise), suivi grossesse (maintien scolarité circulaire MINESEC).
    \item \textbf{Enseignement SVTEEHB :} Temps curriculaire protégé (3h/semaine en 1ère C/TI), approche par compétences, situations concrètes camerounaises, évaluation des savoir-être (respect, consentement, responsabilité).
    \item \textbf{Liaison formelle École -- Structure de santé :} Convention CSB/SAJ -- Lycée : permanence infirmière/maïeuticien(ne) au lycée 1 fois/semaine, visite de classe au CSB/SAJ, numéros d'urgence affichés (Infirmerie, Vie Scolaire, 1510, 116, CSB local).
    \item \textbf{Implication parents / COPA :} Réunions d'information sur le programme SVTEEHB, déconstruction des rumeurs (« On apprend le sexe à l'école »), charte de non-violence / non-discrimination (grossesse, VIH, orientation).
    \item \textbf{Procédure claire violences / grossesse :} Signalement (Procureur, Délégué Régional MINESEC, MINAS), protection de la victime, maintien à l'école / réintégration post-partum (droit garanti par la loi).
\end{itemize}
\end{definition}

\begin{savaistu}[Numéros utiles au Cameroun -- À afficher en classe, dans l'infirmerie, au tableau]
\begin{center}
\begin{tabularx}{\linewidth}{|l|X|l|}\hline
\textbf{Service} & \textbf{Numéro / Contact} & \textbf{Disponibilité} \\ \hline
\textbf{NumVert Santé Reproductive (MINSANTÉ/UNFPA)} & \textbf{1510} & 24h/24, anonyme, gratuit \\ \hline
\textbf{NumVert Violences Basées sur le Genre (MINPROFF/UNFPA)} & \textbf{116} & 24h/24, anonyme, gratuit \\ \hline
\textbf{CAMNAFAW (Cliniques, info, orientation)} & 222 20 20 20 / 6XX XX XX XX & Heures ouvrées + urgences \\ \hline
\textbf{ACMS (Info produits, dépôts communautaires)} & 222 30 30 30 / *135\# & USSD, SMS \\ \hline
\textbf{RENATA (Réseau National des Associations de Tantines -- Filles mères)} & 6XX XX XX XX & Accueil, réinsertion, plaidoyer \\ \hline
\textbf{ALVF (Association de Lutte contre les Violences Faite aux Femmes)} & 222 XX XX XX / 6XX XX XX XX & Juridique, psycho, hébergement \\ \hline
\end{tabularx}
\end{center}
\emph{Nota : Les numéros mobiles changent ; vérifier l'annuaire actuel du MINSANTÉ ou du SAJ le plus proche.}
\end{savaistu}

\begin{aretenir}[Synthèse : Les 5 piliers d'une prévention efficace pour toi, jeune Camerounais(e) en Première]
\begin{enumerate}
    \item \textbf{Connais ton corps et tes droits :} Physiologie, cycle, consentement (libre, éclairé, réversible, enthousiaste), loi (accès soins 15 ans, avortement si vie en danger/viol/inceste).
    \item \textbf{Choisis ta stratégie ABCD :} Abstinence (report) OU Fidélité + Dépistage mutuel OU Préservatif systématique (Double protection = Préservatif + Implant/DIU/Injectable/Pilule). \textbf{Aucun jugement, juste ton choix éclairé.}
    \item \textbf{Identifie TON SAJ :} Où ? Horaires ? Numéro ? Y es-tu déjà allé(e) ? (Si non, va-y cette semaine, même sans problème, pour repérer les lieux et le personnel).
    \item \textbf{Maîtrise la check-list post-risque :} CU (J0--J3/5) $\rightarrow$ PPE VIH (J0--J3) / Traitement IST (si symptômes) $\rightarrow$ Contraception durable + Double protection $\rightarrow$ Dépistage VIH J42 $\rightarrow$ Soutien psycho si besoin. \textbf{NumVert 1510 en mémoire.}
    \item \textbf{Développe tes Life Skills :} Négociation DESC, refus assertif, recherche d'aide, pensée critique (médias, rumeurs, pression pairs), empathie, responsabilité partagée.
\end{enumerate}
\end{aretenir}

\begin{exercice}[Mise en situation intégrée -- Préparation au Probatoire]
\textbf{Contexte :} Lors d'une sortie de fin d'année, ton(ta) meilleur(e) ami(e) te confie : « J'ai couché avec X hier soir, on n'avait pas de capote, il/elle a dit qu'il/elle se retirait à temps. J'ai super peur d'être enceinte / d'avoir chopé un truc. Je n'ose pas aller à l'infirmerie, le surveillant général est copain avec mon père. »
\begin{enumerate}
    \item Identifie les \textbf{comportements à risque} présents dans cette situation (au moins 3).
    \item En utilisant la \textbf{check-list post-risque (6.5)}, détaille l'ordre chronologique des actions concrètes que tu lui conseilles de faire \textbf{aujourd'hui (J0)}, \textbf{demain (J1)}, \textbf{dans 6 semaines (J42)}. Précise les structures camerounaises accessibles (SAJ, CSB, CEDA, NumVert).
    \item Rédige un \textbf{script de négociation DESC} que ton ami(e) pourrait utiliser avec X pour exiger le préservatif ou refuser un nouveau rapport non protégé.
    \item Explique pourquoi la \textbf{double protection} est la norme recommandée pour lui/elle, et cite deux méthodes LARC adaptées aux nullipares.
    \item Quel \textbf{rôle l'école (infirmerie, club santé, SVTEEHB, liaison CSB)} peut-elle jouer pour que cette situation ne se reproduise pas et que ton ami(e) soit soutenu(e) sans risque de divulgation à ses parents ?
\end{enumerate}
\end{exercice}

\begin{corrige}[Exercice -- Mise en situation intégrée]
\begin{enumerate}
    \item \textbf{Comportements à risque :} (1) Rapport sexuel non protégé (pas de préservatif valide). (2) Recours au coït interrompu (retrait) comme unique « méthode » -- inefficace contraception + zéro protection IST/VIH. (3) Pression / manipulation du partenaire (« il/elle a dit... ») -- absence de négociation / consentement éclairé. (4) Peur d'accéder aux services (barrière confidentialité perçue). (5) Absence de contraception durable.
    \item \textbf{Plan d'action chronologique :}
    \begin{itemize}
        \item \textbf{J0 (Aujourd'hui -- URGENT) :} Se rendre au \textbf{SAJ le plus proche} (CSB, CEDA, CAMNAFAW, infirmerie si ouverte/confidentielle) pour : \textbf{Contraception d'urgence} (LNG 1,5 mg si < 72h, UPA 30 mg si < 120h, ou DIU Cu si < 120h -- \textbf{meilleure option}). Demander \textbf{PPE VIH} (Truvada 28j) si risque VIH évalué élevé (partenaire statut inconnu/positif non traité) -- délai max 72h, idéal < 4h. Dépistage \textbf{IST syndromique} si symptômes, sinon prévoir contrôle. Repartir avec \textbf{préservatifs gratuits} et RDV pour contraception durable.
        \item \textbf{J1 (Demain) :} Pose \textbf{contraception durable LARC} (Implant sous-cutané 3--5 ans OU DIU Cu 10 ans -- Catégorie OMS 1 pour nullipares) + Conseil \textbf{Double protection} (Préservatif systématique + Implant/DIU). Information sur \textbf{numéro vert 1510} pour questions/anxiété.
        \item \textbf{J42 (6 semaines) :} \textbf{Test VIH 4ème génération (Ag/Ac)} au CEDA ou SAJ (anonyme, gratuit, conseil pré/post-test). Si négatif = rassurant (>99\% fiable). Si positif = prise en charge ARV immédiate (Test \& Treat). Dépistage IST de contrôle si traitement initial.
    \end{itemize}
    \item \textbf{Script DESC :} « \textbf{D :} Hier, on n'avait pas de préservatif valide et tu t'es retiré. Ce n'est pas une protection. \textbf{E :} J'ai très peur d'être enceinte et d'avoir attrapé une maladie. Je ne dors plus. \textbf{S :} Pour la prochaine fois, on achète des préservatifs \textbf{avant} (kiosque, pharmacie, SAJ gratuit). Moi, je vais me faire poser l'implant/DIU cette semaine. Si on n'a pas de capote, \textbf{pas de pénétration}. On peut faire autre chose. \textbf{C :} Comme ça, on profite à fond, sans stress, en se respectant. » $\rightarrow$ Si X refuse / se moque / menace : « Non, c'est non. Ma santé passe avant. »
    \item \textbf{Double protection :} Norme car : (1) Fertilité max à cet âge, coût grossesse énorme (scolarité, santé maternelle). (2) Risque IST/VIH réel (partenaires multiples, biologie). (3) Observance imparfaite possible (oubli, rupture). \textbf{LARC nullipares : Implant (Nexplanon\textsuperscript{\textregistered}/Jadelle\textsuperscript{\textregistered}) et DIU Cu (TT380A/MLCu380).} Efficacité \textasciitilde 99,9\%, oubli impossible, réversibles, discrets, gratuits en SAJ/CSB.
    \item \textbf{Rôle école :} \textbf{Infirmerie :} Accueil confidentiel J0 (CU, orientation, préservatifs), sans informer parents (secret professionnel, circulaire). \textbf{Club Santé / PE :} Soutien pair, accompagnement au SAJ, distribution capotes, déstigmatisation. \textbf{SVTEEHB :} Cours sur CU, PPE, double protection, consentement, numéros utiles -- connaissances pour agir. \textbf{Liaison CSB/SAJ :} Permanence infirmière/maïeuticien(ne) au lycée, visite de classe au CSB pour lever la peur de l'inconnu. \textbf{Vie Scolaire / Proviseur :} Application circulaire maintien scolarité, protection contre discrimination, procédure signalement violences si contrainte/viol avéré.
\end{enumerate}
\end{corrige}

\newpage

\coursec{Activité d'intégration}

\courssub{Objectifs de l'activité}

Cette activité d'intégration te place dans une situation réaliste et fréquente de la vie d'un lycéen camerounais. À travers l'étude du cas « Le dilemme de Chantal et Kevin », tu vas mobiliser l'ensemble des savoirs et savoir-faire de la leçon. À la fin de cette activité, tu dois être capable de :

\begin{itemize}
\item \textbf{identifier} tous les comportements à risque pour la santé reproductive présents dans une situation donnée (savoir-faire 1) ;
\item \textbf{expliquer} précisément les conséquences physiologiques, sanitaires et sociales de chaque comportement identifié (savoir-faire 2) ;
\item \textbf{proposer} des alternatives responsables et concrètes, adaptées au contexte camerounais (savoir-faire 3) ;
\item \textbf{choisir et justifier} une méthode de contraception adaptée à un profil donné, en argumentant sur l'efficacité, la protection contre les IST, l'observance et l'accessibilité ;
\item \textbf{t'orienter} vers les structures de santé réelles du Cameroun (centres de santé intégrés, hôpitaux, centres de dépistage, CAMNAFAW, etc.).
\end{itemize}

L'activité se déroule en groupes de 4 à 5 élèves, avec restitution orale et production d'une affiche de synthèse. Elle est évaluée à l'aide d'une grille critériée portant sur les savoirs, les savoir-faire et les attitudes (écoute, respect, non-jugement, travail en équipe).

\courssub{Situation}

\textbf{Le dilemme de Chantal et Kevin}

Chantal, 16 ans, est élève en classe de Première C au lycée de Nkongsamba. Kevin, 19 ans, est en classe de Terminale dans le même établissement. Ils sortent ensemble depuis 3 mois.

Depuis quelques semaines, Kevin insiste de plus en plus pour qu'ils aient des rapports sexuels, affirmant que « c'est la preuve d'amour ». Chantal hésite : elle a peur de tomber enceinte — sa cousine a abandonné ses études après une grossesse à 17 ans — et elle a entendu parler des infections sexuellement transmissibles (IST), notamment du VIH/Sida, lors d'une campagne de sensibilisation au centre de santé de quartier.

Pour la rassurer, Kevin propose deux « solutions » qu'il dit avoir apprises de ses amis : le \textbf{retrait} (« je me retirerai avant ») et la \textbf{« période sûre »} (« il y a des jours où on ne peut pas tomber enceinte »). De son côté, Amina, la meilleure amie de Chantal, lui conseille de garder sous la main une \textbf{pilule du lendemain} « au cas où », qu'on peut acheter sans ordonnance dans certaines pharmacies de la ville.

Un samedi soir, lors d'une fête d'anniversaire dans le quartier, Chantal et Kevin consomment tous deux des boissons alcoolisées. Sous l'emprise de l'alcool, ils ont un \textbf{rapport sexuel non protégé}, sans préservatif et sans aucune contraception. Le rapport a lieu 4 jours après la fin des règles de Chantal, dont le cycle menstruel dure environ 28 jours.

Le lendemain matin, Chantal panique. Il est 8 h, soit environ 12 heures après le rapport. Elle se pose mille questions : puis-je être enceinte ? Ai-je pu attraper une IST ? Que faire, où aller, et comment éviter que cela se reproduise ?

\textbf{Données utiles :} cycle de Chantal : 28 jours ; ovulation estimée vers le 14\textsuperscript{e} jour du cycle ; durée de survie des spermatozoïdes dans les voies génitales féminines : jusqu'à 5 jours ; durée de vie de l'ovule après l'ovulation : 12 à 24 heures ; délai maximal d'efficacité de la contraception d'urgence (lévonorgestrel) : 72 h (idéalement le plus tôt possible) ; fenêtre de dépistage VIH fiable : à partir de 6 semaines après la prise de risque (test rapide de 4\textsuperscript{e} génération) ; taux d'échec du retrait en usage courant : environ 20 à 22 grossesses pour 100 femmes sur un an ; préservatif masculin bien utilisé : environ 98 \% d'efficacité contraceptive et seule méthode protégeant aussi des IST.

\courssub{Consignes}

Travail en groupes de 4 à 5 élèves. Durée : 2 h (préparation) + restitution orale (10 min par groupe) avec affiche de synthèse.

\begin{enumerate}
\item \textbf{Identifier} TOUS les comportements à risque présents dans cette histoire (au minimum 5). Pour chacun, citer la phrase ou l'élément de la situation qui le révèle.
\item Pour \textbf{chaque} comportement à risque identifié, \textbf{expliquer} la conséquence physiologique, sanitaire ou sociale précise qu'il peut entraîner, en t'appuyant sur les données numériques de la situation et sur les rappels de physiologie (cycle menstruel, fécondation).
\item \textbf{Proposer} 3 alternatives responsables concrètes que Chantal aurait pu activer : l'une \textbf{AVANT} le rapport, l'une \textbf{PENDANT} (au moment de la relation), l'une \textbf{APRÈS} le rapport.
\item Si Chantal devient sexuellement active de façon consentie et choisie, \textbf{choisir} la méthode contraceptive la plus adaptée à son cas et la \textbf{justifier} par 3 arguments portant sur : l'efficacité, la protection contre les IST, l'observance (facilité d'utilisation régulière) et l'accès (coût, disponibilité au Cameroun).
\item \textbf{Indiquer} la structure exacte (nom et type) où Chantal peut se rendre dès le lendemain matin au Cameroun pour : (a) obtenir une contraception d'urgence ; (b) effectuer un dépistage du VIH ; (c) se faire poser un implant contraceptif ; (d) recevoir un conseil psychologique et un accompagnement.
\end{enumerate}

\begin{attention}[Règles de travail en groupe]
Cette activité aborde des sujets sensibles. Les règles sont : respect de la parole de chacun, interdiction de tout jugement moral ou moquerie, vocabulaire scientifique exigé (on parle d'ovulation, de fécondation, d'IST — pas de rumeurs), confidentialité : ce qui se dit dans le groupe reste dans le groupe.
\end{attention}

\courssub{Ressources à mobiliser}

Pour réussir cette activité, tu dois mobiliser les savoirs de la leçon :

\begin{itemize}
\item \textbf{Physiologie de la reproduction :} le cycle menstruel (environ 28 jours) comporte une phase folliculaire, une ovulation (vers le 14\textsuperscript{e} jour) et une phase lutéale ; la fécondation n'est possible que pendant une fenêtre fertile d'environ 6 jours (5 jours avant l'ovulation + le jour de l'ovulation), car les spermatozoïdes survivent jusqu'à 5 jours et l'ovule 12 à 24 h.
\item \textbf{Comportements à risque :} rapports sexuels non protégés, pression psychologique et rapports non pleinement consentis, consommation d'alcool (qui altère le jugement), recours à des méthodes inefficaces (retrait, « période sûre »), mésusage de la contraception d'urgence.
\item \textbf{Conséquences :} grossesse précoce non désirée (risques obstétricaux chez l'adolescente, interruption des études, rejet familial et social), IST dont le VIH/Sida, avortement clandestin et ses complications (hémorragies, infections, stérilité, décès).
\item \textbf{Moyens de contraception :} préservatif masculin et féminin (double protection : grossesse + IST), pilule combinée ou microprogestative (efficace mais exige une observance quotidienne et ne protège pas des IST), injectables, implant (efficacité supérieure à 99 \% pendant 3 à 5 ans, faible contrainte d'observance), stérilet, contraception d'urgence (à réserver aux situations exceptionnelles, efficace surtout dans les 24 premières heures).
\end{itemize}

\courssub{Démarche et corrigé commenté}

\begin{exoresolu}[Corrigé commenté de l'étude de cas « Le dilemme de Chantal et Kevin »]
\textbf{Rappel des questions :}
\begin{enumerate}
\item Identifier tous les comportements à risque (minimum 5) en citant les éléments du texte.
\item Expliquer les conséquences précises de chaque risque en utilisant les données physiologiques.
\item Proposer 3 alternatives responsables : AVANT, PENDANT, APRÈS.
\item Choisir et justifier la méthode contraceptive la plus adaptée (3 arguments).
\item Indiquer les structures de santé camerounaises exactes pour : CU, dépistage VIH, pose d'implant, soutien psychologique.
\end{enumerate}

\tcblower

\textbf{Question 1 — Identification des comportements à risque.}

On relève au moins 6 comportements à risque :
\begin{enumerate}
\item \textbf{La pression psychologique exercée par Kevin} (« c'est la preuve d'amour ») : un rapport obtenu sous pression n'est pas un choix libre et éclairé ; c'est un facteur de risque en soi (atteinte au consentement).
\item \textbf{Le recours envisagé au retrait} (coït interrompu) : méthode très peu fiable (environ 20 à 22 \% d'échec en usage courant), car le liquide pré-éjaculatoire peut contenir des spermatozoïdes et le contrôle de l'éjaculation est aléatoire.
\item \textbf{La croyance en une « période sûre »} (méthode Ogino-Knaus) : l'ovulation peut survenir plus tôt ou plus tard que prévu, surtout chez une adolescente dont les cycles sont souvent irréguliers ; il n'existe pas de période totalement inféconde.
\item \textbf{Le mésusage envisagé de la pilule du lendemain} « au cas où » : la contraception d'urgence n'est pas une méthode de contraception régulière ; son efficacité diminue avec le temps et elle ne protège pas des IST.
\item \textbf{La consommation d'alcool} : elle altère le jugement, diminue la capacité à négocier le préservatif et favorise les rapports non protégés.
\item \textbf{Le rapport sexuel non protégé} (sans préservatif ni contraception) : c'est le comportement à risque majeur, cumulant risque de grossesse et risque d'IST.
\end{enumerate}

\textbf{Question 2 — Conséquences précises de chaque risque.}

\begin{itemize}
\item \textbf{Pression psychologique :} rapport non pleinement consenti $\Rightarrow$ souffrance psychologique, perte d'estime de soi, et exposition à des rapports non désirés répétés.
\item \textbf{Retrait :} échec fréquent $\Rightarrow$ grossesse non désirée ; aucune protection contre les IST (le contact des muqueuses suffit à transmettre le VIH, la chlamydiose, la gonococcie, la syphilis...).
\item \textbf{« Période sûre » :} raisonnons avec les données du cycle de Chantal.
  \begin{itemize}
  \item Jour 1 = 1\textsuperscript{er} jour des règles.
  \item Fin des règles $\approx$ Jour 5 (durée moyenne 5 jours).
  \item Rapport = 4 jours après la fin des règles = \textbf{Jour 9 du cycle}.
  \item Ovulation estimée = \textbf{Jour 14}.
  \item Survie spermatozoïdes = \textbf{5 jours maximum}.
  \end{itemize}
  Des spermatozoïdes déposés au Jour 9 peuvent survivre jusqu'au Jour 14 (9 + 5 = 14), jour de l'ovulation. \textbf{La fécondation est donc physiologiquement possible.} La « période sûre » de Kevin n'en est pas une.
  
  \begin{popfigure}
  \begin{tikzpicture}[scale=0.9, every node/.style={font=\small}]
    % Axe principal
    \draw[thick, ->] (0,0) -- (16,0) node[right] {Jours du cycle};
    \foreach \x/\label in {1/1, 5/5, 9/9, 14/14, 28/28}
      \draw (\x,0.1) -- (\x,-0.1) node[below] {\label};
    
    % Règles
    \draw[popred, line width=3pt] (1,0.8) -- (5,0.8) node[midway, above=2pt] {Règles (J1--J5)};
    
    % Rapport
    \draw[dashed, poporange, thick] (9,-0.5) -- (9,1.5);
    \node[poporange, below] at (9,-0.7) {Rapport (J9)};
    
    % Survie spermatozoïdes
    \draw[popblue, thick, <->] (9,1.5) -- (14,1.5) node[midway, above=2pt] {Survie spermatozoïdes (5 j max)};
    
    % Ovulation
    \draw[popgreen!70!black, thick] (14,-0.5) -- (14,1.5);
    \node[popgreen!70!black, above] at (14,1.7) {Ovulation (J14)};
    
    % Fenêtre fertile à risque
    \fill[poporange!20] (9,-0.3) rectangle (14.5,0.3);
    \node[poporange!80!black] at (11.5,0.5) {Fenêtre fertile à risque};
    
    % Phase lutéale
    \draw[poppurple, line width=2pt] (15,0.8) -- (28,0.8) node[midway, above=2pt] {Phase lutéale};
  \end{tikzpicture}
  \legende{Chronologie du cycle de Chantal (28 jours) montrant que le rapport au Jour 9 tombe dans la fenêtre de survie des spermatozoïdes jusqu'à l'ovulation au Jour 14. Le risque de fécondation est réel.}
  \end{popfigure}

\item \textbf{Mésusage de la contraception d'urgence :} efficacité décroissante (environ 95 \% dans les 24 h, moins de 60 \% à 72 h pour le lévonorgestrel), effets secondaires (nausées, troubles du cycle), fausse sécurité, et aucune protection contre les IST.
\item \textbf{Alcool :} baisse du contrôle de soi $\Rightarrow$ oubli ou refus du préservatif, rapports non désirés, exposition cumulée grossesse + IST.
\item \textbf{Rapport non protégé :} 
  \begin{itemize}
  \item \emph{Sanitaire :} risque de grossesse (fenêtre fertile démontrée ci-dessus) et de contamination par le VIH ou une autre IST (chlamydiose, gonococcie, syphilis, hépatite B, HPV...).
  \item \emph{Social :} en cas de grossesse, risque d'interruption des études (comme la cousine de Chantal), de rejet familial, de stigmatisation, de mariage forcé ou de recours à un avortement clandestin aux complications graves (hémorragie, infection, stérilité, voire décès maternel).
  \end{itemize}
\end{itemize}

\textbf{Question 3 — Trois alternatives responsables.}

\begin{itemize}
\item \textbf{AVANT :} Chantal peut affirmer son droit de dire non et fixer ses limites (consentement) ; si elle choisit librement d'être sexuellement active, elle doit consulter \textbf{au préalable} un centre de santé (CSI, hôpital de district) ou une structure \textbf{CAMNAFAW} pour choisir une contraception adaptée et se procurer des préservatifs. Elle peut en parler à une personne de confiance (infirmière scolaire, éducatrice, personnel de santé).
\item \textbf{PENDANT :} exiger l'utilisation d'un \textbf{préservatif masculin ou féminin} à chaque rapport, correctement mis dès le début et jusqu'au bout ; refuser tout rapport sous l'emprise de l'alcool ou de substances ; s'assurer du consentement mutuel explicite.
\item \textbf{APRÈS :} se rendre \textbf{le plus tôt possible} (idéalement dans les 12--24 h, au plus tard dans les 72 h) dans une pharmacie ou un centre de santé pour une \textbf{contraception d'urgence} (lévonorgestrel 1,5 mg en dose unique) ; effectuer un \textbf{dépistage du VIH} (test rapide TROD, à refaire après la fenêtre de 6 semaines pour lever le doute) et un dépistage des autres IST ; puis choisir une contraception régulière efficace.
\end{itemize}

\textbf{Question 4 — Méthode la plus adaptée et justification.}

Le choix recommandé pour une adolescente sexuellement active est la \textbf{double protection : préservatif (masculin ou féminin) à chaque rapport + implant contraceptif sous-cutané} (ou, à défaut, injectable progestatif trimestriel). Justification :
\begin{enumerate}
\item \textbf{Efficacité contraceptive maximale :} l'implant (étinogestrel ou lévonorgestrel) dépasse 99 \% d'efficacité (indice de Pearl < 0,1) pendant 3 à 5 ans, sans dépendre d'un geste quotidien ni de la mémoire ; le préservatif sécurise chaque rapport en cas de défaillance exceptionnelle.
\item \textbf{Protection contre les IST (argument décisif) :} seul le préservatif (masculin ou féminin) protège simultanément des grossesses non désirées \textbf{et} des IST, dont le VIH. C'est indispensable pour une jeune fille dont le partenaire n'a pas fait de dépistage récent et dont la fidélité n'est pas garantie.
\item \textbf{Observance et accès au Cameroun :} l'implant ne demande aucune observance quotidienne (contrairement à la pilule, dont l'oubli est fréquent à cet âge et source d'échecs) ; il est disponible à coût subventionné (souvent gratuit ou < 1000 FCFA) dans les services de planification familiale des centres de santé intégrés (CSI), hôpitaux de district et cliniques \textbf{CAMNAFAW} présentes dans les grandes villes (Yaoundé, Douala, Nkongsamba, Bafoussam, Garoua, Maroua, etc.).
\end{enumerate}

\textbf{Question 5 — Où aller dès demain matin au Cameroun.}

\begin{itemize}
\item \textbf{(a) Contraception d'urgence :} une \textbf{pharmacie} agréée (délivrance possible sans ordonnance, coût modique $\approx$ 500--1000 FCFA) ou le \textbf{centre de santé intégré (CSI)} le plus proche — Chantal doit y aller dès ce matin (il reste $\approx$ 60 h de délai utile, l'efficacité est maximale dans les 12--24 h).
\item \textbf{(b) Dépistage VIH :} un \textbf{Centre de Conseil et de Dépistage (CCD)} ou \textbf{Centre de Dépistage Volontaire (CDV)} d'un hôpital de district ou d'un CSI — dépistage gratuit, confidentiel, avec pré et post-conseil ; premier test immédiat (TROD 4\textsuperscript{e} génération), test de confirmation après 6 semaines (fenêtre sérologique).
\item \textbf{(c) Pose d'un implant :} le service de \textbf{planification familiale (PF)} d'un CSI, d'un hôpital de district (ex: Hôpital de District de Nkongsamba, Hôpital Laquintinie de Douala, Hôpital Central de Yaoundé), ou une clinique \textbf{CAMNAFAW} (Association Camerounaise pour le Bien-Être Familial). La pose est réalisée par un personnel qualifié (sage-femme, infirmier PF, médecin) après conseil et consentement.
\item \textbf{(d) Conseil psychologique et accompagnement :} la \textbf{cellule d'écoute / service social} du CSI ou de l'hôpital ; le \textbf{centre d'écoute CAMNAFAW} ; une ONG spécialisée dans la santé des adolescents (ex: \emph{Association pour la Promotion de la Santé de l'Adolescent - APSA}, \emph{RENATA} pour les grossesses précoces) ; en cas de rapport forcé ou sous contrainte (viol), dépôt de plainte possible auprès de la \textbf{Brigade des Mœurs} (commissariat) et prise en charge médico-légale urgente (certificat médical, prophylaxie post-exposition VIH si < 72 h).
\end{itemize}
\end{exoresolu}

\courssub{Bilan de l'activité}

Cette étude de cas a permis de mettre en évidence plusieurs acquis essentiels :

\begin{itemize}
\item \textbf{Les risques sont cumulatifs et interconnectés :} dans l'histoire de Chantal, l'alcool a levé les inhibitions, ce qui a conduit au rapport non protégé, lequel expose simultanément à la grossesse et aux IST. Un comportement à risque en entraîne presque toujours un autre.
\item \textbf{Les « fausses bonnes idées » sont dangereuses :} le calcul de la fenêtre fertile a démontré scientifiquement que la « période sûre » invoquée par Kevin n'existe pas (J9 + 5 j = J14 = ovulation). Le retrait et le mésusage de la pilule du lendemain relèvent du même constat : seules les données physiologiques permettent de démasquer ces idées reçues.
\item \textbf{La prévention repose sur trois temps :} anticiper (information, contraception choisie, préservatifs disponibles), se protéger au moment du rapport (préservatif systématique), et réagir vite après un accident (contraception d'urgence dans les 72 h, dépistage).
\item \textbf{Le consentement est non négociable :} « c'est la preuve d'amour » est un argument de manipulation ; dire non est un droit, et un rapport sous pression ou sous alcool n'est pas un choix responsable.
\item \textbf{Le système de santé camerounais offre des solutions concrètes et accessibles :} pharmacies, centres de santé intégrés (CSI), centres de conseil et de dépistage (CCD/CDV), services de planification familiale et CAMNAFAW. Savoir où aller, c'est déjà se protéger.
\end{itemize}

\begin{aretenir}[L'essentiel de l'activité]
Face à une situation à risque pour la santé reproductive, la démarche responsable comporte quatre réflexes : \cle{identifier} les comportements dangereux, \cle{expliquer} leurs conséquences à l'aide des connaissances physiologiques, \cle{agir} avant, pendant et après (information, préservatif, contraception d'urgence, dépistage), et \cle{s'orienter} vers les structures de santé compétentes. La double protection (préservatif + contraception efficace longue durée type implant) reste la stratégie la plus sûre pour un jeune sexuellement actif.
\end{aretenir}

\newpage

\coursec{Méthodes et exercices résolus}

\coursec{Lutte contre les comportements préjudiciables à la santé reproductive}

\courssub{Rappels physiologiques : les bases de la fertilité}

\begin{definition}[Appareil reproducteur féminin et cycles]
L'appareil reproducteur féminin comprend les \textbf{ovaires} (gonades : production d'ovocytes et d'hormones), les \textbf{trompes de Fallope} (site de la fécondation), l'\textbf{utérus} (cavité utérine bordée d'endomètre, site de l'implantation), le \textbf{col utérin} et le \textbf{vagin}. La fertilité repose sur deux cycles synchronisés d'environ 28 jours (variable de 21 à 35 jours) :
\begin{itemize}
\item Le \textbf{cycle ovarien} : croissance folliculaire (phase folliculaire) $\rightarrow$ \textbf{ovulation} (libération de l'ovocyte II) $\rightarrow$ formation du \textbf{corps jaune} (phase lutéale).
\item Le \textbf{cycle utérin} (endomètre) : \textbf{menstruations} (jours 1--5) $\rightarrow$ \textbf{phase proliférative} (épaississement sous œstrogènes) $\rightarrow$ \textbf{phase sécrétoire} (vascularisation, glandes sous progestérone) propice à la nidation.
\end{itemize}
\end{definition}

\begin{propriete}[Fenêtre de fertilité et fécondation]
L'ovulation survient vers le \textbf{14\textsuperscript{e} jour} (pic de LH). L'ovocyte survit \textbf{12 à 24 h}. Les spermatozoïdes survivent \textbf{3 à 5 jours} dans les voies génitales féminines. La \textbf{période fertile} s'étend donc du \textbf{4--5 jours avant l'ovulation au lendemain de l'ovulation}. La fécondation (rencontre spermatozoïde/ovocyte) a lieu dans l'ampoule de la trompe. Le zygote migre vers l'utérus (3--4 jours) et s'implante (nidation) vers le \textbf{6--7\textsuperscript{e} jour} post-fécondation.
\end{propriete}

\begin{popfigure}
\begin{tikzpicture}[scale=0.85, every node/.style={font=\small}]
% Axes
\draw[->] (0,0) -- (12,0) node[right] {Jours du cycle};
\draw[->] (0,0) -- (0,5.5) node[above] {Niveau hormonal (arb. u.)};
\foreach \x in {0,7,14,21,28} \draw (\x/2.33,0.1) -- (\x/2.33,-0.1) node[below] {\x};
% FSH
\draw[popblue, thick, domain=0:12, samples=100] plot (\x/2.33, {2.5*exp(-(\x-3)^2/20) + 0.5});
\node[popblue] at (1.5, 3.5) {FSH};
% LH
\draw[poporange, thick, domain=0:12, samples=100] plot (\x/2.33, {4*exp(-(\x-14)^2/8) + 0.3});
\node[poporange] at (6, 4.5) {LH};
% Oestradiol (E2)
\draw[poppink, thick, domain=0:12, samples=100] plot (\x/2.33, {3*exp(-(\x-13)^2/18) + 0.5});
\draw[poppink, thick, domain=14:28, samples=100] plot (\x/2.33, {2*exp(-(\x-21)^2/18) + 0.5});
\node[poppink] at (3, 3.8) {Œstradiol};
% Progestérone (P4)
\draw[popgreen, thick, domain=14:28, samples=100] plot (\x/2.33, {3.5*exp(-(\x-21)^2/18) + 0.3});
\node[popgreen] at (9, 4) {Progestérone};
% Phases
\node[align=center] at (1.5, -0.7) {Menstruations\\(J1--J5)};
\node[align=center] at (4.5, -0.7) {Phase folliculaire};
\node[align=center] at (7.5, -0.7) {Ovulation\\J14};
\node[align=center] at (10.5, -0.7) {Phase lutéale};
\end{tikzpicture}
\legende{Profils hormonaux au cours d'un cycle ovarien type de 28 jours. Le pic de LH déclenche l'ovulation (J14). L'œstradiol domine la phase folliculaire, la progestérone la phase lutéale.}
\end{popfigure}

\begin{experience}[Observation de la glaire cervicale : signe d'ovulation]
\textbf{Matériel :} Échantillons de glaire (ou images microscopiques) prélevés à J8, J14, J22. \textbf{Protocole :} Étaler sur lame, observer au microscope (×100). \textbf{Observations :} J8 : glaire épaisse, filante courte, mailles serrées (barrière). J14 : glaire fluide, très filante (\emph{filant} > 10 cm), mailles larges en \emph{feuilles de fougère} (arborisation) -- facilite le passage des spermatozoïdes. J22 : glaire épaisse, pâteuse, non filante (bouchon). \textbf{Interprétation :} Sous l'effet des œstrogènes (pic pré-ovulatoire), la glaire devient perméable ; sous progestérone, elle redevient imperméable. C'est un indicateur naturel de fertilité (méthode de Billings), mais insuffisant seul pour contraception.
\end{experience}

\begin{aretenir}[Repères physiologiques clés pour la prévention]
\begin{itemize}
\item \textbf{Puberté :} premières règles (ménarche) $\neq$ maturité pelvienne complète (bassin finit de croître vers 18--20 ans).
\item \textbf{Ovulation :} imprévisible chez l'adolescente (cycles anovulatoires ou irréguliers fréquents les 2 premières années).
\item \textbf{Rapport non protégé :} risque de grossesse \textbf{dès le premier rapport} si ovulation ; risque d'IST \textbf{à chaque rapport}.
\item \textbf{Double protection :} seul le \cle{préservatif} (masculin ou féminin) protège simultanément contre la grossesse non désirée et les IST (VIH, HPV, blennorragie, chlamydiose, syphilis...).
\end{itemize}
\end{aretenir}

\courssub{Comportements à risque : identification et déterminants}

\begin{methode}[Identifier un comportement à risque dans une situation donnée]
\begin{enumerate}
\item \textbf{Lire attentivement la situation} et repérer les faits : âge des personnes, nature des rapports (protégés ou non), fréquence, contexte (pression, ignorance, alcool, transaction).
\item \textbf{Confronter chaque fait aux critères de risque} : un comportement est à risque s'il expose à une grossesse non désirée, à une infection sexuellement transmissible (IST) dont le VIH/Sida, ou à une atteinte de l'intégrité physique ou psychique.
\item \textbf{Nommer précisément le comportement} avec le vocabulaire scientifique : \emph{rapport sexuel non protégé}, \emph{partenaires multiples}, \emph{rapport précoce avant maturité du bassin}, \emph{rapport transactionnel}, \emph{recours à l'avortement clandestin}.
\item \textbf{Justifier l'identification} en reliant le fait observé au risque encouru : « le rapport est non protégé (fait), donc il expose à une grossesse non désirée et aux IST (risque) ».
\item \textbf{Ne pas conclure sans preuve} : un seul élément du texte suffit à identifier, mais il faut toujours le citer ou le reformuler explicitement.
\end{enumerate}
\textbf{Réflexe Probatoire} : la réponse type est « Fait observé + nom du comportement + risque associé ». Une identification sans justification ne vaut que la moitié des points.
\end{methode}

\begin{exoresolu}[Identifier les comportements à risque dans un récit de vie]
\textbf{Énoncé.} Dans un quartier de Douala, Amina, 16 ans, élève en classe de Seconde, entretient depuis six mois une relation avec un commerçant de 28 ans qui lui offre de l'argent et des cadeaux. Ils ont des rapports sexuels sans préservatif ; Amina pense que « se laver après le rapport » suffit à éviter une grossesse. Sa camarade Brigitte, 17 ans, a déjà interrompu deux grossesses chez une « matrone » du quartier, à l'aide de tisanes et d'un manioc introduit dans le vagin.
\begin{enumerate}
\item Relever dans ce texte trois comportements à risque pour la santé reproductive.
\item Pour chacun, préciser le ou les risques encourus.
\end{enumerate}
\tcblower
\textbf{Solution détaillée.}

\textbf{1. Identification des comportements à risque.}

\textbf{Comportement 1 : les rapports sexuels non protégés d'Amina.} Le texte indique explicitement qu'ils ont des rapports « sans préservatif ». Or le préservatif est le seul moyen qui protège à la fois des grossesses non désirées et des IST. De plus, le partenaire est un adulte dont on ignore le statut sérologique, et la relation est \emph{transactionnelle} (argent et cadeaux), ce qui réduit le pouvoir de négociation de la jeune fille.

\textbf{Comportement 2 : la croyance en une méthode inefficace (lavage après le rapport).} Le lavage vaginal après le rapport ne détruit ni les spermatozoïdes déjà présents dans l'utérus (ils peuvent y parvenir en quelques minutes) ni les agents infectieux. Cette fausse croyance entretient la prise de risque.

\textbf{Comportement 3 : les avortements clandestins de Brigitte.} Deux interruptions de grossesse réalisées hors structure sanitaire, par une personne non qualifiée, avec des moyens non médicaux (tisanes, corps étranger vaginal). C'est un avortement clandestin, comportement à très haut risque.

\textbf{2. Risques encourus.}

\begin{center}
\begin{tabularx}{\linewidth}{|l|X|}
\hline
\rowcolor{popblueL} \textbf{Comportement} & \textbf{Risques encourus} \\
\hline
Rapports non protégés (Amina) & Grossesse précoce non désirée ; IST dont VIH/Sida, blennorragie, chlamydiose ; déscolarisation si grossesse \\
\hline
Croyance au lavage & Maintien de l'exposition au risque (faux sentiment de protection) ; irritations et infections vaginales favorisées par les douches répétées \\
\hline
Avortements clandestins (Brigitte) & Hémorragie, infection (septicémie), perforation utérine, stérilité définitive, décès ; séquelles psychologiques ; risque pénal \\
\hline
\end{tabularx}
\end{center}

\textbf{Conclusion.} Trois comportements à risque sont identifiables : les rapports non protégés, la croyance en une méthode de « protection » inefficace et les avortements clandestins. Chacun expose les jeunes filles à des conséquences sanitaires graves (grossesse précoce, IST, complications d'avortement) et sociales (déscolarisation, stigmatisation).
\end{exoresolu}

\courssub{Grossesses précoces : conséquences sanitaires et sociales}

\begin{methode}[Expliquer les conséquences sanitaires et sociales]
\begin{enumerate}
\item \textbf{Rappeler le mécanisme biologique} qui relie le comportement à la conséquence : fécondation possible dès l'ovulation (rapport non protégé $\rightarrow$ grossesse), fragilité anatomique d'un bassin immature (grossesse précoce $\rightarrow$ dystocie), geste septique en milieu non stérile (avortement clandestin $\rightarrow$ infection).
\item \textbf{Distinguer les plans} : conséquences \emph{sanitaires} (sur le corps et la santé), \emph{sociales} (scolarité, famille, insertion), \emph{psychologiques} (culpabilité, dépression), parfois \emph{juridiques}.
\item \textbf{Hiérarchiser} : commencer par les conséquences immédiates (hémorragie, infection) puis les conséquences à long terme (stérilité, fistule obstétricale, déscolarisation).
\item \textbf{Appuyer sur des données} quand elles sont fournies : citer les chiffres du document et les exploiter (calcul de pourcentage, comparaison).
\item \textbf{Rédiger en phrases complètes} avec des connecteurs logiques : « parce que », « ce qui entraîne », « par conséquent ».
\end{enumerate}
\textbf{Formule à retenir} : pourcentage $= \dfrac{\text{effectif de la partie}}{\text{effectif total}} \times 100$.
\end{methode}

\begin{exoresolu}[Exploiter des données sur les grossesses précoces]
\textbf{Énoncé.} Dans un hôpital de district de la région de l'Extrême-Nord du Cameroun, on a enregistré en un an 800 accouchements. Parmi eux, 240 concernaient des mères âgées de 13 à 17 ans. Chez ces jeunes mères, on a observé 96 cas de complications majeures (dystocie, hémorragie, fistule obstétricale), contre 80 cas chez les 560 mères de 20 à 34 ans.
\begin{enumerate}
\item Calculer la proportion de grossesses précoces parmi les accouchements.
\item Comparer la fréquence des complications chez les jeunes mères et chez les mères de 20 à 34 ans.
\item Expliquer, par un mécanisme physiologique, pourquoi les grossesses précoces sont plus à risque.
\end{enumerate}
\tcblower
\textbf{Solution détaillée.}

\textbf{1. Proportion de grossesses précoces.}
\begin{align*}
p &= \frac{240}{800} \times 100 = 30\,\%.
\end{align*}
30 \% des accouchements concernent des adolescentes de 13 à 17 ans : c'est une proportion très élevée qui traduit l'ampleur des grossesses précoces.

\textbf{2. Comparaison des fréquences de complications.}

Chez les jeunes mères (13--17 ans) :
\[ f_1 = \frac{96}{240} \times 100 = 40\,\%. \]

Chez les mères de 20 à 34 ans :
\[ f_2 = \frac{80}{560} \times 100 \approx 14,3\,\%. \]

Rapport : $\dfrac{40}{14,3} \approx 2,8$. Les complications majeures sont donc près de \textbf{trois fois plus fréquentes} chez les adolescentes que chez les femmes adultes.

\textbf{3. Explication physiologique.}

Chez l'adolescente, la croissance du \textbf{bassin osseux} n'est pas achevée : le détroit pelvien est souvent trop étroit pour le passage de la tête fœtale, ce qui provoque des \textbf{dystocies} (accouchements difficiles) avec travail prolongé. Un travail prolongé comprime les tissus mous (vessie, vagin) contre les os du bassin ; privés de sang, ces tissus meurent et laissent une communication anormale : c'est la \textbf{fistule obstétricale}, source d'incontinence urinaire et/ou fécale permanente et d'exclusion sociale. De plus, l'organisme de l'adolescente, encore en croissance, entre en compétition nutritionnelle avec le fœtus (anémie maternelle, retard de croissance intra-utérin), et l'immaturité utérine favorise les hémorragies du post-partum (atonie utérine).

\textbf{Conclusion.} Avec 30 \% de grossesses précoces et une fréquence de complications presque triplée (40 \% contre 14,3 \%), les données confirment que la grossesse précoce est un problème majeur de santé publique, expliqué notamment par l'immaturité anatomique du bassin de l'adolescente.
\end{exoresolu}

\courssub{Avortements clandestins : un drame de santé publique}

\begin{methode}[Traiter un sujet sur l'avortement clandestin]
\begin{enumerate}
\item \textbf{Définir} : interruption volontaire de grossesse pratiquée hors cadre médical légal, par du personnel non qualifié ou avec des moyens non médicaux (tisanes, corps étrangers, médicaments détournés sans surveillance).
\item \textbf{Expliquer la chaîne des complications} : geste non aseptique $\rightarrow$ infection (métrite, septicémie) ; traumatisme instrumental $\rightarrow$ perforation utérine, hémorragie ; évacuation incomplète $\rightarrow$ rétention placentaire, infection secondaire ; cicatrisation pathologique $\rightarrow$ synéchies, stérilité.
\item \textbf{Élargir aux conséquences sociales et psychologiques} : stigmatisation, abandon scolaire, culpabilité, dépression ; risque pénal pour la patiente et le praticien (Code pénal camerounais).
\item \textbf{Dégager les alternatives} : prévention par l'éducation et la contraception, prise en charge médicale des grossesses à risque dans les structures sanitaires, soutien psychologique et social.
\item \textbf{Rester factuel et respectueux} : le sujet est sensible ; employer le vocabulaire médical, jamais de jugement moral dans une copie de sciences.
\end{enumerate}
\end{methode}

\begin{exoresolu}[Avortements clandestins : analyse d'une étude hospitalière]
\textbf{Énoncé.} Un service de gynécologie d'un hôpital de Yaoundé a admis en un an 150 patientes pour complications d'avortements clandestins. Le tableau suivant répartit les complications observées :

\begin{center}
\begin{tabularx}{\linewidth}{|X|c|}
\hline
\rowcolor{popblueL} \textbf{Complication} & \textbf{Nombre de cas} \\
\hline
Infections graves (métrite, septicémie) & 60 \\
\hline
Hémorragies & 45 \\
\hline
Perforations utérines & 15 \\
\hline
Décès & 6 \\
\hline
Autres complications & 24 \\
\hline
\end{tabularx}
\end{center}

\begin{enumerate}
\item Calculer le pourcentage de décès parmi les patientes admises.
\item Expliquer, par un mécanisme, l'origine des infections graves.
\item Citer deux conséquences à long terme pour une patiente qui survit à une perforation utérine ou à une infection grave.
\item Proposer deux mesures de prévention à l'échelle d'un établissement scolaire.
\end{enumerate}
\tcblower
\textbf{Solution détaillée.}

\textbf{1. Pourcentage de décès.}
\[ p = \frac{6}{150} \times 100 = 4\,\%. \]
4 \% des patientes admises pour complications d'avortement clandestin en sont décédées : une mortalité considérable pour un acte évitable.

\textbf{2. Mécanisme des infections graves.} L'avortement clandestin est réalisé dans des conditions d'\textbf{asepsie nulle} : instruments non stérilisés (ou corps étrangers comme des tiges ou du manioc), environnement souillé, opérateur non formé. Les micro-organismes (streptocoques, staphylocoques, \emph{Clostridium}) introduits dans le vagin franchissent le col utérin et colonisent la cavité utérine, dont la muqueuse est fragilisée et saignante : c'est la \textbf{métrite}. Si l'infection gagne le sang, elle devient une \textbf{septicémie}, potentiellement mortelle. L'évacuation incomplète (restes de tissus) entretient l'infection.

\textbf{3. Conséquences à long terme.} (a) La \textbf{stérilité} : les infections et les cicatrices utérines (synéchies) obstruent les trompes ou empêchent la nidation ; (b) les \textbf{séquelles chirurgicales} : une perforation peut imposer une hystérectomie (ablation de l'utérus) d'urgence, supprimant définitivement toute possibilité de grossesse. On peut ajouter les séquelles psychologiques (dépression, anxiété, trouble de stress post-traumatique).

\textbf{4. Mesures de prévention en milieu scolaire.} (a) Organiser des \textbf{séances d'éducation à la santé reproductive} animées par des professionnels de santé (information sur la fertilité, les risques, la contraception) ; (b) mettre en place un \textbf{dispositif d'écoute et d'orientation} (assistante sociale, infirmerie scolaire, relais vers un centre de santé ou un CJAS -- Centre Jeunes Amour et Santé) pour que toute élève en difficulté trouve un conseil confidentiel avant de recourir à une solution dangereuse.

\textbf{Conclusion.} Avec 4 \% de décès et 40 \% d'infections graves, l'avortement clandestin est un drame de santé publique. Sa prévention repose sur l'éducation, l'accès à la contraception et l'orientation vers les structures de santé légales.
\end{exoresolu}

\courssub{Moyens de contraception : panorama et choix responsable}

\begin{definition}[Contraception]
Ensemble des moyens permettant d'éviter une grossesse. On distingue :
\begin{itemize}
\item \textbf{Méthodes barrières / mécaniques :} préservatif masculin (latex/polyuréthane), préservatif féminin, diaphragme, cape cervicale. Seuls les préservatifs protègent des IST.
\item \textbf{Méthodes hormonales :} pilule combinée (œstrogène + progestatif) ou micro-progestative (progestatif seul), injectable (tous les 3 mois), implant sous-cutané (3--5 ans), patch, anneau vaginal. Bloquent l'ovulation, épaississent la glaire, atrophient l'endomètre. \textbf{Nécessitent un avis médical} (contre-indications : thrombose, cancer hormono-dépendant, tabagisme > 35 ans pour pilule combinée).
\item \textbf{Méthodes intra-utérines (DIU / stérilet) :} au cuivre (effet spermicide, 5--10 ans) ou hormonal (lévonorgestrel, 3--5 ans). Posés par un professionnel.
\item \textbf{Méthodes naturelles :} abstinence périodique (calendrier, température, Billings). Taux d'échec élevé, cycles irréguliers = peu fiables, surtout à l'adolescence.
\item \textbf{Méthodes définitives :} ligature des trompes (femme), vasectomie (homme). Irréversibles, pour personnes ayant réalisé leur projet parental.
\item \textbf{Contraception d'urgence :} pilule du lendemain (lévonorgestrel ou ulipristal) jusqu'à 72--120 h après rapport non protégé. \textbf{Pas une méthode régulière}.
\end{itemize}
\end{definition}

\begin{methode}[Analyser un tableau de moyens contraceptifs]
\begin{enumerate}
\item \textbf{Classer les méthodes} en catégories : \emph{mécaniques/barrières} (préservatifs, diaphragme), \emph{hormonales} (pilule, injectable, implant, patch), \emph{intra-utérines} (stérilet), \emph{naturelles} (abstinence périodique, méthode des températures — peu fiables), \emph{définitives} (ligature des trompes, vasectomie).
\item \textbf{Comparer selon trois critères} : efficacité (taux d'échec en usage courant), protection contre les IST (seul le préservatif en offre une), réversibilité et mode d'obtention (ordonnance, centre de santé).
\item \textbf{Relier le choix à la situation} : un adolescent exposé aux IST a besoin du préservatif ; une femme qui ne désire plus d'enfant peut opter pour une méthode définitive après conseil médical.
\item \textbf{Signaler les limites} : aucune méthode hormonale ne protège des IST ; les méthodes naturelles ont un taux d'échec élevé ; toute contraception hormonale exige un avis médical.
\item \textbf{Rédiger la réponse} en citant le critère qui motive le choix.
\end{enumerate}
\end{methode}

\begin{exoresolu}[Choisir une contraception adaptée à un profil]
\textbf{Énoncé.} Le tableau ci-dessous présente quatre méthodes contraceptives et leur taux d'échec en usage courant (pour 100 femmes sur un an).

\begin{center}
\begin{tabularx}{\linewidth}{|l|X|c|c|}
\hline
\rowcolor{popblueL} \textbf{Méthode} & \textbf{Mode d'action} & \textbf{Taux d'échec} & \textbf{Protège des IST} \\
\hline
Préservatif masculin & Barrière empêchant le contact sperme-voies génitales & 13 \% & Oui \\
\hline
Pilule combinée & Blocage de l'ovulation par les hormones & 7 \% & Non \\
\hline
Implant sous-cutané & Libération continue de progestatif bloquant l'ovulation & 0,1 \% & Non \\
\hline
Méthode du calendrier & Abstinence pendant la période supposée fertile & 23 \% & Non \\
\hline
\end{tabularx}
\end{center}

\begin{enumerate}
\item Quelle méthode conseiller à un couple de jeunes non dépistés pour le VIH ? Justifier.
\item Quelle méthode est la plus efficace contre la grossesse ? Pourquoi n'est-elle pas suffisante à elle seule pour un jeune couple ?
\item Pourquoi la méthode du calendrier est-elle déconseillée aux adolescentes ?
\end{enumerate}
\tcblower
\textbf{Solution détaillée.}

\textbf{1. Conseil pour un couple non dépisté.} Le \textbf{préservatif masculin} est la méthode à conseiller, car c'est la \emph{seule} du tableau qui protège des IST en même temps que de la grossesse (colonne « Protège des IST » : Oui). En l'absence de dépistage, le risque infectieux est réel ; la barrière mécanique empêche le contact entre les sécrétions génitales des deux partenaires. Son efficacité contraceptive (13 \% d'échec en usage courant, souvent par mauvaise utilisation) est améliorée par un usage correct et systématique.

\textbf{2. Méthode la plus efficace.} L'\textbf{implant sous-cutané}, avec seulement 0,1 \% d'échec, est la méthode la plus efficace contre la grossesse : il libère en continu un progestatif qui bloque l'ovulation pendant plusieurs années, sans dépendre d'un geste quotidien. Cependant, il ne protège \emph{pas} des IST (colonne : Non). Pour un jeune couple non dépisté, il est donc insuffisant seul : il faudrait lui associer le préservatif (stratégie dite de « double protection »).

\textbf{3. Limites de la méthode du calendrier.} Avec 23 \% d'échec, c'est la méthode la moins fiable. Elle suppose un cycle régulier et une ovulation prévisible ; or chez l'adolescente, les cycles sont souvent irréguliers pendant les années qui suivent les premières règles, et la date d'ovulation varie. De plus, les spermatozoïdes survivent jusqu'à 5 jours dans les voies génitales féminines, ce qui élargit la période fertile réelle. Elle est donc déconseillée aux adolescentes.

\textbf{Conclusion.} Le choix d'une contraception dépend du profil : préservatif pour la double protection, implant pour l'efficacité maximale contre la grossesse (associé au préservatif si risque d'IST), et rejet des méthodes naturelles peu fiables chez les jeunes.
\end{exoresolu}

\courssub{Stratégies de prévention et alternatives responsables}

\begin{methode}[Proposer des alternatives responsables à un comportement à risque]
\begin{enumerate}
\item \textbf{Reformuler le problème} en une phrase : qui est exposé, à quel risque, dans quel contexte.
\item \textbf{Proposer des alternatives réalistes et graduées} :
\begin{itemize}
\item \emph{Prévention primaire} : information et éducation à la santé reproductive, abstinence ou report de la première relation, communication parent-enfant ;
\item \emph{Protection} : usage correct et systématique du préservatif (masculin ou féminin), contraception adaptée (pilule, injectable, implant) après conseil dans un centre de santé ;
\item \emph{Recours aux services} : consultation dans un centre de santé ou un centre jeunes (CJAS), dépistage volontaire du VIH, prise en charge médicale légale en cas de grossesse à risque.
\end{itemize}
\item \textbf{Justifier chaque proposition} par son effet : « le préservatif, utilisé correctement, empêche à la fois la grossesse et la transmission des IST ».
\item \textbf{Adapter au contexte local} : disponibilité des structures (hôpitaux de district, centres de santé intégrés, campagnes de sensibilisation), gratuité ou faible coût du dépistage et des préservatifs lors des campagnes (ex. Journée Mondiale de Lutte contre le Sida, vacances « Zéro Grossesse à l'École »).
\item \textbf{Conclure par un message de responsabilisation} : la santé reproductive est un choix éclairé, pas une fatalité.
\end{enumerate}
\textbf{Réflexe} : une alternative non justifiée ou irréaliste (hors de portée du jeune) ne rapporte pas tous les points.
\end{methode}

\begin{exoresolu}[Construire une réponse responsable à une situation à risque]
\textbf{Énoncé.} Reprenons la situation d'Amina (exercice précédent) : 16 ans, rapports non protégés avec un partenaire adulte, croyance au lavage vaginal protecteur. En tant qu'élève de Première sensibilisé à la santé reproductive, proposer à Amina trois alternatives responsables, en justifiant chacune.
\tcblower
\textbf{Solution détaillée.}

\textbf{Alternative 1 : s'informer et rompre avec les fausses croyances.} Amina doit d'abord comprendre que le lavage vaginal ne protège ni de la grossesse ni des IST : les spermatozoïdes atteignent l'utérus en quelques minutes et les agents infectieux pénètrent les muqueuses pendant le rapport. Elle peut s'informer auprès d'un centre de santé, d'une infirmière scolaire ou d'une association de jeunes. \emph{Justification} : l'information exacte est la première barrière contre les comportements à risque.

\textbf{Alternative 2 : exiger l'usage du préservatif à chaque rapport, ou reporter les relations sexuelles.} Le préservatif, masculin ou féminin, est le seul moyen de « double protection » : il empêche la rencontre entre spermatozoïdes et ovocyte (contraception) et bloque les agents infectieux (prévention des IST dont le VIH). Si le partenaire refuse, Amina a le droit de refuser le rapport ; le report des relations sexuelles (abstinence) reste la protection la plus sûre à son âge. \emph{Justification} : à 16 ans, une grossesse exposerait Amina à des complications obstétricales (bassin immature) et à une déscolarisation quasi certaine.

\textbf{Alternative 3 : se rendre dans un centre de santé pour un dépistage et un conseil personnalisé.} Ayant déjà eu des rapports non protégés, Amina devrait effectuer un \textbf{dépistage volontaire du VIH} et des IST (souvent gratuit lors des campagnes), et recevoir un conseil contraceptif adapté (pilule, injectable ou implant) si elle maintient une vie sexuelle. \emph{Justification} : le dépistage précoce permet une prise en charge rapide (ARV si séropositivité) ; la contraception médicalisée, choisie avec un professionnel, est fiable et réversible.

\textbf{Conclusion.} Face aux rapports non protégés, les alternatives responsables sont : l'information exacte, la double protection par le préservatif (ou le report des relations) et le recours aux services de santé pour le dépistage et la contraception. Ces choix préservent la santé, la scolarité et l'avenir d'Amina.
\end{exoresolu}

\begin{savaistu}[Contexte camerounais : ressources et cadres légaux]
\begin{itemize}
\item \textbf{Cadre légal :} La loi n° 2016/007 du 12 juillet 2016 relative au Code pénal (art. 337--339) réprime l'avortement, mais l'autorise si la grossesse met en danger la vie de la mère (avis de deux médecins) ou résulte d'un viol (attestation du procureur).
\item \textbf{Structures de proximité :} Centres de Santé Intégrés (CSI), Hôpitaux de District (HD), Hôpitaux Régionaux. Existence de \textbf{CJAS} (Centres Jeunes Amour et Santé) dans plusieurs régions (Yaoundé, Douala, Bafoussam, Garoua, Maroua...) offrant accueil, écoute, dépistage, contraception, gratuitement ou à coût réduit.
\item \textbf{Campagnes nationales :} « Vacances sans grossesses », « Zéro grossesse à l'école », Journée Mondiale de la Contraception (26 sept), Journée Mondiale du Sida (1er déc). Distribution gratuite de préservatifs (masculins/féminins) et de lubrifiants.
\item \textbf{Numéros utiles :} 117 (Police/Secours), 1510 (Allô Santé - Ministère de la Santé Publique), numéros verts des ONG (ex. CAMNAFAW, Association Camerounaise pour le Bien-Être Familial).
\end{itemize}
\end{savaistu}

\begin{attention}[Les 5 erreurs qui coûtent des points au Probatoire]
\begin{enumerate}
\item \textbf{Identifier sans justifier.} Écrire « comportement à risque : rapports non protégés » sans citer le fait du texte ni le risque encouru ne vaut que la moitié des points. Toujours rédiger : fait + nom du comportement + risque.
\item \textbf{Confondre contraception et protection contre les IST.} Affirmer que la pilule ou l'implant « protège des maladies » est faux : seul le \cle{préservatif} assure la double protection. C'est l'erreur la plus fréquente.
\item \textbf{Oublier le $\times 100$ ou mal poser le pourcentage.} Un pourcentage se calcule sur l'effectif \emph{du groupe concerné}, pas sur l'effectif total : pour comparer les complications, on calcule $\frac{96}{240}$ et $\frac{80}{560}$, et non $\frac{96}{800}$. Vérifier aussi que le résultat est cohérent (inférieur à 100 \%).
\item \textbf{Employer un vocabulaire vague ou moralisateur.} « Elle fait des bêtises », « c'est mal » : à bannir. Le correcteur attend les termes scientifiques : \emph{rapport non protégé}, \emph{grossesse précoce}, \emph{dystocie}, \emph{septicémie}, \emph{fistule obstétricale}, \emph{avortement clandestin}.
\item \textbf{Proposer des alternatives irréalistes ou non justifiées.} « Il faut arrêter » n'est pas une alternative. Chaque proposition doit être concrète (préservatif, dépistage, consultation, éducation), adaptée au contexte (centre de santé, CJAS, campagne de dépistage gratuit) et justifiée par son effet sur le risque.
\end{enumerate}
\end{attention}

\newpage

\coursec{Exercices}

\courssub{Je vérifie mes connaissances}

\begin{exercice}[QCM : efficacité des méthodes contraceptives]
Pour chaque question, choisir la (ou les) bonne(s) réponse(s).
\begin{enumerate}
\item L'indice de Pearl mesure :
\begin{enumerate}
\item le nombre de spermatozoïdes par éjaculat ;
\item le nombre de grossesses survenues chez 100 femmes utilisant une méthode pendant un an ;
\item la durée moyenne d'un cycle menstruel ;
\item le taux d'hormones dans le sang.
\end{enumerate}
\item Parmi les méthodes suivantes, laquelle protège à la fois des grossesses non désirées et des infections sexuellement transmissibles (IST) ?
\begin{enumerate}
\item la pilule oestroprogestative ;
\item le stérilet au cuivre ;
\item le préservatif masculin ;
\item la méthode des températures.
\end{enumerate}
\item La fécondation a normalement lieu :
\begin{enumerate}
\item dans l'utérus ;
\item dans le tiers externe de la trompe de Fallope ;
\item dans l'ovaire ;
\item dans le vagin.
\end{enumerate}
\item La contraception d'urgence (pilule du lendemain) est d'autant plus efficace qu'elle est prise :
\begin{enumerate}
\item au moins 7 jours après le rapport ;
\item le plus tôt possible après le rapport non protégé ;
\item uniquement le jour de l'ovulation ;
\item après un test de grossesse positif.
\end{enumerate}
\end{enumerate}
\end{exercice}

\begin{exercice}[Vrai ou faux — justifier]
Répondre par VRAI ou FAUX, puis justifier chaque réponse en une ou deux phrases.
\begin{enumerate}
\item Une fille ne peut pas tomber enceinte lors de son tout premier rapport sexuel.
\item L'interruption volontaire de grossesse réalisée hors du cadre médical légal expose à des hémorragies et à des infections pouvant entraîner la stérilité ou la mort.
\item Le retrait (coït interrompu) est une méthode contraceptive fiable.
\item Une grossesse chez une adolescente de 14 ans présente plus de risques médicaux qu'une grossesse chez une femme de 25 ans.
\end{enumerate}
\end{exercice}

\begin{exercice}[Définitions]
Définir brièvement et précisément chacun des termes suivants :
\begin{enumerate}
\item contraception ;
\item indice de Pearl ;
\item grossesse précoce ;
\item avortement clandestin ;
\item infection sexuellement transmissible (IST).
\end{enumerate}
\end{exercice}

\begin{exercice}[Calcul direct : lecture d'un indice de Pearl]
On donne les indices de Pearl (usage typique) de trois méthodes : préservatif masculin : 13 ; pilule oestroprogestative : 8 ; implant sous-cutané : 0,1.
\begin{enumerate}
\item Pour 100 femmes utilisant le préservatif pendant un an, combien de grossesses non désirées peut-on attendre en moyenne ?
\item Pour 1000 femmes utilisant l'implant pendant un an, combien de grossesses peut-on attendre ?
\item Classer ces trois méthodes de la plus efficace à la moins efficace.
\end{enumerate}
\end{exercice}

\courssub{Je m'entraîne}

\begin{exercice}[Identifier les comportements à risque]
Voici quatre situations observées dans un lycée de Yaoundé :
\begin{itemize}
\item Situation A : Marlyse, 16 ans, a des rapports non protégés car son partenaire affirme que « le retrait suffit ».
\item Situation B : Jean, 17 ans, utilise systématiquement un préservatif avec sa partenaire.
\item Situation C : Nadège, 15 ans, enceinte de 3 mois, absorbe une décoction vendue au marché pour « faire revenir les règles ».
\item Situation D : Alain, 18 ans, multiplie les partenaires sans protection et refuse le dépistage du VIH.
\end{itemize}
\begin{enumerate}
\item Identifier, pour chaque situation, s'il s'agit d'un comportement à risque ou d'un comportement responsable.
\item Pour chaque comportement à risque, citer deux conséquences sanitaires possibles.
\item Proposer une alternative responsable pour les situations A, C et D.
\end{enumerate}
\end{exercice}

\begin{exercice}[Classer des méthodes contraceptives]
On donne le tableau suivant (indices de Pearl en usage typique) :

\begin{center}
\begin{tabularx}{\linewidth}{|l|X|}\hline
\rowcolor{popblueL} \textbf{Méthode} & \textbf{Indice de Pearl (grossesses pour 100 femmes-an)} \\ \hline
Implant sous-cutané & 0,1 \\ \hline
Stérilet au cuivre & 0,8 \\ \hline
Pilule oestroprogestative & 8 \\ \hline
Préservatif masculin & 13 \\ \hline
Méthode Ogino (abstinence périodique) & 24 \\ \hline
\end{tabularx}
\end{center}

\begin{enumerate}
\item Classer ces méthodes de la plus efficace à la moins efficace contre la grossesse.
\item Quelle(s) méthode(s) protège(nt) également des IST ? Justifier.
\item Expliquer, à l'aide des connaissances sur le cycle ovarien, pourquoi la méthode Ogino est la moins fiable, en particulier chez une adolescente aux cycles irréguliers.
\end{enumerate}
\end{exercice}

\begin{exercice}[Fenêtre de fécondité]
Le cycle ovarien d'une jeune femme dure régulièrement 28 jours, l'ovulation ayant lieu au 14\textsuperscript{e} jour. On rappelle que les spermatozoïdes survivent environ 3 jours dans les voies génitales féminines et que l'ovule n'est fécondable que pendant environ 24 heures.
\begin{enumerate}
\item Déterminer la fenêtre de fécondité (jours du cycle pendant lesquels un rapport non protégé peut aboutir à une grossesse).
\item Représenter cette fenêtre sur un axe gradué de J1 à J28.
\item En déduire pourquoi un rapport non protégé au 12\textsuperscript{e} jour du cycle est à haut risque de grossesse.
\end{enumerate}
\end{exercice}

\begin{exercice}[Conduite à tenir après un oubli de pilule]
Amina, 19 ans, prend une pilule oestroprogestative (plaquette de 21 comprimés). Elle constate un oubli du 14\textsuperscript{e} comprimé (milieu de plaquette) datant de plus de 12 heures, et elle a eu un rapport non protégé la veille.
\begin{enumerate}
\item Rappeler le mécanisme d'action contraceptive de la pilule oestroprogestative.
\item Expliquer pourquoi un oubli de plus de 12 heures en milieu de plaquette fait courir un risque de grossesse.
\item Proposer une conduite à tenir raisonnée (reprise du comprimé, protection complémentaire, recours éventuel à la contraception d'urgence) en justifiant chaque étape.
\end{enumerate}
\end{exercice}

\courssub{Je me prépare au Probatoire}

\begin{exercice}[Type Probatoire — Étude de cas : les conséquences d'un avortement clandestin]
Lors d'une consultation au Centre Mère-Enfant de l'hôpital de district de Bafoussam, Sandrine, 16 ans, élève en classe de Première, consulte pour des douleurs pelviennes chroniques. L'anamnèse révèle qu'à 15 ans, à la suite d'une grossesse non désirée, elle a subi un avortement clandestin pratiqué par une « matrone » à l'aide d'une sonde non stérile. Le bilan actuel montre une obstruction des deux trompes de Fallope (stérilité tubaire).
\begin{enumerate}
\item Définir les termes : avortement clandestin, stérilité tubaire.
\item Rappeler le rôle des trompes de Fallope dans la reproduction humaine et le site normal de la fécondation.
\item Expliquer le lien physiopathologique entre l'avortement clandestin subi par Sandrine et sa stérilité tubaire (enchaînement : infection $\rightarrow$ salpingite $\rightarrow$ obstruction).
\item Citer trois autres complications immédiates ou tardives des avortements clandestins.
\item Proposer trois mesures de prévention qui, mises en œuvre dans son lycée et sa communauté, auraient pu éviter ce drame.
\end{enumerate}
\end{exercice}

\begin{exercice}[Type Probatoire — Exploitation de données et argumentaire]
Une enquête réalisée dans trois établissements scolaires de la région du Littoral auprès de 300 élèves de 15 à 19 ans a donné les résultats suivants :

\begin{center}
\begin{tabularx}{\linewidth}{|X|c|}\hline
\rowcolor{popblueL} \textbf{Situation déclarée} & \textbf{Pourcentage des élèves} \\ \hline
Ont déjà eu un rapport sexuel & 45 \% \\ \hline
Ont utilisé un préservatif lors du dernier rapport & 30 \% \\ \hline
Connaissent l'existence d'un centre de santé adapté aux jeunes & 20 \% \\ \hline
Ont déjà été confrontés (eux ou une connaissance) à une grossesse non désirée & 18 \% \\ \hline
\end{tabularx}
\end{center}

\begin{enumerate}
\item Calculer le nombre d'élèves ayant déjà eu un rapport sexuel et le nombre d'élèves ayant utilisé un préservatif lors du dernier rapport.
\item À partir de ces données, dégager deux comportements à risque fréquents chez ces élèves.
\item Expliquer deux conséquences sanitaires et deux conséquences sociales des grossesses précoces en milieu scolaire camerounais.
\item Un élève déclare : « Le préservatif réduit le plaisir, donc je ne l'utilise pas. » Rédiger une réponse structurée comportant trois arguments médicaux et un argument relationnel pour convaincre cet élève de changer de comportement.
\item Proposer, sous forme d'un plan en trois axes, une campagne de sensibilisation à mener dans ces établissements.
\end{enumerate}
\end{exercice}

\begin{exercice}[Type Probatoire — Cycle ovarien et choix contraceptif raisonné]
Le document ci-dessous décrit un cycle ovarien de 28 jours chez une jeune femme : la concentration sanguine des \oe strogènes augmente de J1 à J13, celle de la LH présente un pic brutal à J13, l'ovulation survient à J14, puis la progestérone domine de J15 à J28 ; l'endomètre s'épaissit de J6 à J21 puis se desquame (règles) de J1 à J5 si aucune fécondation n'a eu lieu.
\begin{enumerate}
\item Construire un schéma légendé du cycle ovarien faisant apparaître : l'axe des jours (J1 à J28), les variations des \oe strogènes, de la LH et de la progestérone, le moment de l'ovulation et les phases du cycle utérin.
\item Situer sur ce schéma la fenêtre de fécondité, en justifiant son étendue par la durée de vie des gamètes.
\item Expliquer le rôle du pic de LH dans le déclenchement de l'ovulation.
\item La jeune femme souhaite utiliser la méthode Ogino (abstinence pendant la période supposée fertile). En t'appuyant sur le schéma et sur l'indice de Pearl de cette méthode (24), expliquer pourquoi ce choix est risqué, en particulier si ses cycles sont irréguliers.
\item Proposer deux méthodes contraceptives plus fiables adaptées à une jeune femme nullipare, en précisant pour chacune le mécanisme d'action et l'avantage principal.
\end{enumerate}
\end{exercice}

\newpage

\coursec{Corrigés des exercices}

\courssub{Je vérifie mes connaissances}

\begin{corrige}[Exercice 1 — QCM : efficacité des méthodes contraceptives]
\begin{enumerate}
\item \textbf{Bonne réponse : b.} L'indice de Pearl est un indicateur statistique qui exprime le nombre de grossesses survenues chez 100 femmes utilisant une méthode contraceptive donnée pendant un an (usage typique). Plus l'indice est bas, plus la méthode est efficace.
\item \textbf{Bonne réponse : c.} Le préservatif masculin (et féminin) est la \emph{seule} méthode contraceptive qui offre une double protection : contraceptive (barrière physique aux spermatozoïdes) et prophylactique (barrière aux agents pathogènes des IST, dont le VIH). La pilule, le stérilet et les méthodes naturelles ne protègent pas des IST.
\item \textbf{Bonne réponse : b.} La fécondation (rencontre spermatozoïde-ovocyte II) a lieu physiologiquement dans le \textbf{tiers externe de la trompe de Fallope} (ampoule tubaire). L'ovocyte y est capté par les franges ; les spermatozoïdes y remontent par le col et l'utérus.
\item \textbf{Bonne réponse : b.} La contraception d'urgence (pilule du lendemain à base de lévonorgestrel ou d'acétate d'ulipristal) agit principalement en retardant ou en inhibant l'ovulation. Son efficacité est maximale si elle est prise \textbf{le plus tôt possible} après le rapport non protégé (idéalement dans les 12-24 h), et décroît rapidement avec le temps (généralement jusqu'à 72 h ou 120 h selon le produit). Elle est inefficace si l'ovulation a déjà eu lieu ou si la grossesse est déjà implantée.
\end{enumerate}
\end{corrige}

\begin{corrige}[Exercice 2 — Vrai ou faux — justifier]
\begin{enumerate}
\item \textbf{FAUX.} La fécondité féminine apparaît dès les premières ovulations (puberté), souvent avant même les premières règles (ménarche). Un premier rapport sexuel non protégé pendant la période fertile peut tout à fait entraîner une grossesse.
\item \textbf{VRAI.} L'avortement clandestin, pratiqué hors cadre médical légal, par du personnel non qualifié et avec du matériel non stérile, expose à des complications graves immédiates : hémorragies massives, perforations utérines, infections sévères (septicémie, péritonite). Ces séquelles peuvent détruire l'endomètre et les trompes, entraînant une \textbf{stérilité définitive}, voire le \textbf{décès} de la jeune fille/femme.
\item \textbf{FAUX.} Le retrait (coït interrompu) est une méthode \textbf{peu fiable} (indice de Pearl usage typique $\approx$ 20-22). Deux causes principales : présence de spermatozoïdes dans le liquide pré-éjaculatoire (sécrété par les glandes de Cowper) et difficulté à retirer le pénis \emph{juste} avant l'éjaculation (perte de contrôle). Il ne protège pas des IST.
\item \textbf{VRAI.} La grossesse chez l'adolescente de 14 ans (bassin immature, croissance non achevée, statut nutritionnel souvent fragile) est une \textbf{grossesse à haut risque} : éclampsie, anémie sévère, dystocie fœtopelvienne (césarienne fréquente), accouchement prématuré, bas poids de naissance, mortalité maternelle et périnatale majorées par rapport à une femme de 25 ans (âge physiologique optimal).
\end{enumerate}
\end{corrige}

\begin{corrige}[Exercice 3 — Définitions]
\begin{enumerate}
\item \textbf{Contraception :} Ensemble des méthodes (mécaniques, hormonales, chirurgicales, comportementales) visant à empêcher la survenue d'une grossesse lors de rapports sexuels, de manière réversible ou définitive.
\item \textbf{Indice de Pearl :} Nombre de grossesses non désirées observées chez 100 femmes utilisant une méthode contraceptive donnée pendant une durée d'un an (exprimé en « grossesses pour 100 femmes-an »). Il permet de comparer l'efficacité des méthodes en conditions réelles d'usage (usage typique) ou théoriques (usage parfait).
\item \textbf{Grossesse précoce :} Grossesse survenant chez une adolescente, généralement définie comme avant 18-20 ans (selon les contextes), alors que le développement physique, psychologique et socio-économique n'est pas achevé, exposant la mère et l'enfant à des risques sanitaires et sociaux majeurs.
\item \textbf{Avortement clandestin :} Interruption volontaire de grossesse réalisée en dehors du cadre légal et médical autorisé, par une personne non habilitée, dans des conditions d'asepsie et de sécurité insuffisantes, exposant la femme à des risques vitaux.
\item \textbf{Infection sexuellement transmissible (IST) :} Maladie causée par des agents pathogènes (bactéries, virus, parasites, champignons) se transmettant principalement par voie sexuelle (rapports vaginaux, anaux, oraux non protégés). Exemples : VIH/sida, syphilis, gonorrhée, chlamydiose, hépatite B, HPV, herpès génital.
\end{enumerate}
\end{corrige}

\begin{corrige}[Exercice 4 — Calcul direct : lecture d'un indice de Pearl]
\begin{enumerate}
\item L'indice de Pearl du préservatif masculin (usage typique) est de \textbf{13}. Cela signifie qu'en moyenne, \textbf{13 grossesses non désirées} sont attendues pour 100 femmes utilisant cette méthode pendant un an.
\item L'indice de Pearl de l'implant sous-cutané est de \textbf{0,1} (0,1 grossesse pour 100 femmes-an). Pour 1000 femmes (10 fois plus), on attend proportionnellement : $0,1 \times 10 = \textbf{1 grossesse}$ en moyenne par an.
\item Classement de la plus efficace (indice le plus bas) à la moins efficace (indice le plus haut) :
\begin{enumerate}
\item Implant sous-cutané (0,1) — \emph{très haute efficacité}
\item Pilule oestroprogestative (8) — \emph{haute efficacité}
\item Préservatif masculin (13) — \emph{efficacité modérée (mais double protection)}
\end{enumerate}
\end{enumerate}
\begin{aretenir}
L'indice de Pearl s'interprète toujours « plus il est bas, plus la méthode est efficace ». Attention : l'usage typique (vie réelle) donne des indices plus élevés que l'usage parfait (théorique) à cause des oublis, erreurs d'utilisation, interactions médicamenteuses.
\end{aretenir}
\end{corrige}

\courssub{Je m'entraîne}

\begin{corrige}[Exercice 5 — Identifier les comportements à risque]
\begin{enumerate}
\item \textbf{Analyse des situations :}
\begin{center}
\begin{tabularx}{\linewidth}{|l|X|X|}\hline
\rowcolor{popblueL} \textbf{Situation} & \textbf{Type de comportement} & \textbf{Justification} \\ \hline
A (Marlyse) & \textbf{À risque} & Confiance en la méthode du retrait (peu fiable), absence de protection barrière ni hormonale. \\ \hline
B (Jean) & \textbf{Responsable} & Utilisation systématique du préservatif (double protection grossesse + IST). \\ \hline
C (Nadège) & \textbf{À risque} & Recours à une méthode abortive traditionnelle non stérile, inefficace et dangereuse (avortement clandestin). \\ \hline
D (Alain) & \textbf{À risque} & Multiplication des partenaires, absence de préservatif, refus du dépistage (ignorance sérologie VIH/IST). \\ \hline
\end{tabularx}
\end{center}
\item \textbf{Conséquences sanitaires possibles (deux par situation à risque) :}
\begin{itemize}
\item \textbf{A :} Grossesse non désirée ; contamination par une IST (VIH, HPV, chlamydiae, etc.).
\item \textbf{C :} Hémorragie grave, infection pelvienne (salpingite) $\rightarrow$ stérilité tubaire, péritonite, choc septique, décès ; séquelles psychologiques.
\item \textbf{D :} Infection par le VIH/sida ou autres IST (syphilis, hépatite B, gonorrhée) ; transmission à ses partenaires ; grossesses non désirées chez ses partenaires.
\end{itemize}
\item \textbf{Alternatives responsables :}
\begin{itemize}
\item \textbf{A (Marlyse) :} Adopter une contraception efficace (pilule, implant, stérilet) \textbf{associée} au préservatif masculin pour la double protection ; consulter un centre de santé jeune (ex. Centre de Santé de l'Adolescent) pour un conseil personnalisé.
\item \textbf{C (Nadège) :} Se rendre \textbf{immédiatement} dans une structure de santé habilitée (hôpital de district, centre mère-enfant) pour une prise en charge médicale de l'avortement (aspiration manuelle intra-utérine, médicaments) dans des conditions de sécurité ; bénéficier d'un conseil post-avortement et d'une contraception adaptée.
\item \textbf{D (Alain) :} Utiliser \textbf{systématiquement le préservatif} à chaque rapport ; réduire le nombre de partenaires / privilégier la fidélité mutuelle ; accepter le \textbf{dépistage VIH/IST} régulier (gratuit et confidentiel dans les centres agréés) ; se faire traiter en cas d'IST dépistée (avec son/ses partenaire(s)).
\end{itemize}
\end{enumerate}
\end{corrige}

\begin{corrige}[Exercice 6 — Classer des méthodes contraceptives]
\begin{enumerate}
\item \textbf{Classement par efficacité croissante (indice de Pearl décroissant) :}
\begin{enumerate}
\item Implant sous-cutané (0,1)
\item Stérilet au cuivre (0,8)
\item Pilule oestroprogestative (8)
\item Préservatif masculin (13)
\item Méthode Ogino / abstinence périodique (24)
\end{enumerate}
\item \textbf{Seul le préservatif masculin} (et le préservatif féminin) protège simultanément contre la grossesse \textbf{et} les IST.
\textbf{Justification :} C'est une méthode de \textbf{barrière physique} (latex ou polyuréthane) qui empêche le contact direct entre les muqueuses et les fluides biologiques (sperme, sécrétions vaginales, sang), bloquant ainsi le passage des spermatozoïdes \emph{et} des agents pathogènes (VIH, bactéries, virus). Les autres méthodes (hormonales, DIU, naturelles) n'ont aucun effet barrière.
\item \textbf{Explication de la faible fiabilité de la méthode Ogino (indice 24) :}
La méthode Ogino (calendrier) consiste à s'abstenir pendant la période théorique fertile calculée sur la durée des cycles passés. Elle échoue car :
\begin{itemize}
\item L'ovulation n'est pas fixée au J14 chez toutes les femmes ; elle dépend de la phase folliculaire, très variable.
\item \textbf{Chez l'adolescente, les cycles sont souvent irréguliers} (immaturité de l'axe hypothalamo-hypophysaire), rendant toute prévision hasardeuse.
\item La \textbf{fenêtre de fécondité réelle} est large : survie des spermatozoïdes (3 à 5 jours) + durée de vie de l'ovocyte (24 h) $\approx$ 6 à 7 jours fertiles potentiels.
\item Le stress, la maladie, le voyage, la prise de poids modifient la date d'ovulation \emph{a posteriori}, le calendrier ne le « sait » pas.
\end{itemize}
\end{enumerate}
\end{corrige}

\begin{corrige}[Exercice 7 — Fenêtre de fécondité]
\begin{enumerate}
\item \textbf{Données :} Cycle 28 jours, ovulation J14. Survie spermatozoïdes $\approx$ 3 jours (donnée de l'énoncé). Fécondabilité ovocyte $\approx$ 24 h (1 jour).
\begin{itemize}
\item \textbf{Début de la fenêtre :} J14 (ovulation) $-$ 3 jours (survie max spermatozoïdes) = \textbf{J11}. Un rapport à J11 dépose des spermatozoïdes qui survivront jusqu'à J14 (ovulation).
\item \textbf{Fin de la fenêtre :} J14 (ovulation) $+$ 1 jour (durée de vie ovocyte) = \textbf{J15}. Un rapport à J15 peut féconder l'ovocyte libéré à J14.
\end{itemize}
\textbf{Fenêtre de fécondité : du J11 au J15 inclus (5 jours).}
\item \textbf{Représentation graphique :}
\begin{popfigure}
\begin{tikzpicture}[scale=0.75]
% Axes
\draw[->, thick] (0,0) -- (29,0) node[right] {Jour du cycle};
\draw[->, thick] (0,-0.5) -- (0,2.5) node[above] {};
% Graduations
\foreach \x in {1,5,10,11,12,13,14,15,20,25,28} {
    \draw (\x,0.1) -- (\x,-0.1) node[below, font=\small] {\x};
}
% Fenêtre de fécondité
\fill[popgreen!30] (11,0) rectangle (15,2);
\draw[popgreen, thick] (11,0) -- (11,2) node[midway, left] {Début};
\draw[popgreen, thick] (15,0) -- (15,2) node[midway, right] {Fin};
% Ovulation
\draw[poporange, very thick, dashed] (14,0) -- (14,2.2) node[above] {\textbf{Ovulation J14}};
% Légendes phases
\node[align=center, font=\small] at (7,2.2) {Phase folliculaire};
\node[align=center, font=\small] at (22,2.2) {Phase lutéale};
% Flèches survie
\draw[<->, popblue, thick] (11,1.5) -- node[above, font=\small, popblue] {3 j survie spermatozoïdes} (14,1.5);
\draw[<->, poppink, thick] (14,1) -- node[above, font=\small, poppink] {1 j vie ovocyte} (15,1);
\end{tikzpicture}
\legende{Fenêtre de fécondité (J11--J15) pour un cycle de 28 jours à ovulation J14. Zone verte : période à risque de grossesse.}
\end{popfigure}
\item \textbf{Justification du risque au J12 :} Le J12 se situe \textbf{à l'intérieur de la fenêtre de fécondité} (J11--J15). Les spermatozoïdes déposés lors d'un rapport au J12 survivront au moins jusqu'au J14 (jour de l'ovulation) et pourront féconder l'ovocyte fraîchement libéré. La probabilité de grossesse est donc \textbf{élevée}.
\end{enumerate}
\end{corrige}

\begin{corrige}[Exercice 8 — Conduite à tenir après un oubli de pilule]
\begin{enumerate}
\item \textbf{Mécanisme d'action de la pilule oestroprogestative (combinée) :}
\begin{enumerate}
\item \textbf{Blocage de l'ovulation (effet principal) :} L'apport exogène d'œstrogènes et de progestatif inhibe la sécrétion hypophysaire de FSH et LH (rétrocontrôle négatif) $\rightarrow$ pas de maturation folliculaire $\rightarrow$ pas de pic de LH $\rightarrow$ \textbf{pas d'ovulation}.
\item \textbf{Modification de la glaire cervicale :} Le progestatif la rend épaisse, visqueuse, infranchissable aux spermatozoïdes (effet barrière).
\item \textbf{Modification de l'endomètre :} Atrophie glandulaire, inadapté à l'implantation (effet de secours).
\end{enumerate}
\item \textbf{Risque lié à l'oubli > 12 h en milieu de plaquette (14\textsuperscript{e} comprimé) :}
Le 14\textsuperscript{e} comprimé correspond environ au \textbf{J14 du cycle} (période pré-ovulatoire critique). Un oubli de plus de 12 h fait chuter le taux sanguin d'hormones, levant le frein sur l'hypophyse. Un \textbf{pic de LH échappé} peut survenir $\rightarrow$ \textbf{ovulation possible}. Or, la patiente a eu un rapport non protégé \textbf{la veille (J13)} : des spermatozoïdes sont potentiellement présents et viables dans les trompes (survie 3-5 j) au moment où l'ovulation échappée pourrait se produire. \textbf{Risque réel de grossesse.}
\item \textbf{Conduite à tenir raisonnée (recommandations OMS / sociétés savantes) :}
\begin{enumerate}
\item \textbf{Prendre immédiatement le comprimé oublié} (le 14\textsuperscript{e}), même si cela implique d'en prendre deux en même temps (celui d'hier + celui du jour).
\item \textbf{Continuer la plaquette normalement} (un comprimé par jour) jusqu'à la fin.
\item \textbf{Protection complémentaire obligatoire :} Utiliser des \textbf{préservatifs} (ou s'abstenir) pendant les \textbf{7 jours suivants} (jusqu'au 20\textsuperscript{e} comprimé inclus), le temps que l'effet contraceptif ovulatoire soit rétabli.
\item \textbf{Contraception d'urgence (pilule du lendemain) :} \textbf{Fortement recommandée} ici car rapport non protégé dans les 5 jours précédant l'oubli (veille = J13). La pilule à l'ulipristal (EllaOne\textsuperscript{\textregistered}) est préférable jusqu'à 120 h (5 j) post-rapport ; le lévonorgestrel (NorLevo\textsuperscript{\textregistered}) jusqu'à 72 h. À prendre \textbf{au plus vite}.
\item \textbf{Test de grossesse :} À réaliser 3 semaines plus tard si règles absentes à l'arrêt de la plaquette (ou saignement de privation anormal).
\item \textbf{Conseil :} Éviter les oublis (alarme téléphone, application, routine) ; envisager une méthode « sans oubli » (implant, stérilet) si oublis répétés.
\end{enumerate}
\end{enumerate}
\begin{attention}[Piège fréquent]
Ne pas confondre oubli < 12 h (pas de risque, prendre aussitôt, pas de protection supplémentaire) et oubli > 12 h (risque, conduite ci-dessus). En fin de plaquette (derniers 7 comprimés actifs), l'oubli > 12 h impose de finir la plaquette, \textbf{sauter l'arrêt de 7 jours} (enchaîner directement la plaquette suivante) + protection 7 j. Ici, milieu de plaquette $\neq$ fin de plaquette.
\end{attention}
\end{corrige}

\courssub{Je me prépare au Probatoire}

\begin{corrige}[Exercice 9 — Type Probatoire : Étude de cas — Conséquences d'un avortement clandestin]
\textbf{Barème indicatif : 10 points}
\begin{enumerate}[itemsep=6pt]
\item \textbf{(2 pts) Définitions :}
\begin{itemize}
\item \textbf{Avortement clandestin :} Interruption volontaire de grossesse pratiquée en dehors du cadre légal, par une personne non qualifiée, dans des conditions d'hygiène et de sécurité précaires, exposant la femme à des risques vitaux.
\item \textbf{Stérilité tubaire :} Incapacité à concevoir due à une obstruction (partielle ou totale) ou une destruction fonctionnelle des deux trompes de Fallope, empêchant la rencontre spermatozoïde-ovocyte ou le transport de l'embryon vers l'utérus.
\end{itemize}
\item \textbf{(1,5 pt) Rôle des trompes et site de fécondation :}
Les trompes de Fallope assurent : (1) la captation de l'ovocyte II par les franges fimbriées ; (2) le \textbf{site normal de la fécondation} dans leur tiers externe (ampoule) ; (3) le transport de l'embryon (stades clivage $\rightarrow$ morula $\rightarrow$ blastocyste) vers la cavité utérine grâce à la ciliature et la péristaltique (4-5 jours).
\item \textbf{(3 pts) Lien physiopathologique (enchaînement causal) :}
\begin{enumerate}
\item \textbf{Infection initiale :} Instrument non stérile $\rightarrow$ introduction de germes (souvent polymicrobien : \emph{Chlamydia trachomatis}, \emph{Neisseria gonorrhoeae}, anaérobies, \emph{Escherichia coli}) dans la cavité utérine.
\item \textbf{Endométrite / Salpingite ascendante :} Les germes remontent de l'utérus vers les trompes $\rightarrow$ inflammation aiguë de la muqueuse tubaire (salpingite) : œdème, exsudat purulent, nécrose de l'épithélium cillé.
\item \textbf{Obstruction / Séquelles :} Guérison par fibrose cicatricielle $\rightarrow$ \textbf{obstruction de la lumière tubaire} (syphons, adhérences péritonéales, hydrosalpinx) $\rightarrow$ blocage du passage gamètes/embryon $\rightarrow$ \textbf{stérilité tubaire bilatérale} (stérilité absolue naturelle). Risque majoré de grossesse extra-utérine (GEU) si perméabilité partielle.
\end{enumerate}
\item \textbf{(1,5 pt) Trois autres complications (immédiates ou tardives) :}
\begin{itemize}
\item Hémorragie sévère / choc hémorragique (perforation utérine, lacération col/vagin).
\item Péritonite généralisée / septicémie $\rightarrow$ décès (urgence vitale).
\item Perforation utérine $\rightarrow$ lésions digestives (côlon, intestin grêle) ou vasculaires.
\item \emph{Tardives :} Douleurs pelviennes chroniques (syndrome d'adhérences), grossesses extra-utérines à répétition, stérilité définitive, troubles psychologiques (syndrome post-traumatique), exclusion sociale.
\end{itemize}
\item \textbf{(2 pts) Trois mesures de prévention (lycée / communauté) :}
\begin{enumerate}
\item \textbf{Éducation sexuelle complète (ESC) obligatoire au programme :} Connaissances sur fertilité, contraception, IST, consentement, genre, droits reproductifs ; animée par personnel formé (infirmiers scolaires, ONG partenaires).
\item \textbf{Accès effectif aux services de santé sexuelle et reproductive (SSRAJ) :} Permanence infirmière / sage-femme au lycée ou convention avec Centre de Santé de l'Adolescent (CSA) / Centre de Planification Familiale (CPF) proche ; distribution gratuite et confidentielle de préservatifs et contraceptifs (pilule, injectable, implant) ; dépistage IST/VIH gratuit.
\item \textbf{Environnement protecteur et réduction de la stigmatisation :} Lutte contre les grossesses en milieu scolaire (réintégration des mères élèves), implication des parents/leaders religieux/chefs traditionnels dans le dialogue intergénérationnel, application de la loi sur l'avortement sécurisé (cas viol, danger maternel, malformation fœtale) et prise en charge humaine des complications d'avortement (soins post-abortum sans jugement).
\end{enumerate}
\end{enumerate}
\begin{aretenir}[Points de vigilance Probatoire]
Citer explicitement les germes (\emph{Chlamydia}, \emph{Neisseria}, anaérobies). Décrire l'enchaînement \textbf{infection $\rightarrow$ salpingite $\rightarrow$ fibrose $\rightarrow$ obstruction}. Distinguer stérilité tubaire (mécanique) et stérilité hormonale. Proposer des mesures \textbf{concrètes, multisectorielles} (santé + éducation + communauté), pas seulement « sensibiliser ».
\end{aretenir}
\end{corrige}

\begin{corrige}[Exercice 10 — Type Probatoire : Exploitation de données et argumentaire]
\textbf{Barème indicatif : 12 points}
\begin{enumerate}[itemsep=6pt]
\item \textbf{(2 pts) Calculs :}
Effectif total $N = 300$ élèves.
\begin{itemize}
\item Élèves ayant déjà eu un rapport sexuel : $45\% \times 300 = \textbf{135 élèves}$.
\item Élèves ayant utilisé un préservatif lors du dernier rapport : $30\% \times 300 = \textbf{90 élèves}$.
\end{itemize}
\begin{attention}[Attention]
Le pourcentage (30\%) porte sur l'effectif total (300), pas seulement sur les 135 sexuellement actifs. Si la question portait sur « parmi les actifs », ce serait $90/135 \approx 67\%$. Lire l'intitulé du tableau : « Pourcentage des élèves ».
\end{attention}
\item \textbf{(2 pts) Deux comportements à risque fréquents :}
\begin{enumerate}
\item \textbf{Faible utilisation du préservatif :} Seulement 30\% de l'ensemble des élèves (soit $\approx$ 67\% des sexuellement actifs) l'ont utilisé au dernier rapport $\rightarrow$ exposition aux IST/VIH et grossesses non désirées.
\item \textbf{Méconnaissance des ressources de santé :} Seulement 20\% connaissent un centre de santé adapté aux jeunes $\rightarrow$ retard/absence de recours à la contraception, au dépistage, à la prise en charge d'une grossesse ou d'une IST.
\end{enumerate}
(Autre comportement possible : 18\% confrontés à grossesse non désirée = conséquence des deux précédents).
\item \textbf{(2 pts) Conséquences des grossesses précoces en milieu scolaire camerounais :}
\begin{itemize}
\item \textbf{Sanitaires :} (1) Mortalité maternelle majorée (éclampsie, hémorragies, avortements clandestins, sepsis) ; (2) Morbidité fœtale/néonatale : prématurité, bas poids de naissance ($< 2500$ g), asphyxie néonatale, mortalité périnatale ; (3) Risque de fistules obstétricales (bassin immature).
\item \textbf{Sociales :} (1) Décrochage / exclusion scolaire définitive (souvent non-réintégration malgré les circulaires) $\rightarrow$ perte de capital humain, pauvreté ; (2) Stigmatisation, rejet familial / conjugal forcé (mariage précoce) ; (3) Dépendance économique, perte d'autonomie, cycle intergénérationnel de pauvreté.
\end{itemize}
\item \textbf{(2 pts) Réponse structurée à l'élève :}
\begin{enumerate}
\item \textbf{Argument médical 1 (Protection double unique) :} « Le préservatif est la \emph{seule} méthode qui te protège \textbf{à la fois} contre le VIH/sida, les IST (syphilis, HPV, hépatite B...) \emph{et} la grossesse. La pilule ou l'implant ne te protègent pas des infections. »
\item \textbf{Argument médical 2 (Prévention cancers) :} « En te protégeant du HPV (papillomavirus), le préservatif réduit ton risque futur de \textbf{cancer du col de l'utérus} (pour ta partenaire) et de cancers ORL/anaux (pour toi). »
\item \textbf{Argument médical 3 (Sans effets systémiques) :} « Contrairement aux méthodes hormonales, le préservatif n'a \textbf{aucun effet secondaire métabolique, hormonal ou vasculaire} (pas de prise de poids, pas de risque thrombose, pas d'interaction médicamenteuse). »
\item \textbf{Argument relationnel :} « L'utiliser, c'est faire preuve de \textbf{responsabilité partagée} et de \textbf{respect} envers ta/ton partenaire : tu protèges sa santé et ton avenir commun. C'est un signe de maturité et de confiance, pas une marque de méfiance. »
\end{enumerate}
\item \textbf{(4 pts) Plan de campagne de sensibilisation (3 axes) :}
\begin{center}
\begin{tabularx}{\linewidth}{|l|X|}\hline
\rowcolor{popblueL} \textbf{Axe} & \textbf{Actions concrètes} \\ \hline
\textbf{1. Information \& Éducation (Savoir)} & \begin{itemize}\item Causeries éducatives mensuelles par infirmier(s) scolaire / pairs-éducateurs formés (thèmes : cycle, fertilité, contraception, IST, consentement, violences). \item Affichages permanents (panneaux, flyers) dans les lieux de vie (réfectoire, dortoirs, toilettes) : numéros verts, adresses centres jeunes, mode d'emploi préservatif. \item Intégration modules ESC dans temps forts (Journée mondiale sida, Journée fille, rentrées). \end{itemize} \\ \hline
\textbf{2. Accès aux Services \& Produits (Pouvoir)} & \begin{itemize}\item Distribution gratuite, anonyme, sans jugement de préservatifs (masculins/féminins) dans l'infirmerie, foyers, bureaux vie scolaire. \item Permanence hebdomadaire d'un(e) sage-femme / conseiller(ère) en planification familiale (conseil, prescription pilule/injectable, pose implant/stérilet, dépistage IST/VIH rapide). \item Convention de référencement rapide vers Centre de Santé de l'Adolescent (CSA) / Hôpital de district pour prise en charge grossesse, avortement sécurisé, violences. \end{itemize} \\ \hline
\textbf{3. Mobilisation Communautaire \& Environnement (Soutien)} & \begin{itemize}\item Création/animation de \textbf{Clubs Santé Jeunes} (pairs-éducateurs) pour briser le tabou, relayer l'info, accompagner vers les services. \item Sensibilisation parents / associations parents d'élèves / leaders religieux / chefs de quartier : « Protéger nos filles, c'est protéger l'avenir » (dialogue intergénérationnel). \item Plaidoyer auprès administration scolaire pour application effective circulaire réintégration mères élèves, lutte contre violences basées sur le genre (VBGE) en milieu scolaire. \end{itemize} \\ \hline
\end{tabularx}
\end{center}
\end{enumerate}
\end{corrige}

\begin{corrige}[Exercice 11 — Type Probatoire : Cycle ovarien et choix contraceptif raisonné]
\textbf{Barème indicatif : 14 points}
\begin{enumerate}[itemsep=6pt]
\item \textbf{(4 pts) Schéma légendé du cycle ovarien (28 jours) :}
\begin{popfigure}
\begin{tikzpicture}[scale=0.7]
% Axes principaux
\draw[->, thick] (0,0) -- (29,0) node[right] {Jours du cycle (J1 à J28)};
\draw[->, thick] (0,0) -- (0,5.5) node[above] {Concentration hormonale (arbitraire)};
% Graduations jours
\foreach \x in {1,5,10,11,12,13,14,15,16,17,20,25,28} {
    \draw (\x,0.1) -- (\x,-0.1) node[below, font=\small] {\x};
}
% Courbe Oestrogènes (E2) - poppink
\draw[poppink, very thick, smooth, domain=1:28, samples=100] plot (\x, {2.5*exp(-(\x-13)*(\x-13)/18) + 0.3*exp(-(\x-21)*(\x-21)/8)});
\node[poppink, font=\small] at (10,4.5) {\textbf{\oe strogènes (E2)}};
% Pic LH - poporange
\draw[poporange, very thick, smooth, domain=12:15, samples=50] plot (\x, {3.5*exp(-(\x-13.5)*(\x-13.5)/0.8)});
\draw[poporange, thick] (1,0.3) -- (12,0.3);
\draw[poporange, thick] (16,0.3) -- (28,0.3);
\node[poporange, font=\small] at (13.5,4.2) {\textbf{LH}};
% Pic LH flèche
\draw[poporange, <->, thick] (13.5,0.5) -- (13.5,3.8) node[midway, right] {\textbf{Pic LH J13}};
% Courbe Progestérone (P4) - popgreen
\draw[popgreen, very thick, smooth, domain=14:28, samples=100] plot (\x, {2.8*exp(-(\x-21)*(\x-21)/12)});
\draw[popgreen, thick] (1,0.2) -- (14,0.2);
\node[popgreen, font=\small] at (22,3.5) {\textbf{Progestérone (P4)}};
% Ovulation
\draw[popblue, very thick, dashed] (14,0) -- (14,5) node[above] {\textbf{Ovulation J14}};
% Phases ovariennes
\node[align=center, font=\small, popmuted] at (7,5) {Phase folliculaire};
\node[align=center, font=\small, popmuted] at (21,5) {Phase lutéale (Corps jaune)};
% Cycle utérin (barre sous axe)
\draw[very thick, popdark] (1,-0.8) -- (5,-0.8) node[midway, below, font=\small] {\textbf{Ménstrues (J1--J5)}};
\draw[very thick, poppink] (6,-0.8) -- (14,-0.8) node[midway, below, font=\small] {\textbf{Phase proliférative (J6--J14)}};
\draw[very thick, popgreen] (15,-0.8) -- (28,-0.8) node[midway, below, font=\small] {\textbf{Phase sécrétoire (J15--J28)}};
% Fenêtre fécondité
\fill[popgold!40] (11,-0.5) rectangle (15,-1.1);
\node[popdark, font=\small, align=center] at (13,-1.3) {\textbf{Fenêtre de fécondité (J11--J15)}};
\end{tikzpicture}
\legende{Schéma synthétique du cycle ovarien de 28 jours : variations hormonales (E2, LH, P4), ovulation J14, phases utérines et fenêtre de fécondité.}
\end{popfigure}
\item \textbf{(2 pts) Fenêtre de fécondité sur le schéma :}
Située \textbf{J11 à J15} (zone hachurée or sur le schéma).
\textbf{Justification :} Les spermatozoïdes survivent \textbf{3 à 5 jours} dans les voies génitales féminines (donnée cours : 3 jours). L'ovocyte II est fécondable \textbf{12 à 24 h} après l'ovulation. Un rapport à J11 (3 j avant ovulation J14) laisse des spermatozoïdes viables au moment de l'ovulation. Un rapport à J15 (1 j après) peut encore féconder l'ovocyte. Donc fenêtre = [J14-3 ; J14+1] = J11--J15.
\item \textbf{(2 pts) Rôle du pic de LH :}
Le \textbf{pic de LH (J13)} est le \textbf{déclencheur obligatoire de l'ovulation}. Il induit : (1) la reprise de la méiose II de l'ovocyte I bloqué en prophase I (maturation nucléaire) ; (2) la rupture de la paroi du follicule de Graaf (lyse enzymatique, prostaglandines, contraction cellules thécales) $\rightarrow$ expulsion de l'ovocyte II (ovulation) ; (3) la lutéinisation des cellules de la granulosa et de la thèque interne $\rightarrow$ formation du \textbf{corps jaune} sécrétant progestérone et œstrogènes.
\item \textbf{(3 pts) Pourquoi la méthode Ogino est risquée (indice Pearl 24) :}
\begin{enumerate}
\item \textbf{Fiabilité faible intrinsèque :} Indice de Pearl 24 = 24 grossesses pour 100 femmes-an (usage typique) $\rightarrow$ méthode \textbf{peu efficace}.
\item \textbf{Imprévisibilité de l'ovulation :} La méthode suppose un cycle régulier (ovulation J14 fixe). Or, \textbf{la phase folliculaire est variable} (durée 10-20 j) ; seul l'intervalle ovulation-règles (phase lutéale) est constant ($\approx$ 14 j). Chez l'adolescente, les cycles sont souvent \textbf{irréguliers} (anovulatoires, courts, longs) $\rightarrow$ impossible de prédire J14 \emph{a priori}.
\item \textbf{Fenêtre fertile large vs abstinence théorique :} La fenêtre réelle (J11--J15, 5 j) est souvent plus large que la période d'abstinence calculée sur une moyenne. Un cycle court (ovulation J10) $\rightarrow$ fenêtre J7--J11 ; un cycle long (ovulation J20) $\rightarrow$ fenêtre J17--J21. L'abstinence « calendrier » ne couvre pas la fenêtre réelle $\rightarrow$ \textbf{risque grossesse élevé}.
\end{enumerate}
\item \textbf{(3 pts) Deux méthodes contraceptives plus fiables pour nullipare :}
\begin{center}
\begin{tabularx}{\linewidth}{|l|X|X|}\hline
\rowcolor{popblueL} \textbf{Méthode} & \textbf{Mécanisme d'action principal} & \textbf{Avantage principal pour nullipare} \\ \hline
\textbf{Implant sous-cutané} (ex. Implanon NXT\textsuperscript{\textregistered}, 68 mg étonogestrel) & \textbf{Progestatif seul (LNG/ETG) :} Inhibition ovulation (blocage pic LH) + épaississement glaire cervicale + atrophie endomètre. & \textbf{Efficacité quasi-parfaite (Pearl 0,05-0,1) ; durée 3-5 ans ; « sans oubli » (observance parfaite) ; réversible immédiat ; sans œstrogènes (pas de risque thromboembolique) ; discret.} \\ \hline
\textbf{Stérilet au cuivre (DIU-Cu)} (ex. TCu380A) & \textbf{Effet spermicide local :} Ions Cu$^{2+}$ toxiques pour spermatozoïdes (motilité, viabilité, capacitation) + réaction inflammatoire stérile endométriale hostile à l'implantation. \emph{Pas d'effet hormonal systémique.} & \textbf{Efficacité très haute (Pearl 0,6-0,8) ; durée 5-10 ans ; « sans oubli » ; sans hormones (cycle naturel respecté) ; contraception d'urgence efficace si posé < 5 j post-rapport ; coût global faible.} \\ \hline
\end{tabularx}
\end{center}
\emph{Autre option valable :} Pilule microprogestative (désogestrel 75 µg) si contre-indication aux œstrogènes et refus méthode invasive ; mais observance quotidienne stricte (fenêtre 3-12 h) $\rightarrow$ Pearl usage typique $\approx$ 3-8 (moins bon qu'implant/DIU).
\end{enumerate}
\begin{aretenir}[Points de vigilance Probatoire]
Le schéma doit être \textbf{quantitatif} (axes gradués, pics aux bons jours). La justification de la fenêtre fécondité doit citer explicitement la durée de vie des gamètes. Pour l'argumentaire Ogino, opposer \textbf{variabilité phase folliculaire} vs \textbf{constance phase lutéale}. Pour le choix contraceptif nullipare, valoriser les \textbf{méthodes à action prolongée (LARC : Long-Acting Reversible Contraception)} : implant et DIU sont les \textbf{premiers choix recommandés par l'OMS} pour les adolescentes (efficacité maximale, observance facilitée, réversibilité rapide).
\end{aretenir}
\end{corrige}

\newpage

\coursec{Fiche bilan}

\begin{popfigure}
\begin{center}\tcbox[colback=poporange,colframe=poporange,coltext=white,arc=9pt,boxsep=4pt,left=6pt,right=6pt]{\bfseries\small Santé Reproductive}\end{center}
\vspace{-2pt}
\begin{tcbraster}[raster columns=3,raster equal height=rows,raster column skip=6pt,raster row skip=6pt,enhanced,blankest]
\begin{tcolorbox}[enhanced,colback=white,colframe=poporange,boxrule=0.9pt,arc=3pt,title={Physiologie (Rappels)},fonttitle=\bfseries\scriptsize,coltitle=white,colbacktitle=poporange,left=4pt,right=4pt,top=2pt,bottom=2pt,boxsep=1pt,toptitle=1.5pt,bottomtitle=1.5pt,halign=left]
\begin{itemize}[nosep,leftmargin=*,itemsep=1pt,font=\scriptsize]
\item Cycle ovarien/utérin
\item Fécondation fenêtre fertile
\item Gamétogenèses
\end{itemize}
\end{tcolorbox}
\begin{tcolorbox}[enhanced,colback=white,colframe=popgreen,boxrule=0.9pt,arc=3pt,title={Comportements à risque},fonttitle=\bfseries\scriptsize,coltitle=white,colbacktitle=popgreen,left=4pt,right=4pt,top=2pt,bottom=2pt,boxsep=1pt,toptitle=1.5pt,bottomtitle=1.5pt,halign=left]
\begin{itemize}[nosep,leftmargin=*,itemsep=1pt,font=\scriptsize]
\item Rapports non protégés
\item Multiplicité partenaires
\item Consommation alcool/drogues
\end{itemize}
\end{tcolorbox}
\begin{tcolorbox}[enhanced,colback=white,colframe=poppurple,boxrule=0.9pt,arc=3pt,title={Grossesses précoces},fonttitle=\bfseries\scriptsize,coltitle=white,colbacktitle=poppurple,left=4pt,right=4pt,top=2pt,bottom=2pt,boxsep=1pt,toptitle=1.5pt,bottomtitle=1.5pt,halign=left]
\begin{itemize}[nosep,leftmargin=*,itemsep=1pt,font=\scriptsize]
\item Risques maternels (fistule, éclampsie)
\item Risques fœtaux (prématurité, PMF)
\item Conséquences scolaires/sociales
\end{itemize}
\end{tcolorbox}
\begin{tcolorbox}[enhanced,colback=white,colframe=poppink,boxrule=0.9pt,arc=3pt,title={Avortements clandestins},fonttitle=\bfseries\scriptsize,coltitle=white,colbacktitle=poppink,left=4pt,right=4pt,top=2pt,bottom=2pt,boxsep=1pt,toptitle=1.5pt,bottomtitle=1.5pt,halign=left]
\begin{itemize}[nosep,leftmargin=*,itemsep=1pt,font=\scriptsize]
\item Définition (OMS)
\item Complications (hémorragie, sepsis)
\item Cadre légal Cameroun
\end{itemize}
\end{tcolorbox}
\begin{tcolorbox}[enhanced,colback=white,colframe=popgold,boxrule=0.9pt,arc=3pt,title={Contraception},fonttitle=\bfseries\scriptsize,coltitle=white,colbacktitle=popgold,left=4pt,right=4pt,top=2pt,bottom=2pt,boxsep=1pt,toptitle=1.5pt,bottomtitle=1.5pt,halign=left]
\begin{itemize}[nosep,leftmargin=*,itemsep=1pt,font=\scriptsize]
\item Hormonale (pilule, implant)
\item Barrière (préservatif)
\item IUD, Naturelle, Définitive
\end{itemize}
\end{tcolorbox}
\begin{tcolorbox}[enhanced,colback=white,colframe=popblue!70!popgreen,boxrule=0.9pt,arc=3pt,title={Prévention Alternatives},fonttitle=\bfseries\scriptsize,coltitle=white,colbacktitle=popblue!70!popgreen,left=4pt,right=4pt,top=2pt,bottom=2pt,boxsep=1pt,toptitle=1.5pt,bottomtitle=1.5pt,halign=left]
\begin{itemize}[nosep,leftmargin=*,itemsep=1pt,font=\scriptsize]
\item Approche ABC (Abstinence, Fidélité, Préservatif)
\item Éducation sexuelle complète
\item Accès services santé jeunes
\end{itemize}
\end{tcolorbox}
\end{tcbraster}

\legende{Carte mentale récapitulative de la leçon NCT05 : organisation des savoirs autour de la santé reproductive.}
\end{popfigure}

\begin{bilan}

\courssub{Définitions clés à maîtriser}

\begin{itemize}
  \item \textbf{Santé reproductive (OMS) :} État de complet bien-être physique, mental et social dans tous les aspects relatifs au système reproductif ; pas seulement absence de maladie.
  \item \textbf{Comportement à risque :} Toute conduite augmentant la probabilité d'IST, de grossesse non désirée ou de traumatisme génital (rapports non protégés, partenaires multiples, usage de substances psychoactives).
  \item \textbf{Grossesse précoce :} Grossesse survenant chez une adolescente (< 18--20 ans) dont le développement biophysique et psychosocial n'est pas achevé.
  \item \textbf{Avortement clandestin (OMS) :} Interruption de grossesse pratiquée par une personne dépourvue des qualifications requises ou dans un environnement ne répondant pas aux normes médicales minimales.
  \item \textbf{Contraception :} Ensemble des méthodes permettant d'éviter la fécondation ou l'implantation (contraception d'urgence exceptée).
  \item \textbf{Indice de Pearl :} Nombre de grossesses survenues pour 100 femmes-années d'utilisation d'une méthode contraceptive (plus il est bas, plus la méthode est efficace).
\end{itemize}

\courssub{Repères physiologiques indispensables (rappels)}

\begin{center}
\begin{tabularx}{\linewidth}{|l|X|X|}\hline
\rowcolor{popblueL}
\textbf{Notion} & \textbf{Repère clé} & \textbf{Application prévention} \\ \hline
Cycle menstruel & Durée moyenne 28 j ; Ovulation $\approx$ J14 & Fenêtre fertile $\approx$ J10--J17 (spermatozoïdes 3--5 j, ovocyte 12--24 h) \\ \hline
Fécondation & Trompe (ampoule) ; 1 spermatozoïde / ovocyte & Contraception barrière/hormonale bloque rencontre \\ \hline
Implantation & J6--J7 post-fécondation ; blastocyste & Contraception d'urgence (pilule du lendemain) agit avant \\ \hline
Ménarche / Ménopause & $\approx$ 12--13 ans / $\approx$ 50 ans & Période de fertilité $\approx$ 35--40 ans \\ \hline
\end{tabularx}
\end{center}

\courssub{Panorama contraceptif (choix responsable = méthode adaptée)}

\begin{center}
\begin{tabularx}{\linewidth}{|l|X|X|X|}\hline
\rowcolor{popblueL}
\textbf{Type} & \textbf{Exemples} & \textbf{Mécanisme principal} & \textbf{Indice Pearl (usage typique)} \\ \hline
Hormonale \emph{œstro-progestative} & Pilule combinée, Patch, Anneau & Inhibition ovulation + glaire cervicale & 0,3 -- 7 \\ \hline
Hormonale \emph{progestative seule} & Micro-pilule, Implant (3--5 ans), Injectable (3 mois) & Glaire cervicale + atrophie endomètre (+ ovulation variable) & 0,05 -- 3 \\ \hline
\emph{Dispositif Intra-Utérin} (DIU) & Cuivre (5--10 ans), Hormonal (LNG 3--5 ans) & Spermicide (Cu) / Glaire + endomètre atrophique (LNG) & 0,6 -- 0,8 \\ \hline
Barrière & Préservatif M/F, Diaphragme & Blocage physique spermatozoïdes & 13 -- 21 (M) ; 21 (F) \\ \hline
Naturelle & Ogino-Knaus, Température, Billings & Abstinence période fertile & 15 -- 25 \\ \hline
Définitive & Ligature trompes, Vasectomie & Stérilisation chirurgicale & $<$ 0,5 \\ \hline
Urgence & Lévonorgestrel (J0--J3), Ulipristal (J0--J5), DIU Cu (J0--J5) & Retard ovulation / empêche implantation & Variable (efficacité $\downarrow$ avec délai) \\ \hline
\end{tabularx}
\end{center}

\courssub{Méthodes express (Savoir-faire probatoire)}

\begin{methode}[Identifier un comportement à risque dans un cas concret]
  \begin{enumerate}
    \item Lister les faits : âge, activité sexuelle, protection utilisée, contexte (alcool, contrainte), antécédents.
    \item Repérer l'absence de \cle{préservatif} (risque IST + grossesse) ou de \cle{contraception efficace} (risque grossesse).
    \item Noter les facteurs de vulnérabilité : précocité, méconnaissance, pression pairs/partenaire, inaccessibilité services.
    \item Conclure : « Comportement à risque identifié : rapports non protégés chez une adolescente de 16 ans non scolarisée, exposant à IST/VIH et grossesse précoce. »
  \end{enumerate}
\end{methode}

\begin{methode}[Expliquer les conséquences sanitaires et sociales]
  \begin{enumerate}
    \item \textbf{Sanitaires immédiates :} IST (ulcérations, écoulements, PID/salpingite $\rightarrow$ stérilité), avortement à risque (hémorragie, sepsis, perforation utérine, décès).
    \item \textbf{Sanitaires différées :} Grossesse précoce $\rightarrow$ éclampsie, fistule obstétricale, accouchement dystocique (bassin immature), avortement spontané récurrent ; nouveau-né $\rightarrow$ prématurité, hypotrophie, mortalité néonatale.
    \item \textbf{Sociales :} Décrochage scolaire, stigmatisation, précarité économique, mariage forcé, violences basées sur le genre, transmission intergénérationnelle de la pauvreté.
  \end{enumerate}
\end{methode}

\begin{methode}[Proposer des alternatives responsables (Approche ABC+)]
  \begin{enumerate}
    \item \textbf{A -- Abstinence / Différer :} Valoriser le choix de ne pas avoir de rapports (seule méthode 100\% efficace sur IST et grossesse).
    \item \textbf{B -- Being faithful (Fidélité mutuelle) :} Réduire le nombre de partenaires + dépistage couple avant abandon préservatif.
    \item \textbf{C -- Condom (Préservatif M/F) :} Double protection (IST + grossesse) ; disponible, sans ordonnance, à utiliser à \emph{chaque} rapport.
    \item \textbf{+ Contraception moderne adaptée :} Orientation vers un service de santé (CPE, hôpital, ONG) pour méthode hormonale, DIU, implant selon critères médicaux (MEC OMS).
    \item \textbf{+ Dépistage IST/VIH régulier :} Prise en charge précoce (PEP, PrEP, traitement partner notification).
  \end{enumerate}
\end{methode}

\end{bilan}

\begin{aretenir}[Checklist avant le Probatoire]
  $\square$ Je définis la santé reproductive (OMS) et je cite 3 composantes.
  $\square$ Je décris la fenêtre fertile (J10--J17) et j'explique pourquoi elle est large (survie gamètes).
  $\square$ Je donne 4 exemples de comportements à risque et leurs déterminants (connaissance, accès, genre, pauvreté).
  $\square$ Je cite 3 risques maternels majeurs de la grossesse précoce (éclampsie, fistule, hémorragie délivrance) et 2 risques fœtaux (prématurité, PMF).
  $\square$ Je définis l'avortement clandestin (OMS) et j'énumère 4 complications graves (hémorragie, sepsis, perforation, décès).
  $\square$ Je classe les méthodes contraceptives en 5 grandes familles et je cite un exemple par famille.
  $\square$ Je compare l'indice de Pearl « usage parfait » vs « usage typique » pour le préservatif et la pilule.
  $\square$ J'explique le mécanisme d'action du DIU au cuivre et de l'implant progestatif.
  $\square$ Je sais rédiger une conduite à tenir face à un rapport non protégé (contraception d'urgence + dépistage IST + orientation contraception durable).
  $\square$ J'argumente l'approche ABC+ comme stratégie de réduction des risques adaptée au contexte camerounais.
  $\square$ Je connais les ressources locales : CPE (Centres de Protection de l'Enfance), CTA (Centres de Traitement Ambulatoire VIH), ONG (CAMNAFAW, UNFPA), numéros verts.
\end{aretenir}

\findecours

\end{document}
